% Template:     Tesis LaTeX
% Documento:    Archivo principal
% Versión:      3.4.8 (02/12/2025)

% CREACIÓN DEL DOCUMENTO
\documentclass[
	spanish, % Idioma: spanish, english, etc.	
	letterpaper, oneside
]{book}

% INFORMACIÓN DEL DOCUMENTO
% DESIGN OF ROBUST CONTROL SYSTEMS
% Essentials of Robust Control - Kemin Zhou
\def\documenttitle {Esenciales de Control Robusto}
\def\documentsubtitle {}
\def\degreetitle {
	% Un enfoque teórico - Kemin Zhou
	Un enfoque teórico
	\bigbreak\vspace{0.3cm}
	Kemin Zhou
}

\def\universityname {Universidad Autónoma de Nuevo León}
\def\universityfaculty {Facultad de Ingeniería Mecánica y Eléctrica}
\def\universitydepartment {Ingeniería Eléctrica}
\def\universitydepartmentimage {departamentos/fime-logo-cel}
\def\universitydepartmentimagecfg {height=3cm}
\def\universitylocation {San Nicolás de los Garza, Nuevo León, México}

% INTEGRANTES, PROFESORES Y FECHAS
\def\documentauthor {Joanathan de la Cruz}
\def\documentdate {\the\year}

\def\portrait {
	\begin{center}
	\vspace{1.5cm} ~ \\
	\MakeUppercase{\textbf{\documenttitle}} ~ \\
	\vspace{1.5cm}
	\MakeUppercase{\degreetitle} ~ \\
	\vfill
	\begin{tabular}{c}
		\MakeUppercase{\textbf{\documentauthor}} \\ \\
		\vspace{1.0cm} \\
		% PROFESOR GUÍA: \\
		% PROFESOR 1 \\
		% \vspace{0.5cm} \\
		% MIEMBROS DE LA COMISIÓN: \\
		% PROFESOR 2 \\
		% PROFESOR 3 \\
		% \vspace{0.5cm} \\
		% Este trabajo ha sido parcialmente financiado por: \\
		% NOMBRE INSTITUCIÓN \\
		\vspace{0.5cm} \\
		\MakeUppercase{\universitylocation} \\
		\MakeUppercase{\documentdate}
	\end{tabular}
	\end{center}
}
\def\abstracttable {
	% \begin{tabular}{l}
	% 	RESUMEN DE LA MEMORIA PARA OPTAR \\
	% 	AL TÍTULO DE MAGÍSTER EN CIENCIAS \\
	% 	DE LA INGENIERÍA \\
	% 	POR: \MakeUppercase{\documentauthor} \\
	% 	FECHA: \MakeUppercase{\documentdate} \\
	% 	PROF. GUÍA: NOMBRE PROFESOR
	% \end{tabular}
}

% IMPORTACIÓN DEL TEMPLATE
\input{template}

% INICIO DE LAS PÁGINAS
\begin{document}

% PORTADA
\templatePortrait

% CONFIGURACIÓN DE PÁGINA Y ENCABEZADOS
\templatePagecfg

% RESUMEN O ABSTRACT
\begin{abstractd}
% 	Robustness of control systems to disturbances and uncertainties has always been the central issue in feedback control. Feedback would not be needed for most control systems if there were no disturbances and uncertainties. Developing multivariable robust control methods has been the focal point in the last two decades in the control community. The state-of-the-art $\mathcal{H}_{\infty}$ robust control theory is the result of this effort.

% This book introduces some essentials of robust and $\mathcal{H}_{\infty}$ control theory. It grew from another book by this author, John C. Doyle, and Keith Glover, entitled Robust and Optimal Control, which has been extensively class-tested in many universities around the world. Unlike that book, which is intended primarily as a comprehensive reference of robust and $\mathcal{H}_{\infty}$ control theory, this book is intended to be a text for a graduate course in multivariable control. It is also intended to be a reference for practicing control engineers who are interested in applying the state-of-the-art robust control techniques in their applications. With this objective in mind, I have streamlined the presentation, added more than 50 illustrative examples, included many related MATLAB ${ }^{\text {b }}$ commands ${ }^1$ and more than 150 exercise problems, and added some recent developments in the area of robust control such as gap metric, $\nu$-gap metric, model validation, and mixed $\mu$ problem. In addition, many proofs are completely rewritten and some advanced topics are either deleted completely or do not get an in-depth treatment.

% The prerequisite for reading this book is some basic knowledge of classical control theory and state-space theory. The text contains more material than could be covered in detail in a one-semester or a one-quarter course. Chapter 1 gives a chapter-by-chapter summary of the main results presented in the book, which could be used as a guide for the selection of topics for a specific course. Chapters 2 and 3 can be used as a refresher for some linear algebra facts and some standard linear system theory. A course focusing on $\mathcal{H}_{\infty}$ control should cover at least most parts of Chapters 4-6, 8, 9, 11-13, and Sections 14.1 and 14.2. An advanced $\mathcal{H}_{\infty}$ control course should also include the rest of Chapter 14, Chapter 16, and possibly Chapters 10, 7, and 15. A course focusing on robustness and model uncertainty should cover at least Chapters 4,5, and 8-10. Chapters 17 and 18 can be added to any advanced robust and $\mathcal{H}_{\infty}$ control course if time permits.

% I have tried hard to eliminate obvious mistakes. It is, however, impossible for me to make the book perfect. Readers are encouraged to send corrections, comments, and suggestions to me, preferably by electronic mail, at
% kemin@ee.lsu.edu
% I am also planning to put any corrections, modifications, and extensions on the Internet so that they can be obtained either from the following anonymous ftp:
% ftp gate.ee.lsu.edu cd pub/kemin/books/essentials/
% or from the author's home page:
% http://kilo.ee.lsu.edu/kemin/books/essentials/
% This book would not be possible without the work done jointly for the previous book with Professor John C. Doyle and Professor Keith Glover. I thank them for their influence on my research and on this book. Their serious attitudes toward scientific research have been reference models for me. I am especially grateful to John for having me as a research fellow in Caltech, where I had two very enjoyable years and had opportunities to catch a glimpse of his "BIG PICTURE" of control.

% I want to thank my editor from Prentice Hall, Tom Robbins, who originally proposed the idea for this book and has been a constant source of support for me while writing it. Without his support and encouragement, this project would have been a difficult one. It has been my great pleasure to work with him.

% I would like to express my sincere gratitude to Professor Bruce A. Francis for giving me many helpful comments and suggestions on this book. Professor Francis has also kindly provided many exercises in the book. I am also grateful to Professor Kang-Zhi Liu and Professor Zheng-Hua Luo, who have made many useful comments and suggestions. I want to thank Professor Glen Vinnicombe for his generous help in the preparation of Chapters 16 and 17. Special thanks go to Professor Jianqing Mao for providing me the opportunity to present much of this material in a series of lectures at Beijing University of Aeronautics and Astronautics in the summer of 1996.

% In addition, I would like to thank all those who have helped in many ways in making this book possible, especially Professor Pramod P. Khargonekar, Professor André Tits, Professor Andrew Packard, Professor Jie Chen, Professor Jakob Stoustrup, Professor Hans Henrik Niemann, Professor Malcolm Smith, Professor Tryphon Georgiou, Professor Tongwen Chen, Professor Hitay Özbay, Professor Gary Balas, Professor Carolyn Beck, Professor Dennis S. Bernstein, Professor Mohamed Darouach, Dr. Bobby Bodenheimer, Professor Guoxiang Gu, Dr. Weimin Lu, Dr. John Morris, Dr. Matt Newlin, Professor Li Qiu, Professor Hector P. Rotstein, Professor Andrew Teel, Professor Jagannathan Ramanujam, Dr. Linda G. Bushnell, Xiang Chen, Greg Salomon, Pablo A. Parrilo, and many other people.



	La robustez de los sistemas de control frente a perturbaciones e incertidumbres siempre ha sido el tema central en el control por retroalimentación. La retroalimentación no sería necesaria para la mayoría de los sistemas de control si no hubiera perturbaciones e incertidumbres. El desarrollo de métodos de control robusto multivariable ha sido el punto focal en las últimas dos décadas en la comunidad de control. La teoría de control robusto $\mathcal{H}_{\infty}$ es el resultado de este esfuerzo.

	Este libro introduce algunos aspectos esenciales de la teoría de control robusto y $\mathcal{H}_{\infty}$. Creció a partir de otro libro de este autor, John C. Doyle, y Keith Glover, titulado Robust and Optimal Control, que ha sido ampliamente probado en muchas universidades alrededor del mundo. A diferencia de ese libro, que está destinado principalmente como una referencia completa de la teoría de control robusto y $\mathcal{H}_{\infty}$, este libro está destinado a ser un texto para un curso de posgrado en control multivariable. También está destinado a ser una referencia para ingenieros de control practicantes que estén interesados en aplicar las técnicas de control robusto más avanzadas en sus aplicaciones. Con este objetivo en mente, he simplificado la presentación, agregado más de 50 ejemplos ilustrativos, incluido muchos comandos relacionados de MATLAB ${ }^{\text {b }}$ y más de 150 problemas de ejercicio, y agregado algunos desarrollos recientes en el área de control robusto como la métrica de brecha, métrica $\nu$-gap, validación de modelos y problema mixto $\mu$. Además, muchas pruebas se han reescrito completamente y algunos temas avanzados se han eliminado por completo o no reciben un tratamiento en profundidad.

	El requisito previo para leer este libro es algún conocimiento básico de teoría de control clásica y teoría de espacio de estados. El texto contiene más material del que podría cubrirse en detalle en un curso de un semestre o un trimestre. El Capítulo 1 ofrece un resumen capítulo por capítulo de los principales resultados presentados en el libro, que podría usarse como guía para la selección de temas para un curso específico. Los Capítulos 2 y 3 pueden usarse como un repaso de algunos hechos de álgebra lineal y algo de teoría estándar de sistemas lineales. Un curso centrado en el control $\mathcal{H}_{\infty}$ debería cubrir al menos la mayoría de las partes de los Capítulos 4-6, 8, 9, 11-13, y las Secciones 14.1 y 14.2. Un curso avanzado de control $\mathcal{H}_{\infty}$ también debería incluir el resto del Capítulo 14, el Capítulo 16, y posiblemente los Capítulos 10, 7, y 15. Un curso centrado en la robustez e incertidumbre del modelo debería cubrir al menos los Capítulos 4,5, y 8-10. Los Capítulos 17 y 18 pueden agregarse a cualquier curso avanzado de control robusto y $\mathcal{H}_{\infty}$ si el tiempo lo permite.

	% He tratado de eliminar errores obvios. Sin embargo, es imposible para mí hacer el libro perfecto. Se anima a los lectores a enviar correcciones, comentarios y sugerencias a mí, preferiblemente por correo electrónico, a kemin@ee.isu.edu

	% Planeo poner cualquier corrección, modificación y extensión en Internet para que puedan obtenerse ya sea desde el siguiente ftp anónimo: ftp gate.ee.lsu.edu cd pub/kemin/books/essentials/ o desde la página de inicio del autor: http://kilo.ee.lsu.edu/kemin/books/essentials/ Este libro no sería posible sin el trabajo realizado conjuntamente para el libro anterior con el Profesor John C. Doyle y el Profesor Keith Glover. Les agradezco por su influencia en mi investigación y en este libro. Sus actitudes serias hacia la investigación científica han sido modelos de referencia para mí. Estoy especialmente agradecido a John por haberme tenido como investigador en Caltech, donde tuve dos años muy agradables y tuve oportunidades de vislumbrar su "BIG PICTURE" del control.

\end{abstractd}

% DEDICATORIA
\begin{dedicatory}
	% Based upon the popular Robust and Optimal Control by Zhou, et al. (PH, 1995), this book offers a streamlined approach to robust control that reflects the most recent topics and developments in the field. It features coverage of state-of-the-art topics, including gap metric, v-gap metric, model validation, and real mu.
	Este libro ofrece un enfoque simplificado para el control robusto que refleja los temas y desarrollos más recientes en el campo. Presenta cobertura de temas de vanguardia, incluyendo la métrica de brecha, la métrica $\nu$-gap, validación de modelos y $\mu$ real.
	% Una frase de dedicatoria, \\
	% pueden ser dos líneas. \\

	Quiero agradecer a mi editor de Prentice Hall, Tom Robbins, quien originalmente propuso la idea para este libro y ha sido una fuente constante de apoyo para mí mientras lo escribía. Sin su apoyo y aliento, este proyecto habría sido difícil. Ha sido un gran placer trabajar con él.

	Me gustaría expresar mi sincera gratitud al Profesor Bruce A. Francis por darme muchos comentarios y sugerencias útiles sobre este libro. El Profesor Francis también ha proporcionado amablemente muchos ejercicios en el libro. También estoy agradecido al Profesor Kang-Zhi Liu y al Profesor Zheng-Hua Luo, quienes han hecho muchos comentarios y sugerencias útiles. Quiero agradecer al Profesor Glen Vinnicombe por su generosa ayuda en la preparación de los Capítulos 16 y 17. Agradecimientos especiales al Profesor Jianqing Mao por brindarme la oportunidad de presentar gran parte de este material en una serie de conferencias en la Universidad de Aeronáutica y Astronáutica de Beijing en el verano de 1996.

	Además, me gustaría agradecer a todos los que han ayudado de muchas maneras para hacer posible este libro, especialmente al Profesor Pramod P. Khargonekar, al Profesor André Tits, al Profesor Andrew Packard, al Profesor Jie Chen, al Profesor Jakob Stoustrup, al Profesor Hans Henrik Niemann, al Profesor Malcolm Smith, al Profesor Tryphon Georgiou, al Profesor Tongwen Chen, al Profesor Hitay Özbay, al Profesor Gary Balas, al Profesor Carolyn Beck, al Profesor Dennis S. Bernstein, al Profesor Mohamed Darouach, al Dr. Bobby Bodenheimer, al Profesor Guoxiang Gu, al Dr. Weimin Lu, al Dr. John Morris, al Dr. Matt Newlin, al Profesor Li Qiu, al Profesor Hector P. Rotstein, al Profesor Andrew Teel, al Profesor Jagannathan Ramanujam, a la Dr. Linda G. Bushnell, a Xiang Chen, a Greg Salomon, a Pablo A. Parrilo y a muchas otras personas.

	% I would also like to thank the following agencies for supporting my research: National Science Foundation, Army Research Office (ARO), Air Force of Scientific Research, and the Board of Regents in the State of Louisiana.

	También quiero agradecer a las siguientes agencias por su apoyo a mi investigación: Fundación Nacional de Ciencia, Oficina de Investigación del Ejército (ARO), Fuerza Aérea de Investigación Científica, y el Consejo de Regentes en el Estado de Louisiana.

	% Finally, I would like to thank my wife, Jing, and my son, Eric, for their generous support, understanding, and patience during the writing of this book.

	Finalmente, me gustaría agradecer a mi esposa, Jing, y a mi hijo, Eric, por su generoso apoyo, comprensión y paciencia durante la escritura de este libro.
	% ~ \\
	% \textbf{Saludos

	\textbf{Sapere aude}
\end{dedicatory}

% AGRADECIMIENTOS
% \begin{acknowledgments}
% 	\lipsum[1]
% \end{acknowledgments}

% TABLA DE CONTENIDOS - ÍNDICE
\templateIndex

% CONFIGURACIONES FINALES
\templateFinalcfg

% ======================= INICIO DEL DOCUMENTO =======================


\chapter{Introducción}
Este capítulo ofrece una breve descripción de los problemas considerados en este libro y de los resultados clave presentados en cada capítulo.

\section{¿De qué trata este libro?}
Este libro trata sobre teoría básica de control robusto y $\mathcal{H}_{\infty}$. Consideramos un sistema de control con posibles múltiples fuentes de incertidumbre, ruido y perturbaciones, como se muestra en la Figura 1.1.

\insertimage[\label{img:introduccion-interconexion-general-sistema}]{2024_06_29_ce7c73b22613114bd4d6g-015_724_1022_1427_541}{width=0.75\textwidth}{Interconexión general del sistema.}

Consideramos principalmente dos tipos de problemas:

\begin{itemize}
  \item Problemas de análisis: Dado un controlador, determinar si las señales controladas (incluyendo errores de seguimiento, señales de control, etc.) satisfacen las propiedades deseadas para todos los ruidos, perturbaciones e incertidumbres del modelo admisibles.
  \item Problemas de síntesis: Diseñar un controlador de modo que las señales controladas satisfagan las propiedades deseadas para todos los ruidos, perturbaciones e incertidumbres del modelo admisibles.
\end{itemize}

La mayor parte de nuestro análisis y síntesis se realizará en un marco unificado de transformación fraccional lineal (LFT). Con ese fin, mostraremos que el sistema de la Figura 1.1 puede expresarse en el diagrama general de la Figura 1.2, donde $P$ es la matriz de interconexión, $K$ es el controlador, $\Delta$ es el conjunto de todas las incertidumbres posibles, $w$ es una señal vectorial que incluye ruidos, perturbaciones y señales de referencia, $z$ es una señal vectorial que incluye todas las señales controladas y errores de seguimiento, $u$ es la señal de control y $y$ es la medición.

\insertimage[\label{img:introduccion-marco-lft-general}]{2024_06_29_ce7c73b22613114bd4d6g-016_384_428_1044_838}{width=0.35\textwidth}{Marco LFT general.}

El diagrama de bloques de la Figura 1.2 representa las siguientes ecuaciones:

$$
\begin{aligned}
{\left[\begin{array}{l}
v \\
z \\
y
\end{array}\right] } & =P\left[\begin{array}{l}
\eta \\
w \\
u
\end{array}\right] \\
\eta & =\Delta v \\
u & =K y
\end{aligned}
$$

Denotemos por $T_{z w}$ la matriz de transferencia de $w$ a $z$, y supongamos que la incertidumbre admisible $\Delta$ satisface $\|\Delta\|_{\infty}<1 / \gamma_{u}$ para algún $\gamma_{u}>0$. Entonces, nuestro problema de análisis consiste en responder si el sistema en lazo cerrado es estable para todo $\Delta$ admisible y si $\left\|T_{z w}\right\|_{\infty} \leq \gamma_{p}$ para algún $\gamma_{p}>0$ preespecificado, donde $\left\|T_{z w}\right\|_{\infty}$ es la norma $\mathcal{H}_{\infty}$ definida como $\left\|T_{z w}\right\|_{\infty}=\sup _{\omega} \bar{\sigma}\left(T_{z w}(j \omega)\right)$. El problema de síntesis consiste en diseñar un controlador $K$ de modo que se satisfagan las condiciones de estabilidad robusta y desempeño antes mencionadas.

En la forma más simple, tenemos o bien $\Delta=0$ o bien $w=0$. El primer caso se convierte en el conocido problema de control $\mathcal{H}_{\infty}$ y el segundo en el problema de estabilidad robusta. Los dos\\
problemas son equivalentes cuando $\Delta$ es una incertidumbre no estructurada de un solo bloque, mediante la aplicación del teorema de la pequeña ganancia (véase el Capítulo 8). Esta consecuencia de estabilidad robusta fue probablemente la principal motivación para el desarrollo de los métodos $\mathcal{H}_{\infty}$.

El análisis y la síntesis para sistemas con $\Delta$ de múltiples bloques pueden reducirse, en la mayoría de los casos, a un problema equivalente de $\mathcal{H}_{\infty}$ con escalados adecuados. Por lo tanto, una solución al problema de control $\mathcal{H}_{\infty}$ es la clave para todos los problemas de robustez considerados en este libro. En la siguiente sección, presentaremos un resumen capítulo por capítulo de los principales resultados del libro.

Remitimos al lector al libro Robust and Optimal Control de K. Zhou, J. C. Doyle y K. Glover [1996] para una breve reseña histórica de $\mathcal{H}_{\infty}$ y del control robusto, así como para un tratamiento detallado de algunos temas avanzados.

\section{Destacados de este libro}
En esta sección se destacan los resultados clave de cada capítulo. Los lectores deben consultar los capítulos correspondientes para los enunciados y condiciones exactos.

El Capítulo 2 repasa algunos hechos básicos de álgebra lineal.

El Capítulo 3 repasa conceptos teóricos de sistemas: controlabilidad, observabilidad, estabilizabilidad, detectabilidad, asignación de polos, teoría de observadores, polos y ceros del sistema, y realizaciones en espacio de estados.

El Capítulo 4 introduce los espacios $\mathcal{H}_{2}$ y los espacios $\mathcal{H}_{\infty}$. Se presentan métodos en espacio de estados para calcular normas de matrices de transferencia racionales reales en $\mathcal{H}_{2}$ y $\mathcal{H}_{\infty}$. Por ejemplo, sea

$$
G(s)=\left[\begin{array}{c|c}
A & B \\
\hline C & 0
\end{array}\right] \in \mathcal{R} \mathcal{H}_{\infty}
$$

Entonces

$$
\|G\|_{2}^{2}=\operatorname{trace}\left(B^{*} Q B\right)=\operatorname{trace}\left(C P C^{*}\right)
$$

y

$$
\|G\|_{\infty}=\max \{\gamma: H \text { tiene un valor propio sobre el eje imaginario }\}
$$

donde $P$ y $Q$ son los Gramianos de controlabilidad y observabilidad, y

$$
H=\left[\begin{array}{cc}
A & B B^{*} / \gamma^{2} \\
-C^{*} C & -A^{*}
\end{array}\right]
$$

El Capítulo 5 introduce la estructura de realimentación y discute su estabilidad.

\insertimage[\label{img:introduccion-estructura-realimentacion}]{2024_06_29_ce7c73b22613114bd4d6g-018_266_683_474_710}{width=0.75\textwidth}{Estructura de realimentación considerada en el análisis de estabilidad interna.}

Definimos que el sistema en lazo cerrado anterior es internamente estable si y solo si

$$
\left[\begin{array}{cc}
I & -\hat{K} \\
-P & I
\end{array}\right]^{-1}=\left[\begin{array}{cc}
(I-\hat{K} P)^{-1} & \hat{K}(I-P \hat{K})^{-1} \\
P(I-\hat{K} P)^{-1} & (I-P \hat{K})^{-1}
\end{array}\right] \in \mathcal{R} \mathcal{H}_{\infty}
$$

También se presentan caracterizaciones alternativas de la estabilidad interna usando factorizaciones coprimas.

El Capítulo 6 considera las propiedades de los sistemas de realimentación y las limitaciones de diseño. También se estudian en este capítulo las formulaciones de los problemas de control óptimo $\mathcal{H}_{2}$ y $\mathcal{H}_{\infty}$, así como la selección de funciones de ponderación.

El Capítulo 7 considera el problema de reducir el orden de un sistema dinámico lineal multivariable usando el método de truncamiento balanceado. Supongamos que

$$
G(s)=\left[\begin{array}{cc|c}
A_{11} & A_{12} & B_{1} \\
A_{21} & A_{22} & B_{2} \\
\hline C_{1} & C_{2} & D
\end{array}\right] \in \mathcal{R} \mathcal{H}_{\infty}
$$

es una realización balanceada con Gramianos de controlabilidad y observabilidad $P=Q=\Sigma=$ $\operatorname{diag}\left(\Sigma_{1}, \Sigma_{2}\right)$

$$
\begin{aligned}
& \Sigma_{1}=\operatorname{diag}\left(\sigma_{1} I_{s_{1}}, \sigma_{2} I_{s_{2}}, \ldots, \sigma_{r} I_{s_{r}}\right) \\
& \Sigma_{2}=\operatorname{diag}\left(\sigma_{r+1} I_{s_{r+1}}, \sigma_{r+2} I_{s_{r+2}}, \ldots, \sigma_{N} I_{s_{N}}\right)
\end{aligned}
$$

Entonces el sistema truncado $G_{r}(s)=\left[\begin{array}{c|c}A_{11} & B_{1} \\ \hline C_{1} & D\end{array}\right]$ es estable y satisface una cota de error aditiva:

$$
\left\|G(s)-G_{r}(s)\right\|_{\infty} \leq 2 \sum_{i=r+1}^{N} \sigma_{i}
$$

También se discute el método de truncamiento balanceado con ponderación en frecuencia.

El Capítulo 8 deriva pruebas de estabilidad robusta para sistemas bajo diversas hipótesis de modelado mediante el uso del teorema de la pequeña ganancia. En particular, mostramos que un sistema, mostrado al inicio de la página siguiente, con una incertidumbre no estructurada $\Delta \in \mathcal{R} \mathcal{H}_{\infty}$\\
con $\|\Delta\|_{\infty}<1$ es robustamente estable si y solo si $\left\|T_{z w}\right\|_{\infty} \leq 1$, donde $T_{z w}$ es la matriz de transferencia de $w$ a $z$.

\insertimage[\label{img:introduccion-incertidumbre-no-estructurada}]{2024_06_29_ce7c73b22613114bd4d6g-019_374_631_537_736}{width=0.35\textwidth}{Interconexión con incertidumbre no estructurada para el análisis de estabilidad robusta.}

El Capítulo 9 introduce la LFT en detalle. Mostramos que muchos problemas de control pueden formularse y tratarse en el marco LFT. En particular, mostramos que todo problema de análisis puede expresarse en forma LFT con cierta $\Delta(s)$ estructurada y cierta matriz de interconexión $M(s)$, y que todo problema de síntesis puede expresarse en forma LFT con una planta generalizada $G(s)$ y un controlador $K(s)$ por diseñar.\\
\insertimage[\label{img:introduccion-formulacion-lft}]{2024_06_29_ce7c73b22613114bd4d6g-019_212_994_1260_585}{width=0.75\textwidth}{Formulación general de problemas de análisis y síntesis en el marco LFT.}

El Capítulo 10 considera la estabilidad robusta y el desempeño para sistemas con múltiples fuentes de incertidumbre. Mostramos que un sistema incierto es robustamente estable y satisface cierto criterio de desempeño $\mathcal{H}_{\infty}$ para todo $\Delta_{i} \in \mathcal{R} \mathcal{H}_{\infty}$ con $\left\|\Delta_{i}\right\|_{\infty}<1$ si y solo si el valor singular estructurado $(\mu)$ del modelo de interconexión correspondiente no es mayor que 1.

\insertimage[\label{img:introduccion-estructura-mu}]{2024_06_29_ce7c73b22613114bd4d6g-019_412_700_1824_707}{width=0.35\textwidth}{Estructura de interconexión para análisis de estabilidad robusta y desempeño con $\mu$.}

El Capítulo 11 caracteriza en espacio de estados todos los controladores que estabilizan un sistema dinámico dado $G(s)$. Para una planta generalizada dada

$$
G(s)=\left[\begin{array}{cc}
G_{11}(s) & G_{12}(s) \\
G_{21}(s) & G_{22}(s)
\end{array}\right]=\left[\begin{array}{c|cc}
A & B_{1} & B_{2} \\
\hline C_{1} & D_{11} & D_{12} \\
C_{2} & D_{21} & D_{22}
\end{array}\right]
$$

mostramos que todos los controladores estabilizantes pueden parametrizarse como la matriz de transferencia de $y$ a $u$ que aparece abajo, donde $F$ y $L$ son tales que $A+L C_{2}$ y $A+B_{2} F$ son estables, y donde $Q$ es cualquier matriz de transferencia propia y estable.

\insertimage[\label{img:introduccion-parametrizacion-controladores}]{2024_06_29_ce7c73b22613114bd4d6g-020_727_808_829_648}{width=0.55\textwidth}{Estructura de parametrización de controladores estabilizantes.}

El Capítulo 12 estudia la solución estabilizante de una ecuación algebraica de Riccati (ARE). Una solución de la siguiente ARE

$$
A^{*} X+X A+X R X+Q=0
$$

se denomina solución estabilizante si $A+R X$ es estable. Sea ahora

$$
H:=\left[\begin{array}{cc}
A & R \\
-Q & -A^{*}
\end{array}\right]
$$

y sea $\mathcal{X}_{-}(H)$ el subespacio invariante estable de $H$, y

$$
\mathcal{X}_{-}(H)=\operatorname{Im}\left[\begin{array}{l}
X_{1} \\
X_{2}
\end{array}\right]
$$

donde $X_{1}, X_{2} \in \mathbb{C}^{n \times n}$. Si $X_{1}$ es no singular, entonces $X:=X_{2} X_{1}^{-1}$ queda determinado de manera única por $H$, denotado por $X=\operatorname{Ric}(H)$. Un resultado clave de este capítulo es el llamado lema real acotado,\\
que establece que una matriz de transferencia estable $G(s)$ satisface $\|G(s)\|_{\infty}<\gamma$ si y solo si existe un $X$ tal que $A+B B^{*} X / \gamma^{2}$ es estable y

$$
X A+A^{*} X+X B B^{*} X / \gamma^{2}+C^{*} C=0
$$

La teoría de control $\mathcal{H}_{\infty}$ del Capítulo 14 se derivará con base en este lema.

El Capítulo 13 trata el control óptimo de sistemas lineales invariantes en el tiempo con criterios de desempeño cuadrático (es decir, problemas $\mathcal{H}_{2}$). Consideramos un sistema dinámico descrito por una LFT con

$$
G(s)=\left[\begin{array}{c|cc}
A & B_{1} & B_{2} \\
\hline C_{1} & 0 & D_{12} \\
C_{2} & D_{21} & 0
\end{array}\right]
$$

\insertimage[\label{img:introduccion-lft-control-h2}]{2024_06_29_ce7c73b22613114bd4d6g-021_236_317_961_901}{width=0.35\textwidth}{Planta generalizada en formulación LFT para el problema de control $\mathcal{H}_{2}$.}

Definimos

$$
\begin{aligned}
& R_{1}=D_{12}^{*} D_{12}>0, \quad R_{2}=D_{21} D_{21}^{*}>0 \\
& H_{2}:=\left[\begin{array}{cc}
A-B_{2} R_{1}^{-1} D_{12}^{*} C_{1} & -B_{2} R_{1}^{-1} B_{2}^{*} \\
-C_{1}^{*}\left(I-D_{12} R_{1}^{-1} D_{12}^{*}\right) C_{1} & -\left(A-B_{2} R_{1}^{-1} D_{12}^{*} C_{1}\right)^{*}
\end{array}\right] \\
& J_{2}:=\left[\begin{array}{cc}
\left(A-B_{1} D_{21}^{*} R_{2}^{-1} C_{2}\right)^{*} & -C_{2}^{*} R_{2}^{-1} C_{2} \\
-B_{1}\left(I-D_{21}^{*} R_{2}^{-1} D_{21}\right) B_{1}^{*} & -\left(A-B_{1} D_{21}^{*} R_{2}^{-1} C_{2}\right)
\end{array}\right] \\
& X_{2}:=\operatorname{Ric}\left(H_{2}\right) \geq 0, \quad Y_{2}:=\operatorname{Ric}\left(J_{2}\right) \geq 0 \\
& F_{2}:=-R_{1}^{-1}\left(B_{2}^{*} X_{2}+D_{12}^{*} C_{1}\right), \quad L_{2}:=-\left(Y_{2} C_{2}^{*}+B_{1} D_{21}^{*}\right) R_{2}^{-1}
\end{aligned}
$$

Entonces, el controlador óptimo $\mathcal{H}_{2}$ (es decir, el controlador que minimiza $\left\|T_{z w}\right\|_{2}$ ) viene dado por

$$
K_{\mathrm{opt}}(s):=\left[\begin{array}{c|c}
A+B_{2} F_{2}+L_{2} C_{2} & -L_{2} \\
\hline F_{2} & 0
\end{array}\right]
$$

El Capítulo 14 considera primero un problema de control $\mathcal{H}_{\infty}$ con la planta generalizada $G(s)$ dada en el Capítulo 13, pero con algunas simplificaciones adicionales: $R_{1}=I, R_{2}=I$, $D_{12}^{*} C_{1}=0$ y $B_{1} D_{21}^{*}=0$. Mostramos que existe un controlador admisible tal que $\left\|T_{z w}\right\|_{\infty}<\gamma$ si y solo si se cumplen las siguientes tres condiciones:

(i) $H_{\infty} \in \operatorname{dom}(\operatorname{Ric})$ and $X_{\infty}:=\operatorname{Ric}\left(H_{\infty}\right)>0$, where

$$
H_{\infty}:=\left[\begin{array}{cc}
A & \gamma^{-2} B_{1} B_{1}^{*}-B_{2} B_{2}^{*} \\
-C_{1}^{*} C_{1} & -A^{*}
\end{array}\right]
$$

(ii) $J_{\infty} \in \operatorname{dom}($ Ric $)$ and $Y_{\infty}:=\operatorname{Ric}\left(J_{\infty}\right)>0$, where

$$
J_{\infty}:=\left[\begin{array}{cc}
A^{*} & \gamma^{-2} C_{1}^{*} C_{1}-C_{2}^{*} C_{2} \\
-B_{1} B_{1}^{*} & -A
\end{array}\right]
$$

(iii) $\rho\left(X_{\infty} Y_{\infty}\right)<\gamma^{2}$.

Además, un controlador admisible tal que $\left\|T_{z w}\right\|_{\infty}<\gamma$ viene dado por

$$
K_{\mathrm{sub}}=\left[\begin{array}{c|c}
\hat{A}_{\infty} & -Z_{\infty} L_{\infty} \\
\hline F_{\infty} & 0
\end{array}\right]
$$

where

$$
\begin{gathered}
\hat{A}_{\infty}:=A+\gamma^{-2} B_{1} B_{1}^{*} X_{\infty}+B_{2} F_{\infty}+Z_{\infty} L_{\infty} C_{2} \\
F_{\infty}:=-B_{2}^{*} X_{\infty}, \quad L_{\infty}:=-Y_{\infty} C_{2}^{*}, \quad Z_{\infty}:=\left(I-\gamma^{-2} Y_{\infty} X_{\infty}\right)^{-1}
\end{gathered}
$$

A continuación consideramos con más detalle el problema general de control $\mathcal{H}_{\infty}$. Indicamos cómo pueden relajarse varias hipótesis para abarcar otros problemas más complejos, como los problemas de control singular. También consideramos el control integral en la teoría $\mathcal{H}_{2}$ y $\mathcal{H}_{\infty}$, y mostramos cómo la solución general de $\mathcal{H}_{\infty}$ puede usarse para resolver el problema de filtrado $\mathcal{H}_{\infty}$.

El Capítulo 15 considera el diseño de controladores de orden reducido mediante reducción de controladores. Se presta especial atención a los métodos de reducción que preservan la estabilidad y el desempeño en lazo cerrado. Se presentan métodos que proporcionan condiciones suficientes en términos de reducción de modelo ponderada en frecuencia.

El Capítulo 16 resuelve primero un problema especial de minimización $\mathcal{H}_{\infty}$. Sea $P=\tilde{M}^{-1} \tilde{N}$ una factorización coprima izquierda normalizada. Entonces mostramos que

$$
\begin{aligned}
\inf _{K \text { estabilizante }}\left\|\left[\begin{array}{c}
K \\
I
\end{array}\right](I+P K)^{-1}\left[\begin{array}{ll}
I & P
\end{array}\right]\right\|_{\infty} \\
=\inf _{K \text { estabilizante }}\left\|\left[\begin{array}{c}
K \\
I
\end{array}\right](I+P K)^{-1} \tilde{M}^{-1}\right\|_{\infty}=\left(\sqrt{1-\left\|\left[\begin{array}{ll}
\tilde{N} & \tilde{M}
\end{array}\right]\right\|_{H}^{2}}\right)^{-1} .
\end{aligned}
$$

Esto implica que existe un controlador robustamente estabilizante para

$$
P_{\Delta}=\left(\tilde{M}+\tilde{\Delta}_{M}\right)^{-1}\left(\tilde{N}+\tilde{\Delta}_{N}\right)
$$

con

$$
\left\|\left[\begin{array}{ll}
\tilde{\Delta}_{N} & \tilde{\Delta}_{M}
\end{array}\right]\right\|_{\infty}<\epsilon
$$

si y solo si

$$
\epsilon \leq \sqrt{1-\left\|\left[\begin{array}{ll}
\tilde{N} & \tilde{M}
\end{array}\right]\right\|_{H}^{2}}
$$

Usando este resultado de estabilización, se propone una técnica de diseño por conformación de lazo (loop-shaping). La técnica propuesta usa únicamente el concepto básico de los métodos de loop-shaping, y luego se emplea un controlador de estabilización robusta para el sistema perturbado en factores coprimos normalizados para construir el controlador final.

El Capítulo 17 introduce la métrica gap y la métrica $\nu$-gap. Se discuten la interpretación en el dominio de la frecuencia y las aplicaciones de la métrica $\nu$-gap. También se considera la reducción de orden de controladores en el marco de la métrica gap o $\nu$-gap.

El Capítulo 18 considera brevemente los problemas de validación de modelos y de análisis y síntesis $\mu$ mixtos (reales y complejos).

La mayoría de los cálculos y ejemplos de este libro se realizan usando Matlab. Dado que utilizaremos MATLAB como herramienta computacional principal, se supone que los lectores tienen conocimientos básicos de las operaciones de MatcaB (por ejemplo, cómo ingresar vectores y matrices). También hemos incluido en este libro breves explicaciones de comandos de Matlab, SimulinK ${ }^{\circledR}$, Control System Toolbox y $\mu$ Analysis and Synthesis Toolbox ${ }^{1}$. En particular, este libro está escrito de manera consistente con $\mu$ Analysis and Synthesis Toolbox. (Robust Control Toolbox, LMI Control Toolbox y otros paquetes de software también pueden usarse con este libro.) Por ello, es útil que los lectores tengan acceso a esta toolbox. En este punto se sugiere probar los siguientes programas de demostración de esta toolbox.

\textbf{msdemo1}

\textbf{msdemo2}
Introduciremos muchos más comandos de MatcaB en los capítulos siguientes.

\section{Notas y referencias}
La formulación original del problema de control $\mathcal{H}_{\infty}$ puede encontrarse en Zames [1981]. Actualmente se han establecido relaciones entre $\mathcal{H}_{\infty}$ y muchos otros temas de control: por ejemplo, control sensible al riesgo de Whittle [1990]; juegos diferenciales (véase Başar y Bernhard [1991], Limebeer, Anderson, Khargonekar y Green [1992]; Green y Limebeer [1995]); representación chain-scattering y factorización J-lossless (Green [1992] y Kimura [1997]). Véase también Zhou, Doyle y Glover [1996] para discusiones y referencias adicionales. La teoría en espacio de estados de $\mathcal{H}_{\infty}$ también se ha desarrollado mucho más, generalizando de invariante en el tiempo a variante en el tiempo, de horizonte infinito a horizonte finito, y de dimensión finita a dimensión infinita, e incluso a algunos escenarios no lineales.[\^{}0]

\section{Problemas}

Problema 1.1 Resolveremos primero un problema fácil. Cuando lees un artículo o un libro, a menudo encuentras una afirmación como esta: ``Es fácil...''. Lo que el autor realmente quiso decir fue una de las siguientes opciones: (a) realmente es fácil; (b) parece fácil; (c) es fácil para un experto; (d) el autor no sabe cómo demostrarlo, pero cree que es correcto. Ahora demuestre que cuando digo ``Es fácil'' en este libro, quiero decir que realmente es fácil. (Pista: Si puede demostrarlo después de leer todo el libro, pídale a su jefe un ascenso. Si no puede demostrarlo después de leer todo el libro, tire el libro y escriba uno usted mismo. Recuerde usar algo como ``es fácil...'' si no está seguro de lo que está diciendo.)
\input{2_Algebra_Lineal.tex}
\chapter{Sistemas lineales}

Este capítulo revisa algunos conceptos básicos de teoría de sistemas. Se definen las nociones de controlabilidad, observabilidad, estabilizabilidad y detectabilidad, y se resumen varias caracterizaciones algebraicas y geométricas de estas nociones. Luego se introduce la teoría de observadores. Se estudian las interconexiones y realizaciones de sistemas. Finalmente, se introducen los conceptos de polos y ceros del sistema.

\section{Descripciones de sistemas dinámicos lineales}
Sea un sistema dinámico lineal invariante en el tiempo de dimensión finita (FDLTI) descrito por las siguientes ecuaciones diferenciales lineales con coeficientes constantes:


\begin{align*}
\dot{x} & =A x+B u, \quad x\left(t_{0}\right)=x_{0}  \tag{3.1}\\
y & =C x+D u \tag{3.2}
\end{align*}


donde $x(t) \in \mathbb{R}^{n}$ se llama estado del sistema, $x\left(t_{0}\right)$ se llama condición inicial del sistema, $u(t) \in \mathbb{R}^{m}$ se llama entrada del sistema, y $y(t) \in \mathbb{R}^{p}$ es la salida del sistema. Las matrices $A, B, C$ y $D$ son matrices reales constantes de dimensiones apropiadas. Un sistema dinámico con una sola entrada $(m=1)$ y una sola salida $(p=1)$ se llama sistema SISO (single-input single-output); en caso contrario se llama sistema MIMO (multiple-input multiple-output). La matriz de transferencia correspondiente de $u$ a $y$ se define como

$$
Y(s)=G(s) U(s)
$$

donde $U(s)$ y $Y(s)$ son las transformadas de Laplace de $u(t)$ y $y(t)$ con condición inicial nula $(x(0)=0)$. Por tanto,

$$
G(s)=C(s I-A)^{-1} B+D
$$

Nótese que las ecuaciones del sistema (3.1) y (3.2) pueden escribirse en una forma matricial más compacta:

$$
\left[\begin{array}{l}
\dot{x} \\
y
\end{array}\right]=\left[\begin{array}{cc}
A & B \\
C & D
\end{array}\right]\left[\begin{array}{l}
x \\
u
\end{array}\right]
$$

Para agilizar cálculos con matrices de transferencia, usaremos la siguiente notación:

$$
\left[\begin{array}{l|l}
A & B \\
\hline C & D
\end{array}\right]:=C(s I-A)^{-1} B+D .
$$

En Matlab el sistema también puede escribirse en forma empaquetada usando el comando

$\gg \mathbf{G}=\mathbf{p c k}(\mathbf{A}, \mathbf{B}, \mathbf{C}, \mathbf{D})$ \% empaqueta la realización en forma particionada

$\gg$ seesys(G) \% muestra $G$ en formato particionado

$\gg[\mathbf{A}, \mathbf{B}, \mathbf{C}, \mathbf{D}]=\operatorname{unpck}(\mathbf{G}) \%$ desempaqueta la matriz del sistema

Nótese que

$$
\left[\begin{array}{ll}
A & B \\
C & D
\end{array}\right]
$$

es una matriz real por bloques, no una función de transferencia.

\textbf{Comandos ilustrativos de MATLAB:}
$\gg \mathbf{G}=\mathbf{p c k}([],[],[], 10) \%$ crea una matriz de sistema constante

$\gg[\mathbf{y}, \mathbf{x}, \mathbf{t}]=\operatorname{step}(\mathbf{A}, \mathbf{B}, \mathbf{C}, \mathbf{D}, \mathbf{Iu}) \% \mathrm{Iu}=i$ (respuesta al escalón del canal $i$)

$\gg[\mathbf{y}, \mathbf{x}, \mathbf{t}]=\operatorname{initial}\left(\mathbf{A}, \mathbf{B}, \mathbf{C}, \mathbf{D}, \mathbf{x}_{0}\right) \%$ respuesta por condición inicial $x_{0}$

$\gg[\mathbf{y}, \mathbf{x}, \mathbf{t}]=$ impulse(A, B, C, D, Iu) \% respuesta al impulso del canal Iu

$\gg[\mathbf{y}, \mathbf{x}]=\operatorname{lsim}(\mathbf{A}, \mathbf{B}, \mathbf{C}, \mathbf{D}, \mathbf{U}, \mathbf{T}) \% \mathrm{U}$ es una entrada matricial de tamaño length$(T) \times \operatorname{column}(B)$; $T$ son los puntos de muestreo.

Comandos relacionados de MATLAB: minfo, trsp, cos\_tr, sin\_tr, siggen

\section{Controlabilidad y observabilidad}
Ahora pasamos a algunos conceptos muy importantes en teoría de sistemas lineales.

Definición 3.1 El sistema dinámico descrito por la ecuación (3.1), o el par $(A, B)$, se dice controlable si, para cualquier estado inicial $x(0)=x_{0}$, cualquier $t_{1}>0$ y cualquier estado final $x_{1}$, existe una entrada (a trozos continua) $u(\cdot)$ tal que la solución de la ecuación (3.1) satisface $x\left(t_{1}\right)=x_{1}$. En caso contrario, el sistema o el par $(A, B)$ se dice no controlable.

La controlabilidad (y la observabilidad introducida después) de un sistema puede verificarse mediante criterios algebraicos o geométricos.

Teorema 3.1 Son equivalentes:

(i) $(A, B)$ es controlable.

(ii) La matriz

$$
W_{c}(t):=\int_{0}^{t} e^{A \tau} B B^{*} e^{A^{*} \tau} d \tau
$$

es definida positiva para cualquier $t>0$.

(iii) La matriz de controlabilidad

$$
\mathcal{C}=\left[\begin{array}{lllll}
B & A B & A^{2} B & \ldots & A^{n-1} B
\end{array}\right]
$$

tiene rango completo por filas o, en otras palabras, $\langle A \mid \operatorname{Im} B\rangle:=\sum_{i=1}^{n} \operatorname{Im}\left(A^{i-1} B\right)=\mathbb{R}^{n}$.

(iv) La matriz $[A-\lambda I, B]$ tiene rango completo por filas para todo $\lambda \in \mathbb{C}$.

(v) Sean $\lambda$ y $x$ cualquier valor propio y cualquier vector propio izquierdo correspondiente de $A$ (es decir, $x^{*} A=x^{*} \lambda$); entonces $x^{*} B \neq 0$.

(vi) Los valores propios de $A+B F$ pueden asignarse libremente (con la restricción de que los complejos vayan en pares conjugados) mediante una elección apropiada de $F$.

Ejemplo 3.1 Sea $A=\left[\begin{array}{ll}2 & 0 \\ 0 & 2\end{array}\right]$ y $B=\left[\begin{array}{l}1 \\ 1\end{array}\right]$. Entonces $x_{1}=\left[\begin{array}{l}1 \\ 0\end{array}\right]$ y $x_{2}=\left[\begin{array}{l}0 \\ 1\end{array}\right]$ son vectores propios independientes de $A$ y $x_{i}^{*} B \neq 0, i=1,2$. Sin embargo, esto no debe llevar a concluir que $(A, B)$ es controlable. De hecho, $x=x_{1}-x_{2}$ también es vector propio de $A$ y $x^{*} B=0$, lo cual implica que $(A, B)$ no es controlable. Por tanto, al usar el criterio (v), se debe verificar para todos los vectores propios posibles.

Definición 3.2 Un sistema dinámico no forzado $\dot{x}=A x$ se dice estable si todos los valores propios de $A$ están en el semiplano izquierdo abierto; es decir, $\operatorname{Re} \lambda(A)<0$. Una matriz $A$ con esta propiedad se dice estable o de Hurwitz.

Definición 3.3 El sistema dinámico de la ecuación (3.1), o el par $(A, B)$, se dice estabilizable si existe una realimentación de estado $u=F x$ tal que el sistema sea estable (es decir, $A+B F$ sea estable).

Es más apropiado llamar a esta estabilizabilidad como estabilizabilidad por realimentación de estado, para diferenciarla de la estabilizabilidad por realimentación de salida definida más adelante.

El siguiente teorema es consecuencia del Teorema 3.1.

Teorema 3.2 Son equivalentes:

(i) $(A, B)$ es estabilizable.

(ii) La matriz $[A-\lambda I, B]$ tiene rango completo por filas para todo $\operatorname{Re} \lambda \geq 0$.

(iii) Para todo $\lambda$ y $x$ tales que $x^{*} A=x^{*} \lambda$ y $\operatorname{Re} \lambda \geq 0$, se cumple $x^{*} B \neq 0$.

(iv) Existe una matriz $F$ tal que $A+B F$ es Hurwitz.

Ahora consideramos las nociones duales: observabilidad y detectabilidad del sistema descrito por las ecuaciones (3.1) y (3.2).

Definición 3.4 El sistema dinámico descrito por las ecuaciones (3.1) y (3.2), o por el par $(C, A)$, se dice observable si, para cualquier $t_{1}>0$, el estado inicial $x(0)=x_{0}$ puede determinarse a partir de la historia temporal de la entrada $u(t)$ y la salida $y(t)$ en el intervalo $\left[0, t_{1}\right]$. En caso contrario, el sistema, o $(C, A)$, se dice no observable.

Teorema 3.3 Son equivalentes:

(i) $(C, A)$ es observable.

(ii) La matriz

$$
W_{o}(t):=\int_{0}^{t} e^{A^{*} \tau} C^{*} C e^{A \tau} d \tau
$$

es definida positiva para cualquier $t>0$.

(iii) La matriz de observabilidad

$$
\mathcal{O}=\left[\begin{array}{c}
C \\
C A \\
C A^{2} \\
\vdots \\
C A^{n-1}
\end{array}\right]
$$

tiene rango completo por columnas o $\bigcap_{i=1}^{n} \operatorname{Ker}\left(C A^{i-1}\right)=0$.

(iv) La matriz $\left[\begin{array}{c}A-\lambda I \\ C\end{array}\right]$ tiene rango completo por columnas para todo $\lambda \in \mathbb{C}$.

(v) Sean $\lambda$ y $y$ cualquier valor propio y cualquier vector propio derecho correspondiente de $A$ (es decir, $A y=\lambda y$); entonces $C y \neq 0$.

(vi) Los valores propios de $A+L C$ pueden asignarse libremente (con la restricción de pares conjugados para valores propios complejos) mediante una elección apropiada de $L$.

(vii) $\left(A^{*}, C^{*}\right)$ es controlable.

Definición 3.5 El sistema, o el par $(C, A)$, se dice detectable si $A+L C$ es estable para alguna matriz $L$.

Teorema 3.4 Son equivalentes:

(i) $(C, A)$ es detectable.

(ii) La matriz $\left[\begin{array}{c}A-\lambda I \\ C\end{array}\right]$ tiene rango completo por columnas para todo $\operatorname{Re} \lambda \geq 0$.

(iii) Para todo $\lambda$ y $x$ tales que $A x=\lambda x$ y $\operatorname{Re} \lambda \geq 0$, se cumple $C x \neq 0$.

(iv) Existe una matriz $L$ tal que $A+L C$ es Hurwitz.

(v) $\left(A^{*}, C^{*}\right)$ es estabilizable.

Las condiciones (iv) y (v) de los Teoremas 3.1 y 3.3, y las condiciones (ii) y (iii) de los Teoremas 3.2 y 3.4, suelen llamarse pruebas de Popov-Belevitch-Hautus (PBH). En particular, las siguientes definiciones de controlabilidad y observabilidad modal suelen ser útiles.

Definición 3.6 Sea $\lambda$ un valor propio de $A$ o, equivalentemente, un modo del sistema. Entonces se dice que el modo $\lambda$ es controlable (observable) si $x^{*} B \neq 0$ ($C x \neq 0$) para todos los vectores propios izquierdos (derechos) de $A$ asociados con $\lambda$; es decir, $x^{*} A=\lambda x^{*}$ ($A x=\lambda x$) y $0 \neq x \in \mathbb{C}^{n}$. En caso contrario, el modo se dice no controlable (no observable).

De esto se sigue que un sistema es controlable (observable) si y solo si todo modo es controlable (observable). De manera similar, un sistema es estabilizable (detectable) si y solo si todo modo inestable es controlable (observable).

Comandos ilustrativos de MATLAB:

$\gg \mathcal{C}=\operatorname{ctrb}(\mathrm{A}, \mathrm{B}) ; \mathcal{O}=\operatorname{obsv}(\mathrm{A}, \mathrm{C})$

$\gg \mathbf{W}_{\mathbf{c}}(\infty)=\operatorname{gram}(\mathbf{A}, \mathbf{B}) ; \%$ si $A$ es estable.

$\gg \mathbf{F}=-\mathbf{p l a c e}(\mathbf{A}, \mathbf{B}, \mathbf{P}) \quad \% P$ es un vector con valores propios deseados.

Comandos relacionados de MATLAB: ctrbf, obsvf, canon, strans, acker

\section{Observadores y controladores basados en observadores}
Es claro que si un sistema es controlable y los estados del sistema están disponibles para realimentación, entonces los polos en lazo cerrado pueden asignarse arbitrariamente mediante una realimentación constante. Sin embargo, en la mayoría de aplicaciones prácticas, los estados del sistema no son completamente accesibles y el diseñador solo conoce la salida $y$ y la entrada $u$. Por ello, la estimación de los estados del sistema a partir de la información disponible de $y$ y $u$ suele ser necesaria para alcanzar ciertos objetivos de diseño. En esta sección consideramos este problema de estimación de estado y su aplicación al control por realimentación.

Considere una planta modelada por las ecuaciones (3.1) y (3.2). Un observador es un sistema dinámico con entrada $(u, y)$ y salida (digamos, $\hat{x}$), que estima asintóticamente el estado $x$, es decir, $\hat{x}(t)-x(t) \rightarrow 0$ cuando $t \rightarrow \infty$ para cualquier estado inicial y para toda entrada.

Teorema 3.5 Un observador existe si y solo si $(C, A)$ es detectable. Además, si $(C, A)$ es detectable, entonces un observador de Luenberger de orden completo está dado por


\begin{align*}
\dot{q} & =A q+B u+L(C q+D u-y)  \tag{3.3}\\
\hat{x} & =q \tag{3.4}
\end{align*}


donde $L$ es cualquier matriz tal que $A+L C$ es estable.

Recuerde que, para un sistema dinámico descrito por las ecuaciones (3.1) y (3.2), si $(A, B)$ es controlable y el estado $x$ está disponible para realimentación, entonces existe una realimentación de estado $u=F x$ tal que los polos de lazo cerrado pueden asignarse arbitrariamente. De forma similar, si $(C, A)$ es observable, entonces los polos del observador pueden colocarse arbitrariamente para hacer que el estimador de estado $\hat{x}$ se aproxime a $x$ tan rápido como se desee. Consideremos ahora qué ocurre si los estados del sistema no están disponibles para realimentación y debe usarse el estado estimado. En ese caso, el controlador tiene la siguiente dinámica:

$$
\begin{aligned}
\dot{\hat{x}} & =(A+L C) \hat{x}+B u+L D u-L y \\
u & =F \hat{x}
\end{aligned}
$$

Entonces, las ecuaciones de estado del sistema total están dadas por

$$
\left[\begin{array}{c}
\dot{x} \\
\dot{\hat{x}}
\end{array}\right]=\left[\begin{array}{cc}
A & B F \\
-L C & A+B F+L C
\end{array}\right]\left[\begin{array}{l}
x \\
\hat{x}
\end{array}\right]
$$

Sea $e:=x-\hat{x}$; entonces la ecuación del sistema queda

$$
\left[\begin{array}{c}
\dot{e} \\
\dot{\hat{x}}
\end{array}\right]=\left[\begin{array}{cc}
A+L C & 0 \\
-L C & A+B F
\end{array}\right]\left[\begin{array}{l}
e \\
\hat{x}
\end{array}\right]
$$

y los polos en lazo cerrado constan de dos partes: los polos resultantes de la realimentación de estado $\lambda_{i}(A+B F)$ y los polos resultantes de la estimación de estado $\lambda_{j}(A+L C)$. Si $(A, B)$ es controlable y $(C, A)$ es observable, entonces existen $F$ y $L$ tales que los valores propios de $A+B F$ y $A+L C$ pueden asignarse arbitrariamente. En particular, pueden hacerse estables. Nótese que también puede obtenerse un resultado ligeramente más débil cuando $(A, B)$ y $(C, A)$ son solo estabilizable y detectable, respectivamente.

El controlador anterior se denomina controlador basado en observador y se denota como

$$
u=K(s) y
$$

y

$$
K(s)=\left[\begin{array}{c|c}
A+B F+L C+L D F & -L \\
\hline F & 0
\end{array}\right]
$$

Ahora denotemos la planta en lazo abierto por

$$
G=\left[\begin{array}{c|c}
A & B \\
\hline C & D
\end{array}\right]
$$

entonces el sistema de realimentación en lazo cerrado es el mostrado a continuación:

\insertimage[\label{img:sistemas-lineales-realimentacion-cerrada}]{2024_06_29_ce7c73b22613114bd4d6g-047_194_478_689_821}{width=0.75\textwidth}{Sistema de realimentación en lazo cerrado para planta $G$ y controlador basado en observador.}

En general, si un sistema es estabilizable mediante realimentación de la salida $y$, se dice que es estabilizable por realimentación de salida. De la construcción anterior se ve claramente que un sistema es estabilizable por realimentación de salida si y solo si $(A, B)$ es estabilizable y $(C, A)$ es detectable.

Ejemplo 3.2 Sea $A=\left[\begin{array}{ll}1 & 2 \\ 1 & 0\end{array}\right], B=\left[\begin{array}{l}1 \\ 0\end{array}\right]$, y $C=\left[\begin{array}{ll}1 & 0\end{array}\right]$. Diseñaremos una realimentación de estado $u=F x$ tal que los polos en lazo cerrado estén en $\{-2,-3\}$. Esto puede hacerse escogiendo $F=\left[\begin{array}{ll}-6 & -8\end{array}\right]$ usando

$$
\gg F=-\operatorname{place}(A, B,[-2,-3])
$$

Ahora suponga que los estados no están disponibles para realimentación y queremos construir un observador de modo que los polos del observador estén en $\{-10,-10\}$. Entonces $L=\left[\begin{array}{l}-21 \\ -51\end{array}\right]$ puede obtenerse usando

$$
\gg L=-\operatorname{acker}\left(A', C',[-10,-10]\right)'
$$

y el controlador basado en observador está dado por

$$
K(s)=\frac{-534(s+0.6966)}{(s+34.6564)(s-8.6564)}
$$

Nótese que el propio controlador estabilizante es inestable. Por supuesto, esto puede no ser deseable en la práctica.

\section{Operaciones en sistemas}
En esta sección presentamos algunos hechos sobre interconexión de sistemas. Como estas demostraciones son directas, dejaremos los detalles al lector.

Suponga que $G_{1}$ y $G_{2}$ son dos subsistemas con representaciones en espacio de estados:

$$
G_{1}=\left[\begin{array}{c|c}
A_{1} & B_{1} \\
\hline C_{1} & D_{1}
\end{array}\right], \quad G_{2}=\left[\begin{array}{c|c}
A_{2} & B_{2} \\
\hline C_{2} & D_{2}
\end{array}\right]
$$

Entonces la conexión en serie o en cascada de estos dos subsistemas es un sistema cuya salida del segundo subsistema es la entrada del primero, como se muestra en el siguiente diagrama:

\insertimage[\label{img:sistemas-lineales-conexion-cascada}]{2024_06_29_ce7c73b22613114bd4d6g-048_101_542_871_781}{width=0.75\textwidth}{Conexión en serie o en cascada de dos subsistemas $G_{1}$ y $G_{2}$.}

Esta operación, en términos de las matrices de transferencia de ambos subsistemas, es esencialmente el producto de dos matrices de transferencia. Por lo tanto, una representación para el sistema en cascada puede obtenerse como

$$
\begin{aligned}
G_{1} G_{2} & =\left[\begin{array}{c|c}
A_{1} & B_{1} \\
\hline C_{1} & D_{1}
\end{array}\right]\left[\begin{array}{c|c}
A_{2} & B_{2} \\
\hline C_{2} & D_{2}
\end{array}\right] \\
& =\left[\begin{array}{cc|c}
A_{1} & B_{1} C_{2} & B_{1} D_{2} \\
0 & A_{2} & B_{2} \\
\hline C_{1} & D_{1} C_{2} & D_{1} D_{2}
\end{array}\right]=\left[\begin{array}{cc|c}
A_{2} & 0 & B_{2} \\
B_{1} C_{2} & A_{1} & B_{1} D_{2} \\
\hline D_{1} C_{2} & C_{1} & D_{1} D_{2}
\end{array}\right] .
\end{aligned}
$$

De manera similar, la conexión en paralelo o suma de $G_{1}$ y $G_{2}$ se obtiene como

$$
G_{1}+G_{2}=\left[\begin{array}{c|c}
A_{1} & B_{1} \\
\hline C_{1} & D_{1}
\end{array}\right]+\left[\begin{array}{c|c}
A_{2} & B_{2} \\
\hline C_{2} & D_{2}
\end{array}\right]=\left[\begin{array}{cc|c}
A_{1} & 0 & B_{1} \\
0 & A_{2} & B_{2} \\
\hline C_{1} & C_{2} & D_{1}+D_{2}
\end{array}\right]
$$

Para referencia futura, también introducimos las siguientes definiciones.

Definición 3.7 La transpuesta de una matriz de transferencia $G(s)$, o sistema dual, se define como

$$
G \longmapsto G^{T}(s)=B^{*}\left(s I-A^{*}\right)^{-1} C^{*}+D^{*}
$$

o, equivalentemente,

$$
\left[\begin{array}{c|c}
A & B \\
\hline C & D
\end{array}\right] \longmapsto\left[\begin{array}{c|c}
A^{*} & C^{*} \\
\hline B^{*} & D^{*}
\end{array}\right]
$$

Definición 3.8 El sistema conjugado de $G(s)$ se define como

$$
G \longmapsto G^{\sim}(s):=G^{T}(-s)=B^{*}\left(-s I-A^{*}\right)^{-1} C^{*}+D^{*}
$$

o, equivalentemente,

$$
\left[\begin{array}{c|c}
A & B \\
\hline C & D
\end{array}\right] \longmapsto\left[\begin{array}{c|c}
-A^{*} & -C^{*} \\
\hline B^{*} & D^{*}
\end{array}\right]
$$

En particular, se cumple $G^{*}(j \omega):=[G(j \omega)]^{*}=G^{\sim}(j \omega)$.

Una matriz racional real $\hat{G}(s)$ se llama inversa de una matriz de transferencia $G(s)$ si $G(s) \hat{G}(s)=\hat{G}(s) G(s)=I$. Suponga que $G(s)$ es cuadrada y $D$ es invertible. Entonces

$$
G^{-1}=\left[\begin{array}{c|c}
A-B D^{-1} C & -B D^{-1} \\
\hline D^{-1} C & D^{-1}
\end{array}\right]
$$

Los comandos de MatLAB correspondientes son

$$
\begin{gathered}
G_{1} G_{2} \Longleftrightarrow \operatorname{mmult}\left(\mathbf{G}_{\mathbf{1}}, \mathbf{G}_{\mathbf{2}}\right), \quad\left[\begin{array}{ll}
G_{1} & G_{2}
\end{array}\right] \Longleftrightarrow \operatorname{sbs}\left(\mathbf{G}_{\mathbf{1}}, \mathbf{G}_{\mathbf{2}}\right) \\
G_{1}+G_{2} \Longleftrightarrow \operatorname{madd}\left(\mathbf{G}_{\mathbf{1}}, \mathbf{G}_{\mathbf{2}}\right), \quad G_{1}-G_{2} \Longleftrightarrow \operatorname{msub}\left(\mathbf{G}_{\mathbf{1}}, \mathbf{G}_{\mathbf{2}}\right) \\
{\left[\begin{array}{l}
G_{1} \\
G_{2}
\end{array}\right] \Longleftrightarrow \operatorname{abv}\left(\mathbf{G}_{\mathbf{1}}, \mathbf{G}_{\mathbf{2}}\right), \quad\left[\begin{array}{cc}
G_{1} & \\
& G_{2}
\end{array}\right] \Longleftrightarrow \operatorname{daug}\left(\mathbf{G}_{\mathbf{1}}, \mathbf{G}_{\mathbf{2}}\right)} \\
G^{T}(s) \Longleftrightarrow \operatorname{transp}(\mathbf{G}), \quad G^{\sim}(s) \Longleftrightarrow \operatorname{cjt}(\mathbf{G}), \quad G^{-1}(s) \Longleftrightarrow \operatorname{minv}(\mathbf{G}) \\
\alpha G(s) \Longleftrightarrow \operatorname{mscl}(\mathbf{G}, \alpha), \alpha \text { es un escalar. }
\end{gathered}
$$

Comandos relacionados de MATLAB: append, parallel, feedback, series, cloop, sclin, sclout, sel

\section{Realizaciones de espacio de estados para matrices de transferencia}
En algunos casos, la descripción natural o conveniente de un sistema dinámico es en términos de su función de transferencia matricial. Esto ocurre, por ejemplo, en sistemas muy complejos para los cuales las ecuaciones diferenciales analíticas son demasiado difíciles o complejas de escribir. Por ello, debe realizarse cierta aproximación o identificación de ingeniería; por ejemplo, se obtienen respuestas en frecuencia de entrada y salida a partir de experimentos para construir una matriz de transferencia que aproxime la dinámica del sistema. Dado que el cálculo en espacio de estados es el más conveniente para implementar en computadora, es necesaria una representación apropiada en espacio de estados para la matriz de transferencia resultante.

En general, suponga que $G(s)$ es una matriz de transferencia racional real y propia. Entonces llamamos realización de $G(s)$ a un modelo en espacio de estados $(A, B, C, D)$ tal que

$$
G(s)=\left[\begin{array}{c|c}
A & B \\
\hline C & D
\end{array}\right]
$$

Definición 3.9 Una realización en espacio de estados $(A, B, C, D)$ de $G(s)$ se dice mínima si $A$ tiene la menor dimensión posible.

Teorema 3.6 Una realización en espacio de estados $(A, B, C, D)$ de $G(s)$ es mínima si y solo si $(A, B)$ es controlable y $(C, A)$ es observable.

Ahora describimos varias formas de obtener una realización en espacio de estados para una matriz de transferencia MIMO dada $G(s)$. Primero consideraremos sistemas SIMO (una entrada y múltiples salidas) y MISO (múltiples entradas y una salida).

Sea $G(s)$ un vector columna de funciones de transferencia con $p$ salidas:

$$
G(s)=\frac{\beta_{1} s^{n-1}+\beta_{2} s^{n-2}+\cdots+\beta_{n-1} s+\beta_{n}}{s^{n}+a_{1} s^{n-1}+\cdots+a_{n-1} s+a_{n}}+d, \beta_{i} \in \mathbb{R}^{p}, d \in \mathbb{R}^{p}
$$

Entonces $G(s)=\left[\begin{array}{l|l}A & b \\ \hline C & d\end{array}\right]$ con

$$
\begin{gathered}
A:=\left[\begin{array}{ccccc}
-a_{1} & -a_{2} & \cdots & -a_{n-1} & -a_{n} \\
1 & 0 & \cdots & 0 & 0 \\
0 & 1 & \cdots & 0 & 0 \\
\vdots & \vdots & & \vdots & \vdots \\
0 & 0 & \cdots & 1 & 0
\end{array}\right] \quad b:=\left[\begin{array}{c}
1 \\
0 \\
0 \\
\vdots \\
0
\end{array}\right] \\
C=\left[\begin{array}{lllll}
\beta_{1} & \beta_{2} & \cdots & \beta_{n-1} & \beta_{n}
\end{array}\right]
\end{gathered}
$$

que es la llamada forma canónica controlable o forma canónica del controlador.

De forma dual, considere un sistema de múltiples entradas y una salida:

$$
G(s)=\frac{\eta_{1} s^{n-1}+\eta_{2} s^{n-2}+\cdots+\eta_{n-1} s+\eta_{n}}{s^{n}+a_{1} s^{n-1}+\cdots+a_{n-1} s+a_{n}}+d, \quad \eta_{i}^{*} \in \mathbb{R}^{m}, d^{*} \in \mathbb{R}^{m}
$$

Entonces

$$
G(s)=\left[\begin{array}{c|c}
A & B \\
\hline c & d
\end{array}\right]=\left[\begin{array}{ccccc|c}
-a_{1} & 1 & 0 & \cdots & 0 & \eta_{1} \\
-a_{2} & 0 & 1 & \cdots & 0 & \eta_{2} \\
\vdots & \vdots & \vdots & & \vdots & \vdots \\
-a_{n-1} & 0 & 0 & \cdots & 1 & \eta_{n-1} \\
-a_{n} & 0 & 0 & \cdots & 0 & \eta_{n} \\
\hline 1 & 0 & 0 & \cdots & 0 & d
\end{array}\right]
$$

que es una forma canónica observable o forma canónica del observador.

Para un sistema MIMO, la forma más simple y directa de obtener una realización es realizar cada elemento de la matriz $G(s)$ y luego combinar todas esas realizaciones individuales para formar una realización de $G(s)$. Para ilustrar, consideremos una matriz de transferencia (en bloques) de $2 \times 2$ como

$$
G(s)=\left[\begin{array}{ll}
G_{1}(s) & G_{2}(s) \\
G_{3}(s) & G_{4}(s)
\end{array}\right]
$$

y supongamos que $G_{i}(s)$ tiene una realización en espacio de estados de la forma

$$
G_{i}(s)=\left[\begin{array}{c|c}
A_{i} & B_{i} \\
\hline C_{i} & D_{i}
\end{array}\right], \quad i=1, \ldots, 4
$$

Entonces una realización para $G(s)$ puede obtenerse como $(\mathbf{G}=\mathbf{abv}(\mathbf{sbs}(\mathbf{G}_{\mathbf{1}}, \mathbf{G}_{\mathbf{2}}), \operatorname{sbs}(\mathbf{G}_{\mathbf{3}}, \mathbf{G}_{\mathbf{4}})))$:

$$
G(s)=\left[\begin{array}{cccc|cc}
A_{1} & 0 & 0 & 0 & B_{1} & 0 \\
0 & A_{2} & 0 & 0 & 0 & B_{2} \\
0 & 0 & A_{3} & 0 & B_{3} & 0 \\
0 & 0 & 0 & A_{4} & 0 & B_{4} \\
\hline C_{1} & C_{2} & 0 & 0 & D_{1} & D_{2} \\
0 & 0 & C_{3} & C_{4} & D_{3} & D_{4}
\end{array}\right]
$$

Alternativamente, si la matriz de transferencia $G(s)$ puede factorizarse como producto y/o suma de varias matrices de transferencia con realizaciones simples, entonces una realización para $G$ puede obtenerse usando las fórmulas de cascada o suma dadas en la sección anterior.

Un problema inherente a estos procedimientos de realización es que, en general, la realización obtenida no será mínima. Para obtener una realización mínima, debe realizarse una descomposición de Kalman de controlabilidad y observabilidad para eliminar estados no controlables y/o no observables. (Un método alternativo numéricamente confiable para eliminar estados no controlables y/o no observables es la realización balanceada, que se discutirá más adelante).

Ahora describimos un procedimiento que sí da como resultado una realización mínima mediante expansión en fracciones parciales (la realización resultante se conoce a veces como realización de Gilbert, en honor a E. G. Gilbert).

Sea $G(s)$ una matriz de transferencia de tamaño $p \times m$ y escríbase en la forma

$$
G(s)=\frac{N(s)}{d(s)}
$$

con $d(s)$ un polinomio escalar. Para simplificar, supondremos que $d(s)$ tiene solo raíces reales y distintas, $\lambda_{i} \neq \lambda_{j}$ si $i \neq j$, y

$$
d(s)=\left(s-\lambda_{1}\right)\left(s-\lambda_{2}\right) \cdots\left(s-\lambda_{r}\right)
$$

Entonces $G(s)$ tiene la siguiente expansión en fracciones parciales:

$$
G(s)=D+\sum_{i=1}^{r} \frac{W_{i}}{s-\lambda_{i}}
$$

Suponga

$$
\operatorname{rank} W_{i}=k_{i}
$$

y sean $B_{i} \in \mathbb{R}^{k_{i} \times m}$ y $C_{i} \in \mathbb{R}^{p \times k_{i}}$ dos matrices constantes tales que

$$
W_{i}=C_{i} B_{i}
$$

Entonces una realización para $G(s)$ está dada por

$$
G(s)=\left[\begin{array}{ccc|c}
\lambda_{1} I_{k_{1}} & & & B_{1} \\
& \ddots & & \vdots \\
& & \lambda_{r} I_{k_{r}} & B_{r} \\
\hline C_{1} & \cdots & C_{r} & D
\end{array}\right]
$$

Se sigue inmediatamente de las pruebas PBH que esta realización es controlable y observable, y por tanto es mínima.

Una consecuencia inmediata de esta realización mínima es que una matriz de transferencia con denominador polinomial de orden $r$ no necesariamente tiene una realización en espacio de estados de orden $r$, a menos que cada $W_{i}$ sea de rango uno.

Este enfoque, de hecho, puede generalizarse a casos más complejos donde $d(s)$ tenga raíces complejas y/o repetidas. El lector puede convencerse probando algunos ejemplos simples.

\textbf{Comandos ilustrativos de MATLAB:}
$\gg \mathrm{G}=$ nd2sys(num, den, gain); $\mathrm{G}=$ zp2sys(zeros, poles, gain);

Comandos relacionados de MATLAB: ss2tf, ss2zp, zp2ss, tf2ss, residue, minreal

\section{Polos y ceros de sistemas multivariables}
Se llama matriz polinomial de una variable a una matriz cuyos elementos son polinomios de esa variable.

Definición 3.10 Sea $Q(s)$ una matriz polinomial (o matriz de transferencia) de tamaño $(p \times m)$ en $s$. Entonces el rango normal de $Q(s)$, denotado normalrank$(Q(s))$, es el máximo rango posible de $Q(s)$ para al menos un valor de $s \in \mathbb{C}$.

Para mostrar la diferencia entre el rango normal de una matriz polinomial y el rango de la matriz polinomial evaluada en un punto específico, considere

$$
Q(s)=\left[\begin{array}{cc}
s & 1 \\
s^{2} & 1 \\
s & 1
\end{array}\right]
$$

Entonces $Q(s)$ tiene rango normal 2 porque rank $Q(3)=2$. Sin embargo, $Q(0)$ tiene rango 1.

Los polos y ceros de una matriz de transferencia pueden caracterizarse en términos de sus realizaciones en espacio de estados. Sea

$$
\left[\begin{array}{c|c}
A & B \\
\hline C & D
\end{array}\right]
$$

una realización en espacio de estados de $G(s)$.

Definición 3.11 Los valores propios de $A$ se llaman polos de la realización de $G(s)$.

Para definir ceros, consideremos la siguiente matriz del sistema:

$$
Q(s)=\left[\begin{array}{cc}
A-s I & B \\
C & D
\end{array}\right]
$$

Definición 3.12 Un número complejo $z_{0} \in \mathbb{C}$ se llama cero invariante de la realización del sistema si satisface

$$
\operatorname{rank}\left[\begin{array}{cc}
A-z_{0} I & B \\
C & D
\end{array}\right]<\text { normalrank }\left[\begin{array}{cc}
A-s I & B \\
C & D
\end{array}\right]
$$

Los ceros invariantes no cambian bajo realimentación de estado constante, ya que

$$
\begin{aligned}
\operatorname{rank}\left[\begin{array}{cc}
A+B F-z_{0} I & B \\
C+D F & D
\end{array}\right] & =\operatorname{rank}\left[\begin{array}{cc}
A-z_{0} I & B \\
C & D
\end{array}\right]\left[\begin{array}{cc}
I & 0 \\
F & I
\end{array}\right] \\
& =\operatorname{rank}\left[\begin{array}{cc}
A-z_{0} I & B \\
C & D
\end{array}\right]
\end{aligned}
$$

También es claro que los ceros invariantes no cambian bajo transformación de similitud.

El siguiente lema es evidente.

Lema 3.7 Suponga que $\left[\begin{array}{cc}A-s I & B \\ C & D\end{array}\right]$ tiene rango normal completo por columnas. Entonces $z_{0} \in \mathbb{C}$ es un cero invariante de una realización $(A, B, C, D)$ si y solo si existen $0 \neq x \in \mathbb{C}^{n}$ y $u \in \mathbb{C}^{m}$ tales que

$$
\left[\begin{array}{cc}
A-z_{0} I & B \\
C & D
\end{array}\right]\left[\begin{array}{l}
x \\
u
\end{array}\right]=0
$$

Además, si $u=0$, entonces $z_{0}$ también es un modo no observable.

Prueba. Por definición, $z_{0}$ es un cero invariante si existe un vector $\left[\begin{array}{l}x \\ u\end{array}\right] \neq 0$ tal que

$$
\left[\begin{array}{cc}
A-z_{0} I & B \\
C & D
\end{array}\right]\left[\begin{array}{l}
x \\
u
\end{array}\right]=0
$$

ya que $\left[\begin{array}{cc}A-s I & B \\ C & D\end{array}\right]$ tiene rango normal completo por columnas.

Por otro lado, suponga que $z_{0}$ es un cero invariante; entonces existe un vector $\left[\begin{array}{l}x \\ u\end{array}\right] \neq 0$ tal que

$$
\left[\begin{array}{cc}
A-z_{0} I & B \\
C & D
\end{array}\right]\left[\begin{array}{l}
x \\
u
\end{array}\right]=0
$$

Afirmamos que $x \neq 0$. De lo contrario, $\left[\begin{array}{l}B \\ D\end{array}\right] u=0$ o $u=0$, ya que $\left[\begin{array}{cc}A-s I & B \\ C & D\end{array}\right]$ tiene rango normal completo por columnas (es decir, $\left[\begin{array}{l}x \\ u\end{array}\right]=0$), lo cual es una contradicción.

Finalmente, nótese que si $u=0$, entonces

$$
\left[\begin{array}{c}
A-z_{0} I \\
C
\end{array}\right] x=0
$$

y $z_{0}$ es un modo no observable por la prueba PBH.

Cuando el sistema es cuadrado, los ceros invariantes pueden calcularse resolviendo un problema de valores propios generalizado:

$$
\underbrace{\left[\begin{array}{ll}
A & B \\
C & D
\end{array}\right]}_{M}\left[\begin{array}{l}
x \\
u
\end{array}\right]=z_{0} \underbrace{\left[\begin{array}{ll}
I & 0 \\
0 & 0
\end{array}\right]}_{N}\left[\begin{array}{l}
x \\
u
\end{array}\right]
$$

usando el comando de Matlab: $\operatorname{eig}(\mathrm{M}, \mathbf{N})$.

Lema 3.8 Suponga que $\left[\begin{array}{cc}A-s I & B \\ C & D\end{array}\right]$ tiene rango normal completo por filas. Entonces $z_{0} \in \mathbb{C}$ es un cero invariante de una realización $(A, B, C, D)$ si y solo si existen $0 \neq y \in \mathbb{C}^{n}$ y $v \in \mathbb{C}^{p}$ tales que

$$
\left[\begin{array}{ll}
y^{*} & v^{*}
\end{array}\right]\left[\begin{array}{cc}
A-z_{0} I & B \\
C & D
\end{array}\right]=0
$$

Además, si $v=0$, entonces $z_{0}$ también es un modo no controlable.

Lema 3.9 $G(s)$ tiene rango normal completo por columnas (filas) si y solo si $\left[\begin{array}{cc}A-s I & B \\ C & D\end{array}\right]$ tiene rango normal completo por columnas (filas).

Prueba. Esto se sigue al notar que

$$
\left[\begin{array}{cc}
A-s I & B \\
C & D
\end{array}\right]=\left[\begin{array}{cc}
I & 0 \\
C(A-s I)^{-1} & I
\end{array}\right]\left[\begin{array}{cc}
A-s I & B \\
0 & G(s)
\end{array}\right]
$$

y

$$
\text { normalrank }\left[\begin{array}{cc}
A-s I & B \\
C & D
\end{array}\right]=n+\operatorname{normalrank}(G(s)) \text {. }
$$

Lema 3.10 Sea $G(s) \in \mathcal{R}_{p}(s)$ una matriz de transferencia de tamaño $p \times m$ y sea $(A, B, C, D)$ una realización mínima. Si la entrada del sistema es de la forma $u(t)=u_{0} e^{\lambda t}$, donde $\lambda \in \mathbb{C}$ no es un polo de $G(s)$ y $u_{0} \in \mathbb{C}^{m}$ es un vector constante arbitrario, entonces la salida debida a la entrada $u(t)$ y al estado inicial $x(0)=(\lambda I-A)^{-1} B u_{0}$ es $y(t)=G(\lambda) u_{0} e^{\lambda t}, \forall t \geq 0$. En particular, si $\lambda$ es un cero de $G(s)$, entonces $y(t)=0$.

Prueba. La respuesta del sistema con respecto a la entrada $u(t)=u_{0} e^{\lambda t}$ y la condición inicial $x(0)=(\lambda I-A)^{-1} B u_{0}$ es (en términos de la transformada de Laplace)

$$
\begin{aligned}
Y(s)= & C(s I-A)^{-1} x(0)+C(s I-A)^{-1} B U(s)+D U(s) \\
= & C(s I-A)^{-1} x(0)+C(s I-A)^{-1} B u_{0}(s-\lambda)^{-1}+D u_{0}(s-\lambda)^{-1} \\
= & C(s I-A)^{-1} x(0)+C\left[(s I-A)^{-1}-(\lambda I-A)^{-1}\right] B u_{0}(s-\lambda)^{-1} \\
& \quad+C(\lambda I-A)^{-1} B u_{0}(s-\lambda)^{-1}+D u_{0}(s-\lambda)^{-1} \\
= & C(s I-A)^{-1}\left(x(0)-(\lambda I-A)^{-1} B u_{0}\right)+G(\lambda) u_{0}(s-\lambda)^{-1} \\
= & G(\lambda) u_{0}(s-\lambda)^{-1}
\end{aligned}
$$

Por lo tanto, $y(t)=G(\lambda) u_{0} e^{\lambda t}$.

Ejemplo 3.3 Sea

$$
G(s)=\left[\begin{array}{c|c}
A & B \\
\hline C & D
\end{array}\right]=\left[\begin{array}{ccc|ccc}
-1 & -2 & 1 & 1 & 2 & 3 \\
0 & 2 & -1 & 3 & 2 & 1 \\
-4 & -3 & -2 & 1 & 1 & 1 \\
\hline 1 & 1 & 1 & 0 & 0 & 0 \\
2 & 3 & 4 & 0 & 0 & 0
\end{array}\right]
$$

Entonces los ceros invariantes del sistema pueden encontrarse usando el comando de Matlab

$$
\begin{aligned}
& \gg \mathbf{G}=\operatorname{pck}(\mathbf{A}, \mathbf{B}, \mathbf{C}, \mathbf{D}), \quad \mathrm{z}_{0}=\operatorname{szeros}(\mathbf{G}), \% \text { o } \\
& \gg \mathbf{z}_{0}=\operatorname{tzero}(\mathbf{A}, \mathbf{B}, \mathbf{C}, \mathbf{D})
\end{aligned}
$$

que da $z_{0}=0.2$. Como $G(s)$ es de rango completo por filas, podemos encontrar $y$ y $v$ tales que

$$
\left[\begin{array}{ll}
y^{*} & v^{*}
\end{array}\right]\left[\begin{array}{cc}
A-z_{0} I & B \\
C & D
\end{array}\right]=0
$$

lo cual también puede calcularse con un comando de MATLAB:

$$
\gg \operatorname{null}\left(\left[\mathbf{A}-\mathbf{z}_{\mathbf{0}} * \operatorname{eye}(\mathbf{3}), \mathbf{B} ; \mathbf{C}, \mathbf{D}\right]'\right) \Longrightarrow\left[\begin{array}{l}
y \\
v
\end{array}\right]=\left[\begin{array}{c}
0.0466 \\
0.0466 \\
-0.1866 \\
--0.9702 \\
0.1399
\end{array}\right]
$$

Comandos relacionados de MATLAB: spoles, rifd

\section{Notas y referencias}
Se remite al lector a Brogan [1991], Chen [1984], Kailath [1980] y Wonham [1985] para un tratamiento amplio de la teoría estándar de sistemas lineales.

\section{Problemas}
Problema 3.1 Sean $A \in \mathbb{C}^{n \times n}$ y $B \in \mathbb{C}^{m \times m}$. Demuestre que $X(t)=e^{A t} X(0) e^{B t}$ es la solución de

$$
\dot{X}=A X+X B
$$

Problema 3.2 Dada la respuesta al impulso, $h_{1}(t)$, de un sistema lineal invariante en el tiempo con modelo

$$
\ddot{y}+a_{1} \dot{y}+a_{2} y=u
$$

encuentre la respuesta al impulso, $h(t)$, del sistema

$$
\ddot{y}+a_{1} \dot{y}+a_{2} y=b_{0} \ddot{u}+b_{1} \dot{u}+b_{2} u
$$

Justifique su respuesta y generalice su resultado a sistemas de orden $n$.

Problema 3.3 Suponga que un sistema de segundo orden está dado por $\dot{x}=A x$. Suponga que se conoce que $x(1)=\left[\begin{array}{c}4 \\ -2\end{array}\right]$ para $x(0)=\left[\begin{array}{l}1 \\ 1\end{array}\right]$ y $x(1)=\left[\begin{array}{c}5 \\ -2\end{array}\right]$ para $x(0)=\left[\begin{array}{l}1 \\ 2\end{array}\right]$. Encuentre $x(n)$ para $x(0)=\left[\begin{array}{l}1 \\ 0\end{array}\right]$. ¿Puede determinar $A$?

Problema 3.4 Suponga que $(A, B)$ es controlable. Demuestre que $(F, G)$ con

$$
F=\left[\begin{array}{ll}
A & 0 \\
C & 0
\end{array}\right], \quad G=\left[\begin{array}{l}
B \\
0
\end{array}\right]
$$

es controlable si y solo si

$$
\left[\begin{array}{ll}
A & B \\
C & 0
\end{array}\right]
$$

es una matriz de rango completo por filas.

Problema 3.5 Sean $\lambda_{i}, i=1,2, \ldots, n$ los $n$ valores propios distintos de una matriz $A \in \mathbb{C}^{n \times n}$ y sean $x_{i}$ y $y_{i}$ los vectores propios unitarios derecho e izquierdo correspondientes, tales que $y_{i}^{*} x_{i}=1$. Demuestre que

$$
y_{i}^{*} x_{j}=\delta_{i j}, \quad A=\sum_{i=1}^{n} \lambda_{i} x_{i} y_{i}^{*}
$$

y

$$
C(s I-A)^{-1} B=\sum_{i=1}^{n} \frac{\left(C x_{i}\right)\left(y_{i}^{*} B\right)}{s-\lambda_{i}}
$$

Además, demuestre que el modo $\lambda_{i}$ es controlable si y solo si $y_{i}^{*} B \neq 0$, y el modo $\lambda_{i}$ es observable si y solo si $C x_{i} \neq 0$.

Problema 3.6 Sea $(A, b, c)$ una realización con $A \in \mathbb{R}^{n \times n}, b \in R^{n}$ y $c^{*} \in \mathbb{R}^{n}$. Suponga que $\lambda_{i}(A)+\lambda_{j}(A) \neq 0$ para todo $i, j$. Esta hipótesis garantiza la existencia de $X \in \mathbb{R}^{n \times n}$ tal que

$$
A X+X A+b c=0
$$

Demuestre que $X$ es no singular si y solo si $(A, b)$ es controlable y $(c, A)$ es observable.

Problema 3.7 Calcule los ceros del sistema y las direcciones de cero correspondientes de las siguientes funciones de transferencia

$$
\begin{gathered}
G_{1}(s)=\left[\begin{array}{ll|ll}
1 & 2 & 2 & 2 \\
1 & 0 & 3 & 4 \\
\hline 0 & 1 & 1 & 2 \\
1 & 0 & 2 & 0
\end{array}\right], \quad G_{2}(s)=\left[\begin{array}{cc|cc}
-1 & -2 & 1 & 2 \\
0 & 1 & 2 & 1 \\
\hline 1 & 1 & 0 & 0 \\
1 & 1 & 0 & 0
\end{array}\right], \\
G_{3}(s)=\left[\begin{array}{ll}
\frac{2(s+1)(s+2)}{s(s+3)(s+4)} & \frac{s+2}{(s+1)(s+3)}
\end{array}\right], \quad G_{4}(s)=\left[\begin{array}{cc}
\frac{1}{s+1} & \frac{s+3}{(s+1)(s-2)} \\
\frac{10}{s-2} & \frac{5}{s+3}
\end{array}\right]
\end{gathered}
$$

También encuentre los vectores $x$ y $u$ cuando corresponda, de modo que se cumpla

$$
\left[\begin{array}{cc}
A-z I & B \\
C & D
\end{array}\right]\left[\begin{array}{l}
x \\
u
\end{array}\right]=0 \text { o }\left[\begin{array}{ll}
x^{*} & u^{*}
\end{array}\right]\left[\begin{array}{cc}
A-z I & B \\
C & D
\end{array}\right]=0
$$

\chapter{Espacios $\mathcal{H}_{2}$ y $\mathcal{H}_{\infty}$}
El objetivo más importante de un sistema de control es cumplir ciertas especificaciones de desempeño, además de proporcionar estabilidad interna. Una forma de describir las especificaciones de desempeño de un sistema de control es en términos del tamaño de ciertas señales de interés. Por ejemplo, el desempeño de un sistema de seguimiento podría medirse mediante el tamaño de la señal de error de seguimiento. En este capítulo analizamos varias formas de definir el tamaño de una señal (es decir, varias normas para señales). Por supuesto, la norma más apropiada depende de la situación concreta. Con ese propósito, primero introduciremos los espacios de Hardy $\mathcal{H}_{2}$ y $\mathcal{H}_{\infty}$. También se presentan algunos métodos en espacio de estados para calcular normas de matrices de transferencia racionales reales en $\mathcal{H}_{2}$ y $\mathcal{H}_{\infty}$.

\section{Espacios de Hilbert}
Recordemos el producto interno de vectores definido en un espacio euclidiano $\mathbb{C}^{n}$:

$$
\langle x, y\rangle:=x^{*} y=\sum_{i=1}^{n} \bar{x}_{i} y_{i} \quad \forall x=\left[\begin{array}{c}
x_{1} \\
\vdots \\
x_{n}
\end{array}\right], y=\left[\begin{array}{c}
y_{1} \\
\vdots \\
y_{n}
\end{array}\right] \in \mathbb{C}^{n}
$$

Nótese que muchas nociones métricas importantes y propiedades geométricas, como longitud, distancia, ángulo y energía de sistemas físicos, pueden deducirse de este producto interno. Por ejemplo, la longitud de un vector $x \in \mathbb{C}^{n}$ se define como

$$
\|x\|:=\sqrt{\langle x, x\rangle}
$$

y el ángulo entre dos vectores $x, y \in \mathbb{C}^{n}$ puede calcularse mediante

$$
\cos \angle(x, y)=\frac{\langle x, y\rangle}{\|x\|\|y\|}, \quad \angle(x, y) \in[0, \pi]
$$

Se dice que dos vectores son ortogonales si $\angle(x, y)=\frac{\pi}{2}$.

Ahora consideramos una generalización natural del producto interno en $\mathbb{C}^{n}$ hacia espacios vectoriales más generales (posiblemente de dimensión infinita).

Definición 4.1 Sea $V$ un espacio vectorial sobre $\mathbb{C}$. Un producto interno ${ }^{1}$ en $V$ es una función de valores complejos,

$$
\langle\cdot, \cdot\rangle: V \times V \longmapsto \mathbb{C}
$$

tal que para cualesquiera $x, y, z \in V$ y $\alpha, \beta \in \mathbb{C}$

(i) $\langle x, \alpha y+\beta z\rangle=\alpha\langle x, y\rangle+\beta\langle x, z\rangle$

(ii) $\langle x, y\rangle=\overline{\langle y, x\rangle}$

(iii) $\langle x, x\rangle>0$ si $x \neq 0$.

Un espacio vectorial $V$ con un producto interno se llama espacio con producto interno.

Es claro que el producto interno definido arriba induce una norma $\|x\|:=\sqrt{\langle x, x\rangle}$, por lo que se satisfacen las condiciones de norma del Capítulo 2. En particular, la distancia entre vectores $x$ y $y$ es $d(x, y)=\|x-y\|$.

Dos vectores $x$ y $y$ en un espacio con producto interno $V$ se dicen ortogonales si $\langle x, y\rangle=0$, denotado $x \perp y$. Más generalmente, un vector $x$ se dice ortogonal a un conjunto $S \subset V$, denotado por $x \perp S$, si $x \perp y$ para todo $y \in S$.

El producto interno y la norma inducida por el producto interno tienen las siguientes propiedades familiares.

Teorema 4.1 Sea $V$ un espacio con producto interno y sean $x, y \in V$. Entonces

(i) $|\langle x, y\rangle| \leq\|x\|\|y\|$ (desigualdad de Cauchy-Schwarz). Además, la igualdad se cumple si y solo si $x=\alpha y$ para alguna constante $\alpha$ o $y=0$.

(ii) $\|x+y\|^{2}+\|x-y\|^{2}=2\|x\|^{2}+2\|y\|^{2}$ (ley del paralelogramo).

(iii) $\|x+y\|^{2}=\|x\|^{2}+\|y\|^{2}$ si $x \perp y$.

Un espacio de Hilbert es un espacio con producto interno completo con la norma inducida por su producto interno. Por ejemplo, $\mathbb{C}^{n}$ con el producto interno usual es un espacio de Hilbert (de dimensión finita). Más generalmente, es directo verificar que $\mathbb{C}^{n \times m}$ con el producto interno definido como

$$
\langle A, B\rangle:=\operatorname{trace} A^{*} B=\sum_{i=1}^{n} \sum_{j=1}^{m} \bar{a}_{i j} b_{i j} \quad \forall A, B \in \mathbb{C}^{n \times m}
$$

también es un espacio de Hilbert (de dimensión finita).

Un conocido espacio de Hilbert de dimensión infinita es $\mathcal{L}_{2}[a, b]$, que consiste en todas las funciones cuadrado integrables y medibles de Lebesgue definidas en un intervalo $[a, b]$, con producto interno definido como

$$
\langle f, g\rangle:=\int_{a}^{b} f(t)^{*} g(t) d t
$$

[\^{}3]para $f, g \in \mathcal{L}_{2}[a, b]$. De forma similar, si las funciones son vectoriales o matriciales, el producto interno se define correspondientemente como

$$
\langle f, g\rangle:=\int_{a}^{b} \operatorname{trace}\left[f(t)^{*} g(t)\right] d t
$$

Algunos espacios usados frecuentemente en este libro son $\mathcal{L}_{2}[0, \infty), \mathcal{L}_{2}(-\infty, 0], \mathcal{L}_{2}(-\infty, \infty)$. Más precisamente, se definen como

$\mathcal{L}_{2}=\mathcal{L}_{2}(-\infty, \infty)$: espacio de Hilbert de funciones matriciales sobre $\mathbb{R}$, con producto interno

$$
\langle f, g\rangle:=\int_{-\infty}^{\infty} \operatorname{trace}\left[f(t)^{*} g(t)\right] d t
$$

$\mathcal{L}_{2+}=\mathcal{L}_{2}[0, \infty)$: subespacio de $\mathcal{L}_{2}(-\infty, \infty)$ con funciones nulas para $t<0$.

$\mathcal{L}_{2-}=\mathcal{L}_{2}(-\infty, 0]$: subespacio de $\mathcal{L}_{2}(-\infty, \infty)$ con funciones nulas para $t>0$.

\section{Espacios $\mathcal{H}_{2}$ y $\mathcal{H}_{\infty}$}
Sea $S \subset \mathbb{C}$ un conjunto abierto, y sea $f(s)$ una función compleja definida en $S$:

$$
f(s): S \longmapsto \mathbb{C}
$$

Entonces se dice que $f(s)$ es analítica en un punto $z_{0}$ de $S$ si es derivable en $z_{0}$ y también en cada punto de algún entorno de $z_{0}$. Es un hecho que, si $f(s)$ es analítica en $z_{0}$, entonces $f$ tiene derivadas continuas de todos los órdenes en $z_{0}$. Por lo tanto, una función analítica en $z_{0}$ tiene una representación en serie de potencias en $z_{0}$. El recíproco también es cierto (es decir, si una función tiene una serie de potencias en $z_{0}$, entonces es analítica en $z_{0}$). Se dice que una función $f(s)$ es analítica en $S$ si tiene derivada o es analítica en cada punto de $S$. Una función matricial es analítica en $S$ si cada elemento de la matriz es analítico en $S$. Por ejemplo, todas las matrices de transferencia racionales reales estables son analíticas en el semiplano derecho y $e^{-s}$ es analítica en todo el plano.

Una propiedad bien conocida de las funciones analíticas es el llamado teorema del módulo máximo.

Teorema 4.2 Si $f(s)$ está definida y es continua en un conjunto cerrado y acotado $S$, y es analítica en el interior de $S$, entonces $|f(s)|$ no puede alcanzar su máximo en el interior de $S$ a menos que $f(s)$ sea constante.

El teorema implica que $|f(s)|$ solo puede alcanzar su máximo en la frontera de $S$; es decir,

$$
\max _{s \in S}|f(s)|=\max _{s \in \partial S}|f(s)|
$$

donde $\partial S$ denota la frontera de $S$. A continuación consideramos algunos espacios de funciones complejas (matriciales) de uso frecuente.

\subsection*{Espacio $\mathcal{L}_{2}(j \mathbb{R})$}
$\mathcal{L}_{2}(j \mathbb{R})$, o simplemente $\mathcal{L}_{2}$, es un espacio de Hilbert de funciones matriciales (o escalares) sobre $j \mathbb{R}$ y consiste en todas las funciones matriciales complejas $F$ tales que la siguiente integral está acotada:

$$
\int_{-\infty}^{\infty} \operatorname{trace}\left[F^{*}(j \omega) F(j \omega)\right] d \omega<\infty
$$

El producto interno para este espacio de Hilbert se define como

$$
\langle F, G\rangle:=\frac{1}{2 \pi} \int_{-\infty}^{\infty} \operatorname{trace}\left[F^{*}(j \omega) G(j \omega)\right] d \omega
$$

para $F, G \in \mathcal{L}_{2}$, y la norma inducida por este producto interno está dada por

$$
\|F\|_{2}:=\sqrt{\langle F, F\rangle}
$$

Por ejemplo, todas las matrices de transferencia racionales reales estrictamente propias sin polos sobre el eje imaginario forman un subespacio (no cerrado) de $\mathcal{L}_{2}(j \mathbb{R})$ que se denota por $\mathcal{R} \mathcal{L}_{2}(j \mathbb{R})$ o simplemente $\mathcal{R} \mathcal{L}_{2}$.

\subsection*{Espacio $\mathcal{H}_{2}$ ${ }^{2}$}
$\mathcal{H}_{2}$ es un subespacio (cerrado) de $\mathcal{L}_{2}(j \mathbb{R})$ con funciones matriciales $F(s)$ analíticas en $\operatorname{Re}(s)>0$ (semiplano derecho abierto). La norma correspondiente se define como

$$
\|F\|_{2}^{2}:=\sup _{\sigma>0}\left\{\frac{1}{2 \pi} \int_{-\infty}^{\infty} \operatorname{trace}\left[F^{*}(\sigma+j \omega) F(\sigma+j \omega)\right] d \omega\right\}
$$

Puede demostrarse ${ }^{3}$ que

$$
\|F\|_{2}^{2}=\frac{1}{2 \pi} \int_{-\infty}^{\infty} \operatorname{trace}\left[F^{*}(j \omega) F(j \omega)\right] d \omega
$$

Por lo tanto, podemos calcular la norma en $\mathcal{H}_{2}$ igual que en $\mathcal{L}_{2}$. El subespacio racional real de $\mathcal{H}_{2}$, que consiste en todas las matrices de transferencia racionales reales estables y estrictamente propias, se denota por $\mathcal{R H}_{2}$.

\subsection*{Espacio $\mathcal{H}_{2}^{\perp}$}
$\mathcal{H}_{2}^{\perp}$ es el complemento ortogonal de $\mathcal{H}_{2}$ en $\mathcal{L}_{2}$; es decir, el subespacio (cerrado) de funciones en $\mathcal{L}_{2}$ que son analíticas en el semiplano izquierdo abierto. El subespacio racional real de $\mathcal{H}_{2}^{\perp}$, que consiste en todas las matrices de transferencia racionales estrictamente propias con todos sus polos en el semiplano derecho abierto, se denotará por $\mathcal{R} \mathcal{H}_{2}^{\perp}$. Es fácil ver que si $G$ es una matriz de transferencia racional real, estable y estrictamente propia, entonces $G \in \mathcal{H}_{2}$ y $G^{\sim} \in \mathcal{H}_{2}^{\perp}$. La mayor parte de nuestro estudio en este libro se centrará en el caso racional real.[\^{}4]

Los espacios $\mathcal{L}_{2}$ definidos previamente en el dominio de la frecuencia pueden relacionarse con los espacios $\mathcal{L}_{2}$ definidos en el dominio del tiempo. Recordemos que una función en el espacio $\mathcal{L}_{2}$ del dominio del tiempo admite transformada bilateral de Laplace (o de Fourier). De hecho, puede demostrarse que esta transformada bilateral de Laplace produce un isomorfismo isométrico entre los espacios $\mathcal{L}_{2}$ en el dominio del tiempo y los espacios $\mathcal{L}_{2}$ en el dominio de la frecuencia (esto es lo que se conoce como relaciones de Parseval):

$$
\begin{aligned}
\mathcal{L}_{2}(-\infty, \infty) & \cong \mathcal{L}_{2}(j \mathbb{R}) \\
\mathcal{L}_{2}[0, \infty) & \cong \mathcal{H}_{2} \\
\mathcal{L}_{2}(-\infty, 0] & \cong \mathcal{H}_{2}^{\perp}
\end{aligned}
$$

Como resultado, si $g(t) \in \mathcal{L}_{2}(-\infty, \infty)$ y su transformada bilateral de Laplace es $G(s) \in \mathcal{L}_{2}(j \mathbb{R})$, entonces

$$
\|G\|_{2}=\|g\|_{2}
$$

Por lo tanto, cuando no haya confusión, la notación para funciones en el dominio del tiempo y en el dominio de la frecuencia se usará indistintamente.

\insertimage[\label{img:espacios-relaciones-funciones}]{2024_06_29_ce7c73b22613114bd4d6g-063_765_822_1065_627}{width=0.45\textwidth}{Relaciones entre espacios de funciones.}

Defina una proyección ortogonal

$$
P_{+}: \mathcal{L}_{2}(-\infty, \infty) \longmapsto \mathcal{L}_{2}[0, \infty)
$$

tal que, para cualquier función $f(t) \in \mathcal{L}_{2}(-\infty, \infty)$, tenemos $g(t)=P_{+} f(t)$ con

$$
g(t):=\left\{\begin{array}{cl}
f(t), & \text { para } t \geq 0 \\
0, & \text { para } t<0
\end{array}\right.
$$

En este libro, $P_{+}$ también se usará para denotar la proyección de $\mathcal{L}_{2}(j \mathbb{R})$ sobre $\mathcal{H}_{2}$. De forma similar, definimos $P_{-}$ como otra proyección ortogonal de $\mathcal{L}_{2}(-\infty, \infty)$ sobre $\mathcal{L}_{2}(-\infty, 0]$ (o de $\mathcal{L}_{2}(j \mathbb{R})$ sobre $\mathcal{H}_{2}^{\perp}$). Entonces, las relaciones entre espacios $\mathcal{L}_{2}$ y espacios $\mathcal{H}_{2}$ se muestran en la Figura 4.1.

Otra clase importante de funciones matriciales complejas usadas en este libro son las funciones acotadas en el eje imaginario.

\subsection*{Espacio $\mathcal{L}_{\infty}(j \mathbb{R})$}
$\mathcal{L}_{\infty}(j \mathbb{R})$, o simplemente $\mathcal{L}_{\infty}$, es un espacio de Banach de funciones matriciales (o escalares) que están (esencialmente) acotadas en $j \mathbb{R}$, con norma

$$
\|F\|_{\infty}:=\underset{\omega \in \mathbb{R}}{\operatorname{ess} \sup } \bar{\sigma}[F(j \omega)] .
$$

El subespacio racional de $\mathcal{L}_{\infty}$, denotado por $\mathcal{R} \mathcal{L}_{\infty}(j \mathbb{R})$ o simplemente $\mathcal{R} \mathcal{L}_{\infty}$, consiste en todas las matrices de transferencia racionales reales propias sin polos en el eje imaginario.

\subsection*{Espacio $\mathcal{H}_{\infty}$}
$\mathcal{H}_{\infty}$ es un subespacio (cerrado) de $\mathcal{L}_{\infty}$ con funciones que son analíticas y acotadas en el semiplano derecho abierto. La norma $\mathcal{H}_{\infty}$ se define como

$$
\|F\|_{\infty}:=\sup _{\operatorname{Re}(s)>0} \bar{\sigma}[F(s)]=\sup _{\omega \in \mathbb{R}} \bar{\sigma}[F(j \omega)] .
$$

La segunda igualdad puede verse como una generalización del teorema del módulo máximo para funciones matriciales. (Véase Boyd y Desoer [1985] para una demostración.) El subespacio racional real de $\mathcal{H}_{\infty}$ se denota por $\mathcal{R} \mathcal{H}_{\infty}$, y consiste en todas las matrices de transferencia racionales reales propias y estables.

\subsection*{Espacio $\mathcal{H}_{\infty}^{-}$}
$\mathcal{H}_{\infty}^{-}$ es un subespacio (cerrado) de $\mathcal{L}_{\infty}$ con funciones que son analíticas y acotadas en el semiplano izquierdo abierto. La norma $\mathcal{H}_{\infty}^{-}$ se define como

$$
\|F\|_{\infty}:=\sup _{\operatorname{Re}(s)<0} \bar{\sigma}[F(s)]=\sup _{\omega \in \mathbb{R}} \bar{\sigma}[F(j \omega)] .
$$

El subespacio racional real de $\mathcal{H}_{\infty}^{-}$ se denota por $\mathcal{R H}_{\infty}^{-}$, y consiste en todas las matrices de transferencia propias, racionales reales y antiestables (es decir, funciones con todos los polos en el semiplano derecho abierto).

Sea $G(s) \in \mathcal{L}_{\infty}$ una matriz de transferencia $p \times q$. Entonces se define un operador de multiplicación como

$$
\begin{aligned}
M_{G}: \mathcal{L}_{2} \longmapsto \mathcal{L}_{2} \\
M_{G} f:=G f .
\end{aligned}
$$

Al escribir el mapeo anterior, hemos supuesto que $f$ tiene dimensión compatible. Una descripción más precisa del operador anterior sería

$$
M_{G}: \mathcal{L}_{2}^{q} \longmapsto \mathcal{L}_{2}^{p}
$$

Es decir, $f$ es una función vectorial de dimensión $q$, con cada componente en $\mathcal{L}_{2}$. Sin embargo, en este libro suprimiremos todas las dimensiones y asumiremos que todos los objetos tienen dimensiones compatibles.

Un hecho útil sobre el operador de multiplicación es que la norma de una matriz $G$ en $\mathcal{L}_{\infty}$ es igual a la norma del operador de multiplicación correspondiente.

Teorema 4.3 Sea $G \in \mathcal{L}_{\infty}$ una matriz de transferencia $p \times q$. Entonces $\left\|M_{G}\right\|=\|G\|_{\infty}$.

Observación 4.1 También es cierto que esta norma de operador es igual a la norma del operador restringido a $\mathcal{H}_{2}$ (o $\mathcal{H}_{2}^{\perp}$); es decir,

$$
\left\|M_{G}\right\|=\left\|M_{G} \mid \mathcal{H}_{2}\right\|:=\sup \left\{\|G f\|_{2}: f \in \mathcal{H}_{2},\|f\|_{2} \leq 1\right\}
$$

Esto quedará claro en la prueba, donde se construye un $f \in \mathcal{H}_{2}$.

Prueba. Por definición,

$$
\left\|M_{G}\right\|=\sup \left\{\|G f\|_{2}: f \in \mathcal{L}_{2},\|f\|_{2} \leq 1\right\}
$$

Primero vemos que $\|G\|_{\infty}$ es una cota superior para la norma del operador:

$$
\begin{aligned}
\|G f\|_{2}^{2} & =\frac{1}{2 \pi} \int_{-\infty}^{\infty} f^{*}(j \omega) G^{*}(j \omega) G(j \omega) f(j \omega) d \omega \\
& \leq\|G\|_{\infty}^{2} \frac{1}{2 \pi} \int_{-\infty}^{\infty}\|f(j \omega)\|^{2} d \omega \\
& =\|G\|_{\infty}^{2}\|f\|_{2}^{2}
\end{aligned}
$$

Para mostrar que $\|G\|_{\infty}$ es la menor cota superior, primero elegimos una frecuencia $\omega_{0}$ donde $\bar{\sigma}[G(j \omega)]$ sea máxima; es decir,

$$
\bar{\sigma}\left[G\left(j \omega_{0}\right)\right]=\|G\|_{\infty}
$$

y denotamos la descomposición en valores singulares de $G\left(j \omega_{0}\right)$ por

$$
G\left(j \omega_{0}\right)=\bar{\sigma} u_{1}\left(j \omega_{0}\right) v_{1}^{*}\left(j \omega_{0}\right)+\sum_{i=2}^{r} \sigma_{i} u_{i}\left(j \omega_{0}\right) v_{i}^{*}\left(j \omega_{0}\right)
$$

donde $r$ es el rango de $G\left(j \omega_{0}\right)$ y $u_{i}, v_{i}$ tienen norma unitaria.

Luego suponemos que $G(s)$ tiene coeficientes reales y construiremos una función $f(s) \in \mathcal{H}_{2}$ con coeficientes reales de modo que la norma se alcance aproximadamente. [Quedará claro a continuación que la prueba es mucho más simple si se permite que $f$ tenga coeficientes complejos, lo cual es necesario cuando $G(s)$ tiene coeficientes complejos.]

Si $\omega_{0}<\infty$, escribimos $v_{1}\left(j \omega_{0}\right)$ como

$$
v_{1}\left(j \omega_{0}\right)=\left[\begin{array}{c}
\alpha_{1} e^{j \theta_{1}} \\
\alpha_{2} e^{j \theta_{2}} \\
\vdots \\
\alpha_{q} e^{j \theta_{q}}
\end{array}\right]
$$

donde $\alpha_{i} \in \mathbb{R}$ y $\theta_{i} \in(-\pi, 0]$, y $q$ es la dimensión de columna de $G$. Ahora sea $0 \leq \beta_{i} \leq \infty$ tal que

$$
\theta_{i}=\angle\left(\frac{\beta_{i}-j \omega_{0}}{\beta_{i}+j \omega_{0}}\right)
$$

(con $\beta_{i} \rightarrow \infty$ si $\theta_{i}=0$) y sea $f$ dada por

$$
f(s)=\left[\begin{array}{c}
\alpha_{1} \frac{\beta_{1}-s}{\beta_{1}+s} \\
\alpha_{2} \frac{\beta_{2}-s}{\beta_{2}+s} \\
\vdots \\
\alpha_{q} \frac{\beta_{q}-s}{\beta_{q}+s}
\end{array}\right] \hat{f}(s)
$$

(con 1 en lugar de $\frac{\beta_{i}-s}{\beta_{i}+s}$ si $\theta_{i}=0$), donde una función escalar $\hat{f}$ se elige de forma que

$$
|\hat{f}(j \omega)|= \begin{cases}c & \text { si }\left|\omega-\omega_{0}\right|<\epsilon \text { o }\left|\omega+\omega_{0}\right|<\epsilon \\
0 & \text { en otro caso }\end{cases}
$$

donde $\epsilon$ es un número positivo pequeño y $c$ se elige para que $\hat{f}$ tenga norma 2 unitaria (es decir, $c=\sqrt{\pi / 2 \epsilon}$). Esto, a su vez, implica que $f$ tiene norma 2 unitaria. Entonces

$$
\begin{aligned}
\|G f\|_{2}^{2} & \approx \frac{1}{2 \pi}\left[\bar{\sigma}\left[G\left(-j \omega_{0}\right)\right]^{2} \pi+\bar{\sigma}\left[G\left(j \omega_{0}\right)\right]^{2} \pi\right] \\
& =\bar{\sigma}\left[G\left(j \omega_{0}\right)\right]^{2}=\|G\|_{\infty}^{2}
\end{aligned}
$$

De manera similar, si $\omega_{0}=\infty$, la conclusión se obtiene haciendo $\omega_{0} \rightarrow \infty$ en lo anterior.

\textbf{Comandos ilustrativos de MATLAB:}
$\gg[\mathbf{s v}, \mathbf{w}]=\operatorname{sigma}(\mathbf{A}, \mathbf{B}, \mathbf{C}, \mathbf{D}) ; \quad \%$ respuesta en frecuencia de los valores singulares; o

$\gg \mathbf{w}=\operatorname{logspace}(\mathbf{l}, \mathbf{h}, \mathbf{n}) ; \mathbf{s v}=\operatorname{sigma}(\mathbf{A}, \mathbf{B}, \mathbf{C}, \mathbf{D}, \mathbf{w}) ; \% n$ puntos entre $10^{l}$ y $10^{h}$.

Comandos relacionados de MATLAB: semilogx, semilogy, bode, freqs, nichols, frsp, vsvd, vplot, pkvnorm

\section{Cálculo de normas $\mathcal{L}_{2}$ y $\mathcal{H}_{2}$}
Sea $G(s) \in \mathcal{L}_{2}$ y recordemos que la norma $\mathcal{L}_{2}$ de $G$ se define como

$$
\begin{aligned}
\|G\|_{2} & :=\sqrt{\frac{1}{2 \pi} \int_{-\infty}^{\infty} \operatorname{trace}\left\{G^{*}(j \omega) G(j \omega)\right\} d \omega} \\
& =\|g\|_{2} \\
& =\sqrt{\int_{-\infty}^{\infty} \operatorname{trace}\left\{g^{*}(t) g(t)\right\} d t}
\end{aligned}
$$

donde $g(t)$ denota el núcleo de convolución de $G$.

Es fácil ver que la norma $\mathcal{L}_{2}$ definida anteriormente es finita si y solo si la matriz de transferencia $G$ es estrictamente propia; es decir, $G(\infty)=0$. Por lo tanto, en general asumiremos que la matriz de transferencia es estrictamente propia cada vez que nos refiramos a la norma $\mathcal{L}_{2}$ de $G$ (por supuesto, esto también aplica a normas $\mathcal{H}_{2}$). Una forma directa de calcular la norma $\mathcal{L}_{2}$ es usar integral de contorno. Suponga que $G$ es estrictamente propia; entonces

$$
\begin{aligned}
\|G\|_{2}^{2} & =\frac{1}{2 \pi} \int_{-\infty}^{\infty} \operatorname{trace}\left\{G^{*}(j \omega) G(j \omega)\right\} d \omega \\
& =\frac{1}{2 \pi j} \oint \operatorname{trace}\left\{G^{\sim}(s) G(s)\right\} d s
\end{aligned}
$$

La última integral es una integral de contorno a lo largo del eje imaginario y alrededor de un semicírculo infinito en el semiplano izquierdo; la contribución de este semicírculo es cero porque $G$ es estrictamente propia. Por el teorema del residuo, $\|G\|_{2}^{2}$ es igual a la suma de los residuos de $\operatorname{trace}\left\{G^{\sim}(s) G(s)\right\}$ en sus polos del semiplano izquierdo.

Aunque $\|G\|_{2}$ puede, en principio, calcularse por su definición o por el método recién sugerido, en muchas aplicaciones es útil disponer de caracterizaciones alternativas y aprovechar representaciones en espacio de estados de $G$. El cálculo de la norma de una matriz de transferencia en $\mathcal{R H}_{2}$ es particularmente simple.

Lema 4.4 Considere una matriz de transferencia

$$
G(s)=\left[\begin{array}{c|c}
A & B \\
\hline C & 0
\end{array}\right]
$$

con $A$ estable. Entonces


\begin{equation*}
\|G\|_{2}^{2}=\operatorname{trace}\left(B^{*} Q B\right)=\operatorname{trace}\left(C P C^{*}\right) \tag{4.1}
\end{equation*}


donde $Q$ y $P$ son los Gramianos de observabilidad y controlabilidad, que pueden obtenerse de las siguientes ecuaciones de Lyapunov:

$$
A P+P A^{*}+B B^{*}=0 \quad A^{*} Q+Q A+C^{*} C=0
$$

Prueba. Como $G$ es estable, tenemos

$$
g(t)=\mathcal{L}^{-1}(G)= \begin{cases}C e^{A t} B, & t \geq 0 \\
0, & t<0\end{cases}
$$

y

$$
\begin{aligned}
\|G\|_{2}^{2} & =\int_{0}^{\infty} \operatorname{trace}\left\{g^{*}(t) g(t)\right\} d t=\int_{0}^{\infty} \operatorname{trace}\left\{g(t) g(t)^{*}\right\} d t \\
& =\int_{0}^{\infty} \operatorname{trace}\left\{B^{*} e^{A^{*} t} C^{*} C e^{A t} B\right\} d t=\int_{0}^{\infty} \operatorname{trace}\left\{C e^{A t} B B^{*} e^{A^{*} t} C^{*}\right\} d t
\end{aligned}
$$

El lema se sigue del hecho de que el Gramiano de controlabilidad de $(A, B)$ y el Gramiano de observabilidad de $(C, A)$ pueden representarse como

$$
Q=\int_{0}^{\infty} e^{A^{*} t} C^{*} C e^{A t} d t, \quad P=\int_{0}^{\infty} e^{A t} B B^{*} e^{A^{*} t} d t
$$

que también pueden obtenerse de

$$
A P+P A^{*}+B B^{*}=0 \quad A^{*} Q+Q A+C^{*} C=0
$$

Para calcular la norma $\mathcal{L}_{2}$ de una función de transferencia racional, $G(s) \in \mathcal{R} \mathcal{L}_{2}$, usando el enfoque de espacio de estados, escriba $G(s)=[G(s)]_{+}+[G(s)]_{-}$ con $G_{+} \in \mathcal{R} \mathcal{H}_{2}$ y $G_{-} \in \mathcal{R} \mathcal{H}_{2}^{\perp}$; entonces

$$
\|G\|_{2}^{2}=\left\|[G(s)]_{+}\right\|_{2}^{2}+\left\|[G(s)]_{-}\right\|_{2}^{2}
$$

donde $\left\|[G(s)]_{+}\right\|_{2}$ y $\left\|[G(s)]_{-}\right\|_{2}=\left\|[G(-s)]_{+}\right\|_{2}=\left\|\left([G(s)]_{-}\right)^{\sim}\right\|_{2}$ pueden calcularse con el lema anterior.

Otra caracterización útil de la norma $\mathcal{H}_{2}$ de $G$ es en términos de experimentos hipotéticos de entrada-salida. Sea $e_{i}$ el $i$-ésimo vector de la base estándar de $\mathbb{R}^{m}$, donde $m$ es la dimensión de entrada del sistema. Aplique la entrada impulsiva $\delta(t) e_{i}$ [$\delta(t)$ es el impulso unitario] y denote la salida por $z_{i}(t)$ ($=g(t) e_{i}$). Suponga $D=0$; entonces $z_{i} \in \mathcal{L}_{2+}$ y

$$
\|G\|_{2}^{2}=\sum_{i=1}^{m}\left\|z_{i}\right\|_{2}^{2}
$$

Nótese que esta caracterización de la norma $\mathcal{H}_{2}$ puede generalizarse adecuadamente para sistemas no lineales variante en el tiempo; véase Chen y Francis [1992] para una aplicación de esta norma en control muestreado.

Ejemplo 4.1 Considere una matriz de transferencia

$$
G=\left[\begin{array}{ll}
\frac{3(s+3)}{(s-1)(s+2)} & \frac{2}{s-1} \\
\frac{s+1}{(s+2)(s+3)} & \frac{1}{s-4}
\end{array}\right]=G_{s}+G_{u}
$$

con

$$
G_{s}=\left[\begin{array}{cc|cc}
-2 & 0 & -1 & 0 \\
0 & -3 & 2 & 0 \\
\hline 1 & 0 & 0 & 0 \\
1 & 1 & 0 & 0
\end{array}\right], \quad G_{u}=\left[\begin{array}{cc|cc}
1 & 0 & 4 & 2 \\
0 & 4 & 0 & 1 \\
\hline 1 & 0 & 0 & 0 \\
0 & 1 & 0 & 0
\end{array}\right]
$$

Entonces el comando $\mathbf{h} 2 \operatorname{norm}\left(\mathbf{G}_{\mathbf{s}}\right)$ da $\left\|G_{s}\right\|_{2}=0.6055$, y $\mathbf{h} 2 \operatorname{norm}\left(\mathbf{cjt}\left(\mathbf{G}_{\mathbf{u}}\right)\right)$ da $\left\|G_{u}\right\|_{2}=3.182$. Por tanto,\\
$\|G\|_{2}=\sqrt{\left\|G_{s}\right\|_{2}^{2}+\left\|G_{u}\right\|_{2}^{2}}=3.2393$.

\textbf{Comandos ilustrativos de MATLAB:}
$\gg \mathbf{P}=\operatorname{gram}(\mathbf{A}, \mathbf{B}) ; \mathbf{Q}=\operatorname{gram}\left(\mathbf{A}', \mathbf{C}'\right) ;$ o $\mathbf{P}=\operatorname{lyap}\left(\mathbf{A}, \mathbf{B} * \mathbf{B}'\right) ;$

$\gg\left[\mathbf{G}_{\mathbf{s}}, \mathbf{G}_{\mathbf{u}}\right]=\operatorname{sdecomp}(\mathbf{G}) ; \%$ descompone en partes estable y antiestable.

\section{Cálculo de normas $\mathcal{L}_{\infty}$ y $\mathcal{H}_{\infty}$}
Consideraremos primero, como en el caso de $\mathcal{L}_{2}$, cómo calcular la norma infinito de una matriz de transferencia en $\mathcal{R} \mathcal{L}_{\infty}$. Sea $G(s) \in \mathcal{R} \mathcal{L}_{\infty}$ y recordemos que la norma $\mathcal{L}_{\infty}$ de una función de transferencia matricial racional $G$ se define como

$$
\|G\|_{\infty}:=\sup _{\omega} \bar{\sigma}\{G(j \omega)\}
$$

El cálculo de la norma $\mathcal{L}_{\infty}$ de $G$ es complicado y requiere una búsqueda. Una interpretación en ingeniería de control de la norma infinito de una función de transferencia escalar $G$ es la distancia en el plano complejo desde el origen hasta el punto más lejano de la gráfica de Nyquist de $G$; también aparece como el valor pico en el diagrama de magnitud de Bode de $|G(j \omega)|$. Por lo tanto, la norma infinito de una función de transferencia puede, en principio, obtenerse gráficamente.

Para obtener una estimación, defina una malla fina de frecuencias:

$$
\left\{\omega_{1}, \cdots, \omega_{N}\right\}
$$

Entonces una estimación de $\|G\|_{\infty}$ es

$$
\max _{1 \leq k \leq N} \bar{\sigma}\left\{G\left(j \omega_{k}\right)\right\}
$$

Este valor suele leerse directamente de un diagrama de valores singulares de Bode. La norma $\mathcal{R} \mathcal{L}_{\infty}$ también puede calcularse en espacio de estados.

Lema 4.5 Sea $\gamma>0$ y

\[
G(s)=\left[\begin{array}{c|c}
A & B  \tag{4.2}\\
\hline C & D
\end{array}\right] \in \mathcal{R} \mathcal{L}_{\infty}
\]

Entonces $\|G\|_{\infty}<\gamma$ si y solo si $\bar{\sigma}(D)<\gamma$ y la matriz Hamiltoniana $H$ no tiene valores propios sobre el eje imaginario, donde

\[
H:=\left[\begin{array}{cc}
A+B R^{-1} D^{*} C & B R^{-1} B^{*}  \tag{4.3}\\
-C^{*}\left(I+D R^{-1} D^{*}\right) C & -\left(A+B R^{-1} D^{*} C\right)^{*}
\end{array}\right]
\]

y $R=\gamma^{2} I-D^{*} D$.

Prueba. Sea $\Phi(s)=\gamma^{2} I-G^{\sim}(s) G(s)$. Entonces es claro que $\|G\|_{\infty}<\gamma$ si y solo si $\Phi(j \omega)>0$ para todo $\omega \in \mathbb{R}$. Como $\Phi(\infty)=R>0$ y $\Phi(j \omega)$ es continua en $\omega$, se tiene que $\Phi(j \omega)>0$ para todo $\omega \in \mathbb{R}$ si y solo si $\Phi(j \omega)$ es no singular para todo $\omega \in \mathbb{R} \cup\{\infty\}$; es decir, $\Phi(s)$ no tiene ceros sobre el eje imaginario. Equivalentemente, $\Phi^{-1}(s)$ no tiene polos sobre el eje imaginario. Mediante álgebra simple se obtiene

% \insertimage[\label{img:calculo-hinfinito-realizacion-phi-inversa}]{2024_06_29_ce7c73b22613114bd4d6g-070_166_884_1004_615}{width=0.35\textwidth}{Realización empleada para el análisis de polos de $\Phi^{-1}(s)$ en la prueba del Lema 4.5.}

$$
\Phi^{-1}(s)=\left[\begin{array}{c|c}
H & {\left[\begin{array}{c}
B R^{-1} \\
-C^* D R^{-1}
\end{array}\right]} \\
\hline\left[\begin{array}{ll}
R^{-1} D^* C & R^{-1} B^*
\end{array}\right] & R^{-1}
\end{array}\right] .
$$

Por tanto, la conclusión se sigue si la realización anterior no tiene modos no controlables ni no observables sobre el eje imaginario. Suponga que $j \omega_{0}$ es un valor propio de $H$, pero no un polo de $\Phi^{-1}(s)$. Entonces $j \omega_{0}$ debe ser un modo no observable de $\left(\left[\begin{array}{ll}R^{-1} D^{*} C & R^{-1} B^{*}\end{array}\right], H\right)$ o un modo no controlable de $\left(H,\left[\begin{array}{c}B R^{-1} \\ -C^{*} D R^{-1}\end{array}\right]\right)$. Suponga ahora que $j \omega_{0}$ es un modo no observable de $\left(\left[\begin{array}{ll}R^{-1} D^{*} C & R^{-1} B^{*}\end{array}\right], H\right)$. Entonces existe un $x_{0}=\left[\begin{array}{l}x_{1} \\ x_{2}\end{array}\right] \neq 0$ tal que

$$
H x_{0}=j \omega_{0} x_{0},\left[\begin{array}{ll}
R^{-1} D^{*} C & R^{-1} B^{*}
\end{array}\right] x_{0}=0
$$

Estas ecuaciones se simplifican a

$$
\begin{aligned}
\left(j \omega_{0} I-A\right) x_{1} & =0 \\
\left(j \omega_{0} I+A^{*}\right) x_{2} & =-C^{*} C x_{1} \\
D^{*} C x_{1}+B^{*} x_{2} & =0 .
\end{aligned}
$$

Como $A$ no tiene valores propios sobre el eje imaginario, tenemos $x_{1}=0$ y $x_{2}=0$. Esto contradice nuestra suposición, y por tanto la realización no tiene modos no observables sobre el eje imaginario.

De forma similar, también se llega a contradicción si se supone que $j \omega_{0}$ es un modo no controlable de $\left(H,\left[\begin{array}{c}B R^{-1} \\ -C^{*} D R^{-1}\end{array}\right]\right)$.

\subsection*{Algoritmo de bisección}
El Lema 4.5 sugiere el siguiente algoritmo de bisección para calcular la norma en $\mathcal{R} \mathcal{L}_{\infty}$:

(a) Seleccione una cota superior $\gamma_{u}$ y una cota inferior $\gamma_{l}$ tales que $\gamma_{l} \leq\|G\|_{\infty} \leq \gamma_{u}$;

(b) Si $\left(\gamma_{u}-\gamma_{l}\right) / \gamma_{l} \leq$ el nivel especificado, detenga; $\|G\| \approx\left(\gamma_{u}+\gamma_{l}\right) / 2$. De lo contrario, vaya al siguiente paso;

(c) Fije $\gamma=\left(\gamma_{l}+\gamma_{u}\right) / 2$;

(d) Pruebe si $\|G\|_{\infty}<\gamma$ calculando los valores propios de $H$ para el $\gamma$ dado;

(e) Si $H$ tiene un valor propio en $j \mathbb{R}$, fije $\gamma_{l}=\gamma$; en caso contrario, fije $\gamma_{u}=\gamma$; vuelva al paso (b).

Por supuesto, el algoritmo anterior aplica también al cálculo de la norma $\mathcal{H}_{\infty}$. Por tanto, el cálculo de la norma $\mathcal{L}_{\infty}$ requiere una búsqueda, ya sea sobre $\gamma$ o sobre $\omega$, en contraste con el cálculo de la norma en $\mathcal{L}_{2}$ ($\mathcal{H}_{2}$), que no la requiere. Una situación algo análoga ocurre para matrices constantes con las normas $\|M\|_{2}^{2}=\operatorname{trace}\left(M^{*} M\right)$ y $\|M\|_{\infty}=\bar{\sigma}[M]$. En principio, $\|M\|_{2}^{2}$ puede calcularse exactamente con un número finito de operaciones, al igual que la prueba de si $\bar{\sigma}(M)<\gamma$ (por ejemplo, $\gamma^{2} I-M^{*} M>0$), pero el valor de $\bar{\sigma}(M)$ no. Para calcular $\bar{\sigma}(M)$, debemos usar algún tipo de algoritmo iterativo.

Observación 4.2 Es claro que $\|G\|_{\infty}<\gamma$ si y solo si $\left\|\gamma^{-1} G\right\|_{\infty}<1$. Por ello, no hay pérdida de generalidad al suponer $\gamma=1$. Esta suposición se hará con frecuencia en el resto del libro. También se observa que existen otros algoritmos rápidos para realizar el cálculo de norma anterior; no obstante, este algoritmo de bisección es el más simple.

Pueden darse interpretaciones adicionales para la norma $\mathcal{H}_{\infty}$ de una matriz de transferencia estable. Cuando $G(s)$ es un sistema SISO, la norma $\mathcal{H}_{\infty}$ de $G(s)$ puede verse como el mayor factor posible de amplificación de la respuesta en régimen permanente del sistema ante excitaciones sinusoidales. Por ejemplo, la respuesta en régimen permanente del sistema ante una entrada sinusoidal $u(t)=U \sin \left(\omega_{0} t+\phi\right)$ es

$$
y(t)=U\left|G\left(j \omega_{0}\right)\right| \sin \left(\omega_{0} t+\phi+\angle G\left(j \omega_{0}\right)\right)
$$

y así, el mayor factor posible de amplificación es $\sup \left|G\left(j \omega_{0}\right)\right|$, que es precisamente la norma $\mathcal{H}_{\infty}$ de la función de transferencia.

En el caso MIMO, la norma $\mathcal{H}_{\infty}$ de una matriz de transferencia $G \in \mathcal{R} \mathcal{H}_{\infty}$ también puede interpretarse como el mayor factor posible de amplificación de la respuesta en régimen permanente ante excitaciones sinusoidales en el siguiente sentido: sean las entradas sinusoidales

$$
u(t)=\left[\begin{array}{c}
u_{1} \sin \left(\omega_{0} t+\phi_{1}\right) \\
u_{2} \sin \left(\omega_{0} t+\phi_{2}\right) \\
\vdots \\
u_{q} \sin \left(\omega_{0} t+\phi_{q}\right)
\end{array}\right], \quad \hat{u}=\left[\begin{array}{c}
u_{1} \\
u_{2} \\
\vdots \\
u_{q}
\end{array}\right]
$$

Entonces la respuesta en régimen permanente del sistema puede escribirse como

$$
y(t)=\left[\begin{array}{c}
y_{1} \sin \left(\omega_{0} t+\theta_{1}\right) \\
y_{2} \sin \left(\omega_{0} t+\theta_{2}\right) \\
\vdots \\
y_{p} \sin \left(\omega_{0} t+\theta_{p}\right)
\end{array}\right], \quad \hat{y}=\left[\begin{array}{c}
y_{1} \\
y_{2} \\
\vdots \\
y_{p}
\end{array}\right]
$$

para algunos $y_{i}, \theta_{i}, i=1,2, \ldots, p$, y además

$$
\|G\|_{\infty}=\sup _{\phi_{i}, \omega_{o}, \hat{u}} \frac{\|\hat{y}\|}{\|\hat{u}\|}
$$

donde $\|\cdot\|$ es la norma euclidiana. Los detalles se dejan como ejercicio.

Ejemplo 4.2 Considere un sistema masa-resorte-amortiguador como se muestra en la Figura 4.2.

\insertimage[\label{img:dos-masas-resortes-amortiguadores}]{2024_06_29_ce7c73b22613114bd4d6g-072_626_650_1232_735}{width=0.35\textwidth}{Un sistema de dos masas-resortes-amortiguadores.}

El sistema dinámico puede describirse mediante las siguientes ecuaciones diferenciales:

$$
\left[\begin{array}{c}
\dot{x}_{1} \\
\dot{x}_{2} \\
\dot{x}_{3} \\
\dot{x}_{4}
\end{array}\right]=A\left[\begin{array}{l}
x_{1} \\
x_{2} \\
x_{3} \\
x_{4}
\end{array}\right]+B\left[\begin{array}{l}
F_{1} \\
F_{2}
\end{array}\right]
$$

\insertimage[\label{img:hinfinito-pico-valor-singular}]{2024_06_29_ce7c73b22613114bd4d6g-073_656_1016_412_541}{width=0.65\textwidth}{$\|G\|_{\infty}$ es el pico del mayor valor singular de $G(j \omega)$.}

con

$$
A=\left[\begin{array}{cccc}
0 & 0 & 1 & 0 \\
0 & 0 & 0 & 1 \\
-\frac{k_{1}}{m_{1}} & \frac{k_{1}}{m_{1}} & -\frac{b_{1}}{m_{1}} & \frac{b_{1}}{m_{1}} \\
\frac{k_{1}}{m_{2}} & -\frac{k_{1}+k_{2}}{m_{2}} & \frac{b_{1}}{m_{2}} & -\frac{b_{1}+b_{2}}{m_{2}}
\end{array}\right], \quad B=\left[\begin{array}{cc}
0 & 0 \\
0 & 0 \\
\frac{1}{m_{1}} & 0 \\
0 & \frac{1}{m_{2}}
\end{array}\right]
$$

Suponga que $G(s)$ es la matriz de transferencia de $\left(F_{1}, F_{2}\right)$ a $\left(x_{1}, x_{2}\right)$; es decir,

$$
C=\left[\begin{array}{llll}
1 & 0 & 0 & 0 \\
0 & 1 & 0 & 0
\end{array}\right], \quad D=0
$$

y suponga $k_{1}=1, k_{2}=4, b_{1}=0.2, b_{2}=0.1, m_{1}=1$ y $m_{2}=2$, con unidades apropiadas. Los siguientes comandos de MATLAB generan el diagrama de Bode de valores singulares del sistema anterior, como se muestra en la Figura 4.3.

$\gg \mathrm{G}=\operatorname{pck}(\mathrm{A}, \mathrm{B}, \mathrm{C}, \mathrm{D})$;

$\gg \operatorname{hinfnorm}(\mathrm{G}, \mathbf{0.0001})$ o $\operatorname{linfnorm}(\mathrm{G}, \mathbf{0.0001}) \%$ error relativo $\leq 0.0001$

$\gg \mathbf{w}=\operatorname{logspace}(-\mathbf{1}, \mathbf{1}, \mathbf{200}) ; \% \mathbf{200}$ puntos entre $1=10^{-\mathbf{1}}$ y $10=10^{\mathbf{1}}$;

$\gg \mathbf{Gf}=\mathbf{frsp}(\mathbf{G}, \mathbf{w}) ; \%$ cálculo de respuesta en frecuencia;

$\gg[\mathbf{u}, \mathbf{s}, \mathbf{v}]=\operatorname{vsvd}(\mathbf{Gf}) ; \%$ SVD en cada frecuencia;

$\gg$ vplot('liv, lm', \textbackslash mathbf\{s\}), grid \% grafica ambos valores singulares y la cuadrícula.

Entonces la norma $\mathcal{H}_{\infty}$ de esta matriz de transferencia es $\|G(s)\|_{\infty}=11.47$, lo cual se muestra como el pico del diagrama de Bode del mayor valor singular en la Figura 4.3. Como el pico se alcanza en $\omega_{\max }=0.8483$, excitar el sistema con la siguiente entrada sinusoidal

$$
\left[\begin{array}{l}
F_{1} \\
F_{2}
\end{array}\right]=\left[\begin{array}{l}
0.9614 \sin (0.8483 t) \\
0.2753 \sin (0.8483 t-0.12)
\end{array}\right]
$$

da la respuesta en régimen permanente del sistema como

$$
\left[\begin{array}{l}
x_{1} \\
x_{2}
\end{array}\right]=\left[\begin{array}{l}
11.47 \times 0.9614 \sin (0.8483 t-1.5483) \\
11.47 \times 0.2753 \sin (0.8483 t-1.4283)
\end{array}\right]
$$

Esto muestra que la respuesta del sistema se amplifica 11.47 veces para una señal de entrada en la frecuencia $\omega_{\max }$, lo cual podría ser indeseable si $F_{1}$ y $F_{2}$ son fuerzas de perturbación y $x_{1}$ y $x_{2}$ son posiciones que deben mantenerse estables.

Ejemplo 4.3 Considere una matriz de transferencia de dos por dos

$$
G(s)=\left[\begin{array}{cc}
\frac{10(s+1)}{s^{2}+0.2 s+100} & \frac{1}{s+1} \\
\frac{s+2}{s^{2}+0.1 s+10} & \frac{5(s+1)}{(s+2)(s+3)}
\end{array}\right]
$$

Una realización en espacio de estados de $G$ puede obtenerse usando los siguientes comandos de MATLAB:

$$
\begin{aligned}
& \gg \mathrm{G}11=\operatorname{nd}2\operatorname{sys}([10,10],[1,0.2,100]) \\
& \gg \mathrm{G}12=\operatorname{nd}2\operatorname{sys}(1,[1,1]) \\
& \gg \mathrm{G}21=\operatorname{nd}2\operatorname{sys}([1,2],[1,0.1,10]) \\
& \gg \mathrm{G}22=\operatorname{nd}2\operatorname{sys}([5,5],[1,5,6]) \\
& \gg \mathrm{G}=\operatorname{sbs}(\operatorname{abv}(\mathrm{G}11, \mathrm{G}21), \operatorname{abv}(\mathrm{G}12, \mathrm{G}22))
\end{aligned}
$$

A continuación, definimos una malla de frecuencias para calcular la respuesta en frecuencia de $G$ y los valores singulares de $G(j \omega)$ sobre un rango de frecuencias adecuado.

$\gg \mathbf{w}=\operatorname{logspace}(\mathbf{0},\mathbf{2},\mathbf{200}) ; \% \mathbf{200}$ puntos entre $1=10^{\mathbf{0}}$ y $100=10^{\mathbf{2}}$;

$\gg \mathbf{Gf}=\mathbf{frsp}(\mathbf{G}, \mathbf{w}) ; \%$ cálculo de respuesta en frecuencia;

$\gg[\mathbf{u}, \mathbf{s}, \mathbf{v}]=\mathbf{vsvd}(\mathbf{Gf}) ; \%$ SVD en cada frecuencia;\\
$\gg$ vplot(' \$\textbackslash mathbf\{liv\}, \textbackslash mathbf\{lm\}', \textbackslash mathbf\{s\}), grid \% grafica ambos valores singulares y la cuadrícula;

$\gg$ pkvnorm(s) \% obtiene la norma a partir de la respuesta en frecuencia de los valores singulares.

Los valores singulares de $G(j \omega)$ se grafican en la Figura 4.4, lo cual da una estimación de $\|G\|_{\infty} \approx 32.861$. El algoritmo de bisección en espacio de estados descrito previamente conduce a $\|G\|_{\infty}=50.25 \pm 0.01$, y el comando correspondiente de MATLAB es

$\gg \operatorname{hinfnorm}(\mathbf{G}, \mathbf{0.0001})$ o $\operatorname{linfnorm}(\mathbf{G}, \mathbf{0.0001}) \%$ error relativo $\leq 0.0001$.

\insertimage[\label{img:mayor-menor-valor-singular-gjw}]{2024_06_29_ce7c73b22613114bd4d6g-075_653_1014_828_542}{width=0.65\textwidth}{El mayor y el menor valor singular de $G(j \omega)$.}

Los resultados computacionales anteriores muestran claramente que el método gráfico puede llevar a una respuesta incorrecta para un sistema ligeramente amortiguado si la malla de frecuencias no es suficientemente densa. En efecto, obtendríamos $\|G\|_{\infty} \approx 43.525, 48.286$ y 49.737 con el método gráfico si se usan 400, 800 y 1600 puntos de frecuencia, respectivamente.

Comandos relacionados de MATLAB: linfnorm, vnorm, getiv, scliv, var2con, xtract, xtracti

\section{Notas y referencias}
El concepto básico de espacios de funciones presentado en este capítulo puede encontrarse en cualquier texto estándar de análisis funcional, por ejemplo, Naylor y Sell [1982] y Gohberg y Goldberg [1981]. Las interpretaciones en teoría de sistemas de normas y espacios de funciones pueden encontrarse en Desoer y Vidyasagar [1975]. El algoritmo de bisección para el cálculo de la norma en $\mathcal{L}_{\infty}$ fue desarrollado por primera vez en Boyd, Balakrishnan y Kabamba [1989]. Un algoritmo más eficiente para calcular la norma en $\mathcal{L}_{\infty}$ se presenta en Bruinsma y Steinbuch [1990].

\section{Problemas}
Problema 4.1 Sea $G(s)$ una matriz en $\mathcal{R} \mathcal{H}_{\infty}$. Demuestre que

$$
\left\|\left[\begin{array}{c}
G \\
I
\end{array}\right]\right\|_{\infty}^{2}=\|G\|_{\infty}^{2}+1
$$

Problema 4.2 (relación de Parseval) Sean $f(t), g(t) \in \mathcal{L}_{2}, F(j \omega)=\mathcal{F}\{f(t)\}$, y $G(j \omega)=\mathcal{F}\{g(t)\}$. Demuestre que

$$
\int_{-\infty}^{\infty} f(t) g(t) d t=\frac{1}{2 \pi} \int_{-\infty}^{\infty} F(j \omega) G^{*}(j \omega) d \omega
$$

y

$$
\int_{-\infty}^{\infty}|f(t)|^{2} d t=\frac{1}{2 \pi} \int_{-\infty}^{\infty}|F(j \omega)|^{2} d \omega
$$

Nótese que

$$
F(j \omega)=\int_{-\infty}^{\infty} f(t) e^{-j \omega t} d t, \quad f(t)=\mathcal{F}^{-1}(F(j \omega))=\frac{1}{2 \pi} \int_{-\infty}^{\infty} F(j \omega) e^{j \omega t} d \omega
$$

donde $\mathcal{F}^{-1}$ denota la transformada inversa de Fourier.

Problema 4.3 Suponga que $A$ es estable. Demuestre

$$
\int_{-\infty}^{\infty}(j \omega I-A)^{-1} d \omega=\pi I
$$

Suponga $G(s)=\left[\begin{array}{c|c}A & B \\ \hline C & 0\end{array}\right] \in \mathcal{R} \mathcal{H}_{\infty}$ y sea $Q=Q^{*}$ el Gramiano de observabilidad. Use la fórmula anterior para demostrar que

$$
\frac{1}{2 \pi} \int_{-\infty}^{\infty} G^{\sim}(j \omega) G(j \omega) d \omega=B^{*} Q B
$$

[Pista: Use el hecho de que $G^{\sim}(s) G(s)=F^{\sim}(s)+F(s)$ y $F(s)=B^{*} Q(s I-A)^{-1} B$.]

Problema 4.4 Calcule la norma 2 y la norma infinito de los siguientes sistemas:

$$
G_{1}(s)=\left[\begin{array}{cc}
\frac{1}{s+1} & \frac{s+3}{(s+1)(s-2)} \\
\frac{10}{s-2} & \frac{5}{s+3}
\end{array}\right], \quad G_{2}(s)=\left[\begin{array}{ll|l}
1 & 0 & 1 \\
2 & 3 & 1 \\
\hline 1 & 2 & 0
\end{array}\right]
$$

$$
G_{3}(s)=\left[\begin{array}{cc|c}
-1 & -2 & 1 \\
1 & 0 & 0 \\
\hline 2 & 3 & 0
\end{array}\right], \quad G_{4}(s)=\left[\begin{array}{ccc|cc}
-1 & -2 & -3 & 1 & 2 \\
1 & 0 & 0 & 0 & 1 \\
0 & 1 & 0 & 2 & 0 \\
\hline 1 & 0 & 0 & 1 & 0 \\
0 & 1 & 1 & 0 & 2
\end{array}\right]
$$

Problema 4.5 Sea $r(t)=\sin \omega t$ la señal de entrada de una planta

$$
G(s)=\frac{\omega_{n}^{2}}{s^{2}+2 \xi \omega_{n} s+\omega_{n}^{2}}
$$

con $0<\xi<1 / \sqrt{2}$. Encuentre la respuesta en régimen permanente del sistema $y(t)$. Encuentre también la frecuencia $\omega$ que produce la mayor magnitud de la respuesta en régimen permanente de $y(t)$.

Problema 4.6 Sea $G(s) \in \mathcal{R} \mathcal{H}_{\infty}$ una matriz de transferencia $p \times q$ y $y=G(s) u$. Suponga

$$
u(t)=\left[\begin{array}{c}
u_{1} \sin \left(\omega_{0} t+\phi_{1}\right) \\
u_{2} \sin \left(\omega_{0} t+\phi_{2}\right) \\
\vdots \\
u_{q} \sin \left(\omega_{0} t+\phi_{q}\right)
\end{array}\right], \quad \hat{u}=\left[\begin{array}{c}
u_{1} \\
u_{2} \\
\vdots \\
u_{q}
\end{array}\right]
$$

Demuestre que la respuesta en régimen permanente del sistema está dada por

$$
y(t)=\left[\begin{array}{c}
y_{1} \sin \left(\omega_{0} t+\theta_{1}\right) \\
y_{2} \sin \left(\omega_{0} t+\theta_{2}\right) \\
\vdots \\
y_{p} \sin \left(\omega_{0} t+\theta_{p}\right)
\end{array}\right], \quad \hat{y}=\left[\begin{array}{c}
y_{1} \\
y_{2} \\
\vdots \\
y_{p}
\end{array}\right]
$$

para ciertos $y_{i}$ y $\theta_{i}, i=1,2, \ldots, p$. Demuestre que

$\sup _{\phi_{i}, \omega_{o},\|\hat{u}\|_{2} \leq 1}\|\hat{y}\|_{2}=\|G\|_{\infty}$.

Problema 4.7 Escriba un programa en MATLAB para graficar, en función de $\gamma$, la distancia desde el eje imaginario hasta el valor propio más cercano de la matriz Hamiltoniana para un modelo en espacio de estados dado con $A$ estable. Pruébelo en

$$
\left[\begin{array}{cc}
\frac{s+1}{(s+2)(s+3)} & \frac{s}{s+1} \\
\frac{s^{2}-2}{(s+3)(s+4)} & \frac{s+4}{(s+1)(s+2)}
\end{array}\right]
$$

Lea el valor de la norma $\mathcal{H}_{\infty}$. Compárelo con la función \texttt{hinfnorm} de MATLAB.

Problema 4.8 Sea $G(s)=\frac{1}{\left(s^{2}+2 \xi s+1\right)(s+1)}$. Calcule $\|G(s)\|_{\infty}$ usando el diagrama de Bode y el algoritmo en espacio de estados, respectivamente, para $\xi=1,0.1,0.01,0.001$, y compare los resultados.

\chapter{Estabilidad interna}
Este capítulo introduce la estructura de realimentación y discute su estabilidad, junto con diversas pruebas de estabilidad. La organización de este capítulo es la siguiente: la Sección 5.1 discute la necesidad de introducir una estructura de realimentación y describe la configuración general de realimentación. La Sección 5.2 define el buen planteamiento del lazo de realimentación. La Sección 5.3 introduce la noción de estabilidad interna y varias pruebas de estabilidad. La Sección 5.4 introduce las factorizaciones coprimas estables de matrices racionales. En esta sección también se consideran las condiciones de estabilidad en términos de distintas factorizaciones coprimas.

\section{Estructura de retroalimentación}
Al diseñar sistemas de control, existen varios temas fundamentales que trascienden los límites de aplicaciones específicas. Aunque pueden diferir para cada aplicación y tener distintos niveles de importancia, estos temas son generales en su relación con los objetivos y procedimientos del diseño de control. Un punto central en estos temas es el requisito de proporcionar un desempeño satisfactorio frente a errores de modelado, variaciones del sistema e incertidumbre. De hecho, este requisito fue la motivación original para el desarrollo de sistemas de realimentación. La realimentación solo se requiere cuando el desempeño del sistema no puede lograrse debido a la incertidumbre en las características del sistema. Un tratamiento más detallado de las incertidumbres de modelo y sus representaciones se discutirá en el Capítulo 8.

Por el momento, suponiendo que se nos da un modelo que incluye una representación de incertidumbre que creemos captura adecuadamente las características esenciales de la planta, el siguiente paso en el proceso de diseño del controlador es determinar qué estructura es necesaria para lograr el desempeño deseado. El prefiltrado de señales de entrada (o control en lazo abierto) puede cambiar la respuesta dinámica del conjunto de modelos, pero no puede reducir el efecto de la incertidumbre. Si la incertidumbre es demasiado grande para lograr la precisión de respuesta deseada, entonces se requiere una estructura de realimentación. Sin embargo, el mero hecho de asumir una estructura de realimentación no garantiza una reducción de la incertidumbre, y existen muchos obstáculos para lograr los beneficios de reducción de incertidumbre que aporta la realimentación. En particular, dado que para cualquier conjunto de modelos razonable que represente la incertidumbre de un sistema físico esta incertidumbre se vuelve grande y la fase es completamente desconocida a frecuencias suficientemente altas, la ganancia de lazo debe ser pequeña en esas frecuencias para evitar desestabilizar la dinámica de alta frecuencia del sistema. Aún peor, el sistema con realimentación en realidad incrementa la incertidumbre y la sensibilidad en los rangos de frecuencia donde la incertidumbre es significativamente grande. En otras palabras, debido al tipo de conjuntos requeridos para modelar razonablemente sistemas físicos y a la restricción de que nuestros controladores sean causales, no podemos usar realimentación (ni ninguna otra estructura de control) para hacer que nuestro conjunto de modelos en lazo cerrado sea un subconjunto propio del conjunto de modelos en lazo abierto. A menudo, lo que sí puede lograrse con un uso inteligente de la realimentación es una reducción significativa de incertidumbre para ciertas señales importantes, con un pequeño incremento distribuido en otras señales. Por tanto, el problema de diseño por realimentación se centra en el compromiso involucrado en reducir el impacto global de la incertidumbre. Este compromiso también aparece, por ejemplo, cuando se usa realimentación para reducir el error de mando/perturbación minimizando al mismo tiempo la degradación de respuesta debida al ruido de medición. Para tener valor práctico, una técnica de diseño debe proporcionar medios para realizar estos compromisos. Discutiremos estos compromisos con mayor detalle en el capítulo siguiente.

Para enfocar la discusión, consideraremos la configuración estándar de realimentación mostrada en la Figura 5.1. Esta consiste en la interconexión de la planta $P$ y el controlador $K$, excitada por la referencia $r$, el ruido del sensor $n$, la perturbación de entrada de planta $d_{i}$ y la perturbación de salida de planta $d$. En general, todas las señales se suponen multivariables y todas las matrices de transferencia se suponen de dimensiones apropiadas.

\insertimage[\label{img:configuracion-estandar-retroalimentacion}]{2024_06_29_ce7c73b22613114bd4d6g-080_293_879_1363_604}{width=0.55\textwidth}{Configuración estándar de realimentación.}

\section{Bien posicionamiento del lazo de retroalimentación}
Suponga que la planta $P$ y el controlador $K$ de la Figura 5.1 son matrices de transferencia racionales reales propias y fijas. Entonces, la primera pregunta natural es si la interconexión en realimentación tiene sentido o es físicamente realizable. Para ser más específicos, considere el siguiente ejemplo simple donde

$$
P=-\frac{s-1}{s+2}, \quad K=1
$$

son ambas funciones de transferencia propias. Sin embargo,

$$
u=\frac{(s+2)}{3}(r-n-d)+\frac{s-1}{3} d_{i}
$$

Es decir, las funciones de transferencia desde las señales externas $r-n-d$ y $d_{i}$ hacia $u$ no son propias. Por tanto, el sistema de realimentación no es físicamente realizable.

Definición 5.1 Se dice que un sistema de realimentación está bien planteado si todas las matrices de transferencia en lazo cerrado están bien definidas y son propias.

Ahora suponga que todas las señales externas $r, n, d$ y $d_{i}$ están especificadas y que las matrices de transferencia en lazo cerrado desde ellas hasta $u$ están, respectivamente, bien definidas y son propias. Entonces $y$ y todas las demás señales también están bien definidas, y las matrices de transferencia relacionadas son propias. Además, dado que las matrices de transferencia desde $d$ y $n$ hacia $u$ son las mismas y difieren de la matriz de transferencia desde $r$ hacia $u$ solo por un signo, el sistema está bien planteado si y solo si la matriz de transferencia desde $\left[\begin{array}{c}d_{i} \\ d\end{array}\right]$ hacia $u$ existe y es propia.

Para ser consistentes con la notación usada en el resto del libro, denotaremos


\begin{equation*}
\hat{K}:=-K \tag{5.1}
\end{equation*}


y reagruparamos las señales de entrada externas en el lazo de realimentación como $w_{1}$ y $w_{2}$, y reagruparamos las señales de entrada de la planta y del controlador como $e_{1}$ y $e_{2}$. Entonces el lazo de realimentación con la planta y el controlador puede representarse de manera simple como en la Figura 5.2, y el sistema está bien planteado si y solo si la matriz de transferencia desde $\left[\begin{array}{l}w_{1} \\ w_{2}\end{array}\right]$ hacia $e_{1}$ existe y es propia.

\insertimage[\label{img:diagrama-analisis-estabilidad-interna}]{2024_06_29_ce7c73b22613114bd4d6g-081_262_680_1601_709}{width=0.45\textwidth}{Diagrama para análisis de estabilidad interna.}

Lema 5.1 El sistema de realimentación de la Figura 5.2 está bien planteado si y solo si


\begin{equation*}
I-\hat{K}(\infty) P(\infty) \tag{5.2}
\end{equation*}


es invertible.

Prueba. El sistema de la Figura 5.2 puede representarse en forma de ecuaciones como

$$
\begin{aligned}
& e_{1}=w_{1}+\hat{K} e_{2} \\
& e_{2}=w_{2}+P e_{1}
\end{aligned}
$$

Entonces se obtiene para $e_{1}$ la expresión

$$
(I-\hat{K} P) e_{1}=w_{1}+\hat{K} w_{2}
$$

Así, el buen planteamiento es equivalente a la condición de que $(I-\hat{K} P)^{-1}$ exista y sea propia. Pero esto es equivalente a que el término constante de la función de transferencia $I-\hat{K} P$ sea invertible.

Es directo mostrar que la ecuación (5.2) es equivalente a cualquiera de las dos condiciones siguientes:


\begin{gather*}
{\left[\begin{array}{cc}
I & -\hat{K}(\infty) \\
-P(\infty) & I
\end{array}\right] \text { es invertible; }}  \tag{5.3}\\
I-P(\infty) \hat{K}(\infty) \text { es invertible. }
\end{gather*}


La condición de buen planteamiento se expresa de forma simple en términos de realizaciones en espacio de estados. Introduzca realizaciones de $P$ y $\hat{K}$:

$$
P=\left[\begin{array}{l|l}
A & B \\
\hline C & D
\end{array}\right], \quad \hat{K}=\left[\begin{array}{c|c}
\hat{A} & \hat{B} \\
\hline \hat{C} & \hat{D}
\end{array}\right]
$$

Entonces $P(\infty)=D, \hat{K}(\infty)=\hat{D}$ y la condición de buen planteamiento en la ecuación (5.3) es equivalente a la invertibilidad de $\left[\begin{array}{cc}I & -\hat{D} \\ -D & I\end{array}\right]$. Afortunadamente, en la mayoría de casos prácticos se tiene $D=0$, y por tanto el buen planteamiento para la mayoría de sistemas de control prácticos está garantizado.

\section{Estabilidad interna}
Considere un sistema descrito por el diagrama de bloques estándar de la Figura 5.2 y suponga que está bien planteado.

Definición 5.2 Se dice que el sistema de la Figura 5.2 es internamente estable si la matriz de transferencia


\begin{align*}
{\left[\begin{array}{cc}
I & -\hat{K} \\
-P & I
\end{array}\right]^{-1} } & =\left[\begin{array}{cc}
(I-\hat{K} P)^{-1} & \hat{K}(I-P \hat{K})^{-1} \\
P(I-\hat{K} P)^{-1} & (I-P \hat{K})^{-1}
\end{array}\right]  \tag{5.4}\\
& =\left[\begin{array}{cc}
I+\hat{K}(I-P \hat{K})^{-1} P & \hat{K}(I-P \hat{K})^{-1} \\
(I-P \hat{K})^{-1} P & (I-P \hat{K})^{-1}
\end{array}\right]
\end{align*}


desde $\left(w_{1}, w_{2}\right)$ hacia $\left(e_{1}, e_{2}\right)$ pertenece a $\mathcal{R} \mathcal{H}_{\infty}$.

Nótese que, para verificar estabilidad interna, es necesario (y suficiente) comprobar si cada una de las cuatro matrices de transferencia de la ecuación (5.4) está en $\mathcal{H}_{\infty}$. No puede concluirse estabilidad aunque tres de las cuatro matrices de transferencia en (5.4) estén en $\mathcal{H}_{\infty}$. Por ejemplo, sea la función de transferencia del sistema interconectado

$$
P=\frac{s-1}{s+1}, \quad \hat{K}=-\frac{1}{s-1}
$$

Entonces es fácil calcular

$$
\left[\begin{array}{l}
e_{1} \\
e_{2}
\end{array}\right]=\left[\begin{array}{cc}
\frac{s+1}{s+2} & -\frac{s+1}{(s-1)(s+2)} \\
\frac{s-1}{s+2} & \frac{s+1}{s+2}
\end{array}\right]\left[\begin{array}{l}
w_{1} \\
w_{2}
\end{array}\right]
$$

lo cual muestra que el sistema no es internamente estable aunque tres de las cuatro funciones de transferencia sean estables.

Observación 5.1 La estabilidad interna es un requisito básico para un sistema práctico de realimentación. Esto se debe a que todos los sistemas interconectados pueden estar inevitablemente sujetos a condiciones iniciales no nulas y a ciertos errores (posiblemente pequeños), y en la práctica no puede tolerarse que esos errores en algunas ubicaciones produzcan señales no acotadas en otras ubicaciones del sistema en lazo cerrado. La estabilidad interna garantiza que todas las señales de un sistema estén acotadas siempre que las señales inyectadas (en cualquier ubicación) estén acotadas.

Sin embargo, existen algunos casos especiales bajo los cuales determinar la estabilidad del sistema es simple.

Corolario 5.2 Suponga que $\hat{K} \in \mathcal{R H}_{\infty}$. Entonces el sistema de la Figura 5.2 es internamente estable si y solo si está bien planteado y $P(I-\hat{K} P)^{-1} \in \mathcal{R} \mathcal{H}_{\infty}$.

Prueba. La necesidad es obvia. Para demostrar suficiencia, basta mostrar que $(I-P \hat{K})^{-1} \in \mathcal{R} \mathcal{H}_{\infty}$. Pero esto se sigue de

$$
(I-P \hat{K})^{-1}=I+(I-P \hat{K})^{-1} P \hat{K}
$$

y de que $(I-P \hat{K})^{-1} P, \hat{K} \in \mathcal{R} \mathcal{H}_{\infty}$.

También tenemos lo siguiente:

Corolario 5.3 Suponga que $P \in \mathcal{R} \mathcal{H}_{\infty}$. Entonces el sistema de la Figura 5.2 es internamente estable si y solo si está bien planteado y $\hat{K}(I-P \hat{K})^{-1} \in \mathcal{R} \mathcal{H}_{\infty}$.

Corolario 5.4 Suponga que $P \in \mathcal{R} \mathcal{H}_{\infty}$ y $\hat{K} \in \mathcal{R} \mathcal{H}_{\infty}$. Entonces el sistema de la Figura 5.2 es internamente estable si y solo si $(I-P \hat{K})^{-1} \in \mathcal{R} \mathcal{H}_{\infty}$ o, equivalentemente, si $\operatorname{det}(I-P(s) \hat{K}(s))$ no tiene ceros en el semiplano derecho cerrado.

Nótese que toda la discusión y conclusiones previas aplican por igual a plantas y controladores de dimensión infinita. Para estudiar el caso más general, limitaremos nuestra discusión a sistemas de dimensión finita y definimos

$$
\begin{aligned}
& n_{k}:=\text { número de polos en semiplano derecho abierto (rhp) de } \hat{K}(s) \\
& n_{p}:=\text { número de polos en semiplano derecho abierto (rhp) de } P(s) .
\end{aligned}
$$

Teorema 5.5 El sistema es internamente estable si y solo si está bien planteado y

(i) el número de polos rhp abiertos de $P(s) \hat{K}(s)$ es $n_{k}+n_{p}$;

(ii) $(I-P(s) \hat{K}(s))^{-1}$ es estable.

Prueba. Es fácil mostrar que $P \hat{K}$ y $(I-P \hat{K})^{-1}$ tienen las siguientes realizaciones:

$$
P \hat{K}=\left[\begin{array}{cc|c}
A & B \hat{C} & B \hat{D} \\
0 & \hat{A} & \hat{B} \\
\hline C & D \hat{C} & D \hat{D}
\end{array}\right], \quad(I-P \hat{K})^{-1}=\left[\begin{array}{c|c}
\bar{A} & \bar{B} \\
\hline \bar{C} & \bar{D}
\end{array}\right]
$$

donde

$$
\begin{aligned}
\bar{A} & =\left[\begin{array}{cc}
A & B \hat{C} \\
0 & \hat{A}
\end{array}\right]+\left[\begin{array}{c}
B \hat{D} \\
\hat{B}
\end{array}\right](I-D \hat{D})^{-1}\left[\begin{array}{ll}
C & D \hat{C}
\end{array}\right] \\
\bar{B} & =\left[\begin{array}{c}
B \hat{D} \\
\hat{B}
\end{array}\right](I-D \hat{D})^{-1} \\
\bar{C} & =(I-D \hat{D})^{-1}\left[\begin{array}{ll}
C & D \hat{C}
\end{array}\right] \\
\bar{D} & =(I-D \hat{D})^{-1}
\end{aligned}
$$

Por lo tanto, el sistema es internamente estable si y solo si $\bar{A}$ es estable (véase Problema 5.2).

Ahora suponga que el sistema es internamente estable; entonces $(I-P \hat{K})^{-1} \in \mathcal{R} \mathcal{H}_{\infty}$. Así, solo necesitamos mostrar que, dada la condición (ii), la condición (i) es necesaria y suficiente para la estabilidad interna. Esto se sigue al notar que $(\bar{A}, \bar{B})$ es estabilizable si y solo si

\[
\left(\left[\begin{array}{cc}
A & B \hat{C}  \tag{5.5}\\
0 & \hat{A}
\end{array}\right],\left[\begin{array}{c}
B \hat{D} \\
\hat{B}
\end{array}\right]\right)
\]

es estabilizable; y $(\bar{C}, \bar{A})$ es detectable si y solo si

\[
\left(\left[\begin{array}{ll}
C & D \hat{C}
\end{array}\right],\left[\begin{array}{cc}
A & B \hat{C}  \tag{5.6}\\
0 & \hat{A}
\end{array}\right]\right)
\]

es detectable. Pero las condiciones (5.5) y (5.6) son equivalentes a la condición (i).

La condición (i) del teorema anterior implica que no hay cancelación inestable polo/cero al formar el producto $P \hat{K}$.

De hecho, el teorema anterior es la base de la teoría clásica de control, donde la estabilidad se verifica solo en una función de transferencia de lazo cerrado, con la suposición implícita de que el controlador mismo es estable (y probablemente también de fase mínima; o al menos marginalmente estable y de fase mínima, con la condición de que cualquier polo sobre el eje imaginario del controlador no esté en la misma ubicación que algún cero de la planta).

Ejemplo 5.1 Sean $P$ y $\hat{K}$ matrices de transferencia de dos por dos

$$
P=\left[\begin{array}{cc}
\frac{1}{s-1} & 0 \\
0 & \frac{1}{s+1}
\end{array}\right], \quad \hat{K}=\left[\begin{array}{cc}
\frac{1-s}{s+1} & -1 \\
0 & -1
\end{array}\right]
$$

Entonces

$$
P \hat{K}=\left[\begin{array}{cc}
\frac{-1}{s+1} & \frac{-1}{s-1} \\
0 & \frac{-1}{s+1}
\end{array}\right], \quad(I-P \hat{K})^{-1}=\left[\begin{array}{cc}
\frac{s+1}{s+2} & -\frac{(s+1)^{2}}{(s+2)^{2}(s-1)} \\
0 & \frac{s+1}{s+2}
\end{array}\right]
$$

Así, el sistema en lazo cerrado no es estable, aunque

$$
\operatorname{det}(I-P \hat{K})=\frac{(s+2)^{2}}{(s+1)^{2}}
$$

no tiene ceros en el semiplano derecho cerrado y el número de polos inestables de $P \hat{K}$ es $n_{k}+n_{p}=1$. Por tanto, en general, que $\operatorname{det}(I-P \hat{K})$ no tenga ceros en el semiplano derecho cerrado no implica necesariamente que $(I-P \hat{K})^{-1} \in \mathcal{R} \mathcal{H}_{\infty}$.

\section{Factorización coprima sobre $\mathcal{R} \mathcal{H}_{\infty}$}
Recuerde que dos polinomios $m(s)$ y $n(s)$, con coeficientes reales por ejemplo, se dicen coprimos si su máximo común divisor es 1 (equivalentemente, no tienen ceros en común). De acuerdo con el algoritmo de Euclides ${ }^{1}$, dos polinomios $m$ y $n$ son coprimos si y solo si existen polinomios $x(s)$ y $y(s)$ tales que $x m+y n=1$; a esta ecuación se le llama identidad de Bezout. De manera similar, dos funciones de transferencia $m(s)$ y $n(s)$ en $\mathcal{R H}_{\infty}$ se dicen coprimas sobre $\mathcal{R} \mathcal{H}_{\infty}$ si existen $x, y \in \mathcal{R} \mathcal{H}_{\infty}$ tales que

$$
x m+y n=1
$$

La definición más primitiva, pero equivalente, es que $m$ y $n$ son coprimos si todo divisor común de $m$ y $n$ es invertible en $\mathcal{R} \mathcal{H}_{\infty}$; es decir,

$$
h, m h^{-1}, n h^{-1} \in \mathcal{R} \mathcal{H}_{\infty} \Longrightarrow h^{-1} \in \mathcal{R} \mathcal{H}_{\infty}
$$

Más generalmente, tenemos lo siguiente:

Definición 5.3 Dos matrices $M$ y $N$ en $\mathcal{R} \mathcal{H}_{\infty}$ son coprimas derechas sobre $\mathcal{R} \mathcal{H}_{\infty}$ si tienen el mismo número de columnas y si existen matrices $X_{r}$ y $Y_{r}$ en $\mathcal{R} \mathcal{H}_{\infty}$ tales que

$$
\left[\begin{array}{ll}
X_{r} & Y_{r}
\end{array}\right]\left[\begin{array}{l}
M \\
N
\end{array}\right]=X_{r} M+Y_{r} N=I
$$

De forma similar, dos matrices $\tilde{M}$ y $\tilde{N}$ en $\mathcal{R} \mathcal{H}_{\infty}$ son coprimas izquierdas sobre $\mathcal{R} \mathcal{H}_{\infty}$ si tienen el mismo número de filas y si existen matrices $X_{l}$ y $Y_{l}$ en $\mathcal{R} \mathcal{H}_{\infty}$ tales que

$$
\left[\begin{array}{ll}
\tilde{M} & \tilde{N}
\end{array}\right]\left[\begin{array}{c}
X_{l} \\
Y_{l}
\end{array}\right]=\tilde{M} X_{l}+\tilde{N} Y_{l}=I
$$

Nótese que estas definiciones son equivalentes a decir que la matriz $\left[\begin{array}{c}M \\ N\end{array}\right]$ es invertible por la izquierda en $\mathcal{R} \mathcal{H}_{\infty}$ y que la matriz $\left[\begin{array}{cc}\tilde{M} & \tilde{N}\end{array}\right]$ es invertible por la derecha en $\mathcal{R} \mathcal{H}_{\infty}$. Estas dos ecuaciones suelen llamarse identidades de Bezout.

Ahora sea $P$ una matriz racional real propia. Una factorización coprima derecha (rcf) de $P$ es una factorización $P=N M^{-1}$, donde $N$ y $M$ son coprimas derechas sobre $\mathcal{R} \mathcal{H}_{\infty}$. De manera similar, una factorización coprima izquierda (lcf) tiene la forma $P=\tilde{M}^{-1} \tilde{N}$, donde $\tilde{N}$ y $\tilde{M}$ son coprimas izquierdas sobre $\mathcal{R} \mathcal{H}_{\infty}$. Se dice que una matriz $P(s) \in \mathcal{R}_{p}(s)$ tiene factorización coprima doble si existen una factorización coprima derecha $P=N M^{-1}$, una factorización coprima izquierda $P=\tilde{M}^{-1} \tilde{N}$, y $X_{r}, Y_{r}, X_{l}, Y_{l} \in \mathcal{R} \mathcal{H}_{\infty}$ tales que

\[
\left[\begin{array}{cc}
X_{r} & Y_{r}  \tag{5.7}\\
-\tilde{N} & \tilde{M}
\end{array}\right]\left[\begin{array}{cc}
M & -Y_{l} \\
N & X_{l}
\end{array}\right]=I
\]

Por supuesto, implícito en estas definiciones está el requisito de que $M$ y $\tilde{M}$ sean cuadradas y no singulares.

Teorema 5.6 Suponga que $P(s)$ es una matriz racional real propia y

$$
P=\left[\begin{array}{l|l}
A & B \\
\hline C & D
\end{array}\right]
$$

es una realización estabilizable y detectable. Sean $F$ y $L$ tales que $A+B F$ y $A+L C$ sean ambos estables, y defina

\[
\left[\begin{array}{cc}
M & -Y_{l}  \tag{5.8}\\
N & X_{l}
\end{array}\right]=\left[\begin{array}{c|cc}
A+B F & B & -L \\
\hline F & I & 0 \\
C+D F & D & I
\end{array}\right]
\]

\[
\left[\begin{array}{cc}
X_{r} & Y_{r}  \tag{5.9}\\
-\tilde{N} & \tilde{M}
\end{array}\right]=\left[\begin{array}{c|cc}
A+L C & -(B+L D) & L \\
\hline F & I & 0 \\
C & -D & I
\end{array}\right]
\]

Entonces $P=N M^{-1}=\tilde{M}^{-1} \tilde{N}$ son, respectivamente, rcf y lcf y, además, se satisface la ecuación (5.7).

Prueba. El teorema se sigue verificando la ecuación (5.7).

Observación 5.2 Nótese que si $P$ es estable, entonces podemos tomar $X_{r}=X_{l}=I, Y_{r}=Y_{l}=0$, $N=\tilde{N}=P, M=\tilde{M}=I$.

Observación 5.3 La factorización coprima de una matriz de transferencia puede interpretarse desde el punto de vista de control por realimentación. Por ejemplo, la factorización coprima derecha surge naturalmente al cambiar la variable de control mediante realimentación de estado. Considere las ecuaciones en espacio de estados de una planta $P$:

$$
\begin{aligned}
\dot{x} & =A x+B u \\
y & =C x+D u
\end{aligned}
$$

A continuación, introduzca una realimentación de estado y cambie la variable

$$
v:=u-F x
$$

donde $F$ es tal que $A+B F$ es estable. Entonces obtenemos

$$
\begin{aligned}
\dot{x} & =(A+B F) x+B v \\
u & =F x+v \\
y & =(C+D F) x+D v
\end{aligned}
$$

Evidentemente, de estas ecuaciones, la matriz de transferencia de $v$ a $u$ es

$$
M(s)=\left[\begin{array}{c|c}
A+B F & B \\
\hline F & I
\end{array}\right]
$$

y la de $v$ a $y$ es

$$
N(s)=\left[\begin{array}{c|c}
A+B F & B \\
\hline C+D F & D
\end{array}\right]
$$

Por lo tanto,

$$
u=M v, \quad y=N v
$$

de modo que $y=N M^{-1} u$; es decir, $P=N M^{-1}$.

Ahora veremos cómo las factorizaciones coprimas pueden usarse para obtener caracterizaciones alternativas de las condiciones de estabilidad interna. Considere de nuevo el diagrama estándar de análisis de estabilidad de la Figura 5.2. Comenzamos con cualquier rcf y lcf de $P$ y $\hat{K}$:


\begin{gather*}
P=N M^{-1}=\tilde{M}^{-1} \tilde{N}  \tag{5.10}\\
\hat{K}=U V^{-1}=\tilde{V}^{-1} \tilde{U} \tag{5.11}
\end{gather*}


Lema 5.7 Considere el sistema en la Figura 5.2. Las siguientes condiciones son equivalentes:

\begin{enumerate}
  \item El sistema de realimentación es internamente estable.
  \item $\left[\begin{array}{cc}M & U \\ N & V\end{array}\right]$ es invertible en $\mathcal{R} \mathcal{H}_{\infty}$.
  \item $\left[\begin{array}{cc}\tilde{V} & -\tilde{U} \\ -\tilde{N} & \tilde{M}\end{array}\right]$ es invertible en $\mathcal{R} \mathcal{H}_{\infty}$.
  \item $\tilde{M} V-\tilde{N} U$ es invertible en $\mathcal{R} \mathcal{H}_{\infty}$.
  \item $\tilde{V} M-\tilde{U} N$ es invertible en $\mathcal{R} \mathcal{H}_{\infty}$.
\end{enumerate}

Prueba. Nótese que el sistema es internamente estable si

$$
\left[\begin{array}{cc}
I & -\hat{K} \\
-P & I
\end{array}\right]^{-1} \in \mathcal{R} \mathcal{H}_{\infty}
$$

o, equivalentemente,

\[
\left[\begin{array}{cc}
I & \hat{K}  \tag{5.12}\\
P & I
\end{array}\right]^{-1} \in \mathcal{R} \mathcal{H}_{\infty}
\]

Ahora

$$
\left[\begin{array}{cc}
I & \hat{K} \\
P & I
\end{array}\right]=\left[\begin{array}{cc}
I & U V^{-1} \\
N M^{-1} & I
\end{array}\right]=\left[\begin{array}{cc}
M & U \\
N & V
\end{array}\right]\left[\begin{array}{cc}
M^{-1} & 0 \\
0 & V^{-1}
\end{array}\right]
$$

por lo que

$$
\left[\begin{array}{cc}
I & \hat{K} \\
P & I
\end{array}\right]^{-1}=\left[\begin{array}{cc}
M & 0 \\
0 & V
\end{array}\right]\left[\begin{array}{ll}
M & U \\
N & V
\end{array}\right]^{-1}
$$

Como las matrices

$$
\left[\begin{array}{cc}
M & 0 \\
0 & V
\end{array}\right],\left[\begin{array}{cc}
M & U \\
N & V
\end{array}\right]
$$

son coprimas derechas (este hecho se deja como ejercicio), la ecuación (5.12) se cumple si y solo si

$$
\left[\begin{array}{ll}
M & U \\
N & V
\end{array}\right]^{-1} \in \mathcal{R} \mathcal{H}_{\infty}
$$

Esto prueba la equivalencia de las condiciones 1 y 2. La equivalencia entre 1 y 3 se prueba de forma similar.

Las condiciones 4 y 5 se deducen de 2 y 3 a partir de la ecuación

$$
\left[\begin{array}{cc}
\tilde{V} & -\tilde{U} \\
-\tilde{N} & \tilde{M}
\end{array}\right]\left[\begin{array}{cc}
M & U \\
N & V
\end{array}\right]=\left[\begin{array}{cc}
\tilde{V} M-\tilde{U} N & 0 \\
0 & \tilde{M} V-\tilde{N} U
\end{array}\right]
$$

Como el lado izquierdo de la ecuación anterior es invertible en $\mathcal{R} \mathcal{H}_{\infty}$, el lado derecho también lo es. Por tanto, se satisfacen las condiciones 4 y 5. Solo falta mostrar que cualquiera de 4 o 5 implica 1. Mostremos que 5 implica 1; esto es obvio pues

$$
\begin{aligned}
{\left[\begin{array}{cc}
I & \hat{K} \\
P & I
\end{array}\right]^{-1} } & =\left[\begin{array}{cc}
I & \tilde{V}^{-1} \tilde{U} \\
N M^{-1} & I
\end{array}\right]^{-1} \\
& =\left[\begin{array}{cc}
M & 0 \\
0 & I
\end{array}\right]\left[\begin{array}{cc}
\tilde{V} M & \tilde{U} \\
N & I
\end{array}\right]^{-1}\left[\begin{array}{cc}
\tilde{V} & 0 \\
0 & I
\end{array}\right] \in \mathcal{R} \mathcal{H}_{\infty}
\end{aligned}
$$

si $\left[\begin{array}{cc}\tilde{V} M & \tilde{U} \\ N & I\end{array}\right]^{-1} \in \mathcal{R} \mathcal{H}_{\infty}$ o si se satisface la condición 5.

Combinando el Lema 5.7 y el Teorema 5.6, obtenemos el siguiente corolario.

Corolario 5.8 Sea $P$ una matriz racional real propia y sea $P=N M^{-1}=\tilde{M}^{-1} \tilde{N}$ la correspondiente rcf y lcf sobre $\mathcal{R} \mathcal{H}_{\infty}$. Entonces existe un controlador

$$
\hat{K}_{0}=U_{0} V_{0}^{-1}=\tilde{V}_{0}^{-1} \tilde{U}_{0}
$$

con $U_{0}, V_{0}, \tilde{U}_{0}$ y $\tilde{V}_{0}$ en $\mathcal{R} \mathcal{H}_{\infty}$ tal que

\[
\left[\begin{array}{cc}
\tilde{V}_{0} & -\tilde{U}_{0}  \tag{5.13}\\
-\tilde{N} & \tilde{M}
\end{array}\right]\left[\begin{array}{cc}
M & U_{0} \\
N & V_{0}
\end{array}\right]=\left[\begin{array}{ll}
I & 0 \\
0 & I
\end{array}\right]
\]

Además, sean $F$ y $L$ tales que $A+B F$ y $A+L C$ sean estables. Entonces un conjunto particular de realizaciones en espacio de estados para estas matrices puede darse por


\begin{gather*}
{\left[\begin{array}{cc}
M & U_{0} \\
N & V_{0}
\end{array}\right]=\left[\begin{array}{c|cc}
A+B F & B & -L \\
\hline F & I & 0 \\
C+D F & D & I
\end{array}\right]}  \tag{5.14}\\
{\left[\begin{array}{cc}
\tilde{V}_{0} & -\tilde{U}_{0} \\
-\tilde{N} & \tilde{M}
\end{array}\right]=\left[\begin{array}{c|cc}
A+L C & -(B+L D) & L \\
\hline F & I & 0 \\
C & -D & I
\end{array}\right]} \tag{5.15}
\end{gather*}


Prueba. La idea detrás de la elección de estas matrices es la siguiente. Usando teoría de observadores, encuentre un controlador $\hat{K}_{0}$ que logre estabilidad interna; por ejemplo,

\[
\hat{K}_{0}:=\left[\begin{array}{c|c}
A+B F+L C+L D F & -L  \tag{5.16}\\
\hline F & 0
\end{array}\right]
\]

Realice las factorizaciones

$$
\hat{K}_{0}=U_{0} V_{0}^{-1}=\tilde{V}_{0}^{-1} \tilde{U}_{0}
$$

que son análogas a las realizadas sobre $P$. Entonces el Lema 5.7 implica que cada una de las dos matrices por bloques del lado izquierdo de la ecuación (5.13) debe ser invertible en $\mathcal{R H}_{\infty}$. De hecho, la ecuación (5.13) se satisface comparándola con la ecuación (5.7).

Encontrar una factorización coprima para una función de transferencia escalar es relativamente fácil. Sea $P(s)=\operatorname{num}(s) / \operatorname{den}(s)$, donde $\operatorname{num}(s)$ y $\operatorname{den}(s)$ son los polinomios numerador y denominador de $P(s)$, y sea $\alpha(s)$ un polinomio estable del mismo orden que $\operatorname{den}(s)$. Entonces $P(s)=n(s) / m(s)$ con $n(s)=\operatorname{num}(s) / \alpha(s)$ y $m(s)=\operatorname{den}(s) / \alpha(s)$ es una factorización coprima. Sin embargo, encontrar $x(s) \in \mathcal{H}_{\infty}$ y $y(s) \in \mathcal{H}_{\infty}$ tales que $x(s) n(s)+y(s) m(s)=1$ requiere mucho más trabajo.

Ejemplo 5.2 Sea $P(s)=\frac{s-2}{s(s+3)}$ y $\alpha=(s+1)(s+3)$. Entonces $P(s)=n(s) / m(s)$ con $n(s)=\frac{s-2}{(s+1)(s+3)}$ y $m(s)=\frac{s}{s+1}$ forma una factorización coprima. Para encontrar $x(s) \in \mathcal{H}_{\infty}$ y $y(s) \in \mathcal{H}_{\infty}$ tales que $x(s) n(s)+y(s) m(s)=1$, considere un controlador estabilizante para $P$: $\hat{K}=-\frac{s-1}{s+10}$. Entonces $\hat{K}=u / v$ con $u=\hat{K}$ y $v=1$ es una factorización coprima y

$$
m(s) v(s)-n(s) u(s)=\frac{(s+11.7085)(s+2.214)(s+0.077)}{(s+1)(s+3)(s+10)}=: \beta(s)
$$

Entonces podemos tomar

$$
\begin{aligned}
& x(s)=-u(s) / \beta(s)=\frac{(s-1)(s+1)(s+3)}{(s+11.7085)(s+2.214)(s+0.077)} \\
& y(s)=v(s) / \beta(s)=\frac{(s+1)(s+3)(s+10)}{(s+11.7085)(s+2.214)(s+0.077)}
\end{aligned}
$$

Pueden usarse programas de MATLAB para encontrar las matrices adecuadas $F$ y $L$ en espacio de estados, de modo que se obtenga la factorización coprima deseada. Sea $A \in \mathbb{R}^{n \times n}, B \in \mathbb{R}^{n \times m}$ y $C \in \mathbb{R}^{p \times n}$. Entonces $F$ y $L$ pueden obtenerse mediante

$\gg \mathrm{F}=-\operatorname{lqr}(\mathrm{A}, \mathrm{B}, \text{eye}(n), \text{eye}(m)); \%$ o

$\gg \mathbf{F}=-\operatorname{place}(A, B, Pf); \% Pf =$ polos de $\mathrm{A}+\mathrm{BF}$

$\gg \mathbf{L}=-\operatorname{lqr}\left(\mathbf{A}^{\prime}, \mathbf{C}^{\prime}, \text{eye}(\mathbf{n}), \text{eye}(\mathbf{p})\right)^{\prime} ; \%$ o

$\gg \mathbf{L}=-\operatorname{place}\left(\mathbf{A}^{\prime}, \mathbf{C}^{\prime}, \mathbf{Pl}\right)^{\prime} ; \% \mathrm{Pl}=$ polos de $\mathrm{A}+\mathrm{LC}$.

\section{Notas y referencias}
La presentación de este capítulo se basa principalmente en Doyle [1984]. La discusión sobre estabilidad interna y factorización coprima también puede encontrarse en Francis [1987], Vidyasagar [1985], y Nett, Jacobson y Balas [1984].

\section{Problemas}
Problema 5.1 Recuerde que un sistema de realimentación se dice internamente estable si todas las funciones de transferencia en lazo cerrado son estables. Describa las condiciones de estabilidad interna del siguiente sistema de realimentación:

\insertimage[\label{img:problema-5-1-esquema-realimentacion}]{2024_06_29_ce7c73b22613114bd4d6g-091_268_605_1151_760}{width=0.45\textwidth}{Esquema del sistema de realimentación para el Problema 5.1.}

¿Cómo pueden simplificarse las condiciones de estabilidad si $H(s)$ y $G_{1}(s)$ son ambos estables?

Problema 5.2 Demuestre que $\left[\begin{array}{cc}I & -\hat{K} \\ -P & I\end{array}\right]^{-1} \in \mathcal{R} \mathcal{H}_{\infty}$ si y solo si

$$
\bar{A}:=\left[\begin{array}{cc}
A & B \hat{C} \\
0 & \hat{A}
\end{array}\right]+\left[\begin{array}{c}
B \hat{D} \\
\hat{B}
\end{array}\right](I-D \hat{D})^{-1}\left[\begin{array}{ll}
C & D \hat{C}
\end{array}\right]
$$

es estable.

Problema 5.3 Suponga que $N, M, U, V \in \mathcal{R} \mathcal{H}_{\infty}$ y que $N M^{-1}$ y $U V^{-1}$ son factorizaciones coprimas derechas, respectivamente. Demuestre que

$$
\left[\begin{array}{cc}
M & 0 \\
0 & V
\end{array}\right]\left[\begin{array}{ll}
M & U \\
N & V
\end{array}\right]^{-1}
$$

también es una factorización coprima derecha.

Problema 5.4 Sea $G(s)=\frac{s-1}{(s+2)(s-3)}$. Encuentre una factorización coprima estable $G=n(s)/m(s)$ y $x, y \in \mathcal{R} \mathcal{H}_{\infty}$ tales que $x n+y m=1$.

Problema 5.5 Sea $N(s)=\frac{(s-1)(s+\alpha)}{(s+2)(s+3)(s+\beta)}$ y $M(s)=\frac{(s-3)(s+\alpha)}{(s+3)(s+\beta)}$. Demuestre que $(N, M)$ también es una factorización coprima de $G$ del Problema 5.4 para cualquier $\alpha>0$ y $\beta>0$.

Problema 5.6 Sea $G=N M^{-1}$ una factorización coprima derecha sobre $\mathcal{R} \mathcal{H}_{\infty}$. Se llama factorización coprima normalizada si $N^{\sim} N+M^{\sim} M=I$. Ahora considere una función de transferencia escalar $G$. Entonces puede usarse el siguiente procedimiento para encontrar una factorización coprima normalizada: (a) Sea $G=n/m$ cualquier factorización coprima sobre $\mathcal{R} \mathcal{H}_{\infty}$. (b) Encuentre un factor espectral $w$ estable y de fase mínima tal que $w^{\sim} w=n^{\sim} n+m^{\sim} m$. Sea $N=n/w$ y $M=m/w$; entonces $G=N/M$ es una factorización coprima normalizada. Encuentre una factorización coprima normalizada para el Problema 5.4.

Problema 5.7 El siguiente procedimiento construye una factorización coprima derecha normalizada cuando $G$ es estrictamente propia:

\begin{enumerate}
  \item Obtenga una realización estabilizable y detectable $A, B, C$.
  \item Ejecute en MATLAB el comando $F=-\operatorname{lqr}\left(A, B, C^{\prime} C, I\right)$.
  \item Defina
\end{enumerate}

$$
\left[\begin{array}{c}
N \\
M
\end{array}\right](s)=\left[\begin{array}{c|c}
A+B F & B \\
\hline C & 0 \\
F & I
\end{array}\right]
$$

Verifique que el procedimiento produce factores que satisfacen $G=N M^{-1}$. Ahora aplique el procedimiento a

$$
G(s)=\left[\begin{array}{cc}
\frac{1}{s-1} & \frac{1}{s-2} \\
\frac{2}{s} & \frac{1}{s+2}
\end{array}\right]
$$

Verifique numéricamente que


\begin{equation*}
N(j \omega)^{*} N(j \omega)+M(j \omega)^{*} M(j \omega)=I, \quad \forall \omega \tag{5.17}
\end{equation*}


Problema 5.8 Use el procedimiento del Problema 5.7 para encontrar la factorización coprima derecha normalizada de

$$
\begin{gathered}
G_{1}(s)=\left[\begin{array}{cc}
\frac{1}{s+1} & \frac{s+3}{(s+1)(s-2)} \\
\frac{10}{s-2} & \frac{5}{s+3}
\end{array}\right] \\
G_{2}(s)=\left[\begin{array}{ll}
\frac{2(s+1)(s+2)}{s(s+3)(s+4)} & \frac{s+2}{(s+1)(s+3)}
\end{array}\right]
\end{gathered}
$$

$$
\begin{gathered}
G_{3}(s)=\left[\begin{array}{ccc|ccc}
-1 & -2 & 1 & 1 & 2 & 3 \\
0 & 2 & -1 & 3 & 2 & 1 \\
-4 & -3 & -2 & 1 & 1 & 1 \\
\hline 1 & 1 & 1 & 0 & 0 & 0 \\
2 & 3 & 4 & 0 & 0 & 0
\end{array}\right] \\
G_{4}(s)=\left[\begin{array}{cc|cc}
-1 & -2 & 1 & 2 \\
0 & 1 & 2 & 1 \\
\hline 1 & 1 & 0 & 0 \\
1 & 1 & 0 & 0
\end{array}\right]
\end{gathered}
$$

Problema 5.9 Defina la factorización coprima izquierda normalizada y describa un procedimiento para encontrar tales factorizaciones en matrices de transferencia estrictamente propias.


\chapter{Especificaciones de desempeño y limitaciones}
En este capítulo, consideramos con mayor profundidad las propiedades de los sistemas con realimentación y discutimos cómo lograr el desempeño deseado mediante control por realimentación. También consideramos las formulaciones matemáticas de los problemas de control óptimo $\mathcal{H}_{2}$ y $\mathcal{H}_{\infty}$. Un paso clave en el diseño de control óptimo es la selección de funciones de ponderación. Daremos algunas guías para ese proceso de selección usando ejemplos SISO. También discutiremos con cierto detalle las limitaciones de diseño impuestas por las restricciones de ancho de banda, por los ceros en semiplano derecho del lazo abierto y por los polos en semiplano derecho del lazo abierto, usando la relación ganancia-fase de Bode, la relación integral de sensibilidad de Bode y la fórmula integral de Poisson.

\section{Propiedades de retroalimentación}
En esta sección, discutimos las propiedades de un sistema con realimentación. En particular, consideramos el beneficio de la estructura de realimentación y el concepto de compromisos de diseño para objetivos en conflicto, es decir, cómo lograr los beneficios de la realimentación en presencia de incertidumbres.

\insertimage[\label{img:configuracion-estandar-realimentacion-cap6}]{2024_06_29_ce7c73b22613114bd4d6g-095_282_876_1816_608}{width=0.35\textwidth}{Configuración estándar de realimentación.}

Considere de nuevo el sistema con realimentación mostrado en la Figura 5.1. Por conveniencia, el diagrama del sistema se muestra nuevamente en la Figura 6.1. Para continuar la discusión, es conveniente definir la matriz de transferencia de lazo en la entrada, $L_{i}$, y la matriz de transferencia de lazo en la salida, $L_{o}$, como

$$
L_{i}=K P, \quad L_{o}=P K
$$

respectivamente, donde $L_{i}$ se obtiene al abrir el lazo en la entrada $(u)$ de la planta, mientras que $L_{o}$ se obtiene al abrir el lazo en la salida $(y)$ de la planta. La matriz de sensibilidad de entrada se define como la matriz de transferencia desde $d_{i}$ hasta $u_{p}$:

$$
S_{i}=\left(I+L_{i}\right)^{-1}, \quad u_{p}=S_{i} d_{i}
$$

La matriz de sensibilidad de salida se define como la matriz de transferencia desde $d$ hasta $y$:

$$
S_{o}=\left(I+L_{o}\right)^{-1}, \quad y=S_{o} d
$$

Las matrices de sensibilidad complementaria de entrada y salida se definen como

$$
\begin{gathered}
T_{i}=I-S_{i}=L_{i}\left(I+L_{i}\right)^{-1} \\
T_{o}=I-S_{o}=L_{o}\left(I+L_{o}\right)^{-1}
\end{gathered}
$$

respectivamente. (La palabra complementaria se usa para indicar que $T$ es el complemento de $S$, es decir, $T=I-S$.) La matriz $I+L_{i}$ se llama matriz de diferencia de retorno de entrada, y $I+L_{o}$ se llama matriz de diferencia de retorno de salida.

Es fácil ver que el sistema en lazo cerrado, si es internamente estable, satisface las siguientes ecuaciones:


\begin{align*}
y & =T_{o}(r-n)+S_{o} P d_{i}+S_{o} d  \tag{6.1}\\
r-y & =S_{o}(r-d)+T_{o} n-S_{o} P d_{i}  \tag{6.2}\\
u & =K S_{o}(r-n)-K S_{o} d-T_{i} d_{i}  \tag{6.3}\\
u_{p} & =K S_{o}(r-n)-K S_{o} d+S_{i} d_{i} \tag{6.4}
\end{align*}


Estas cuatro ecuaciones muestran los beneficios fundamentales y los objetivos de diseño inherentes en lazos de realimentación. Por ejemplo, la ecuación (6.1) muestra que los efectos de la perturbación $d$ sobre la salida de la planta pueden hacerse "pequeños" haciendo pequeña la función de sensibilidad de salida $S_{o}$. De forma similar, la ecuación (6.4) muestra que los efectos de la perturbación $d_{i}$ sobre la entrada de la planta pueden hacerse pequeños haciendo pequeña la función de sensibilidad de entrada $S_{i}$. La noción de pequeñez de una matriz de transferencia en cierto rango de frecuencias puede hacerse explícita usando valores singulares dependientes de la frecuencia; por ejemplo, $\bar{\sigma}\left(S_{o}\right)<1$ sobre un rango de frecuencias significaría que los efectos de la perturbación $d$ en la salida de la planta están efectivamente desensibilizados en ese rango de frecuencias.

Por lo tanto, un buen rechazo de perturbaciones en la salida de la planta $(y)$ requeriría que

$$
\begin{aligned}
\bar{\sigma}\left(S_{o}\right) & \left.=\bar{\sigma}\left((I+P K)^{-1}\right)=\frac{1}{\underline{\sigma}(I+P K)} \quad \text { (para perturbación en salida de planta, } d\right) \\
\bar{\sigma}\left(S_{o} P\right) & \left.=\bar{\sigma}\left((I+P K)^{-1} P\right)=\bar{\sigma}\left(P S_{i}\right) \quad \text { (para perturbación en entrada de planta, } d_{i}\right)
\end{aligned}
$$

sean pequeñas, y un buen rechazo de perturbaciones en la entrada de la planta $\left(u_{p}\right)$ requeriría que

$$
\begin{aligned}
\bar{\sigma}\left(S_{i}\right) & \left.=\bar{\sigma}\left((I+K P)^{-1}\right)=\frac{1}{\underline{\sigma}(I+K P)} \quad \text { (para perturbación en entrada de planta, } d_{i}\right) \\
\bar{\sigma}\left(S_{i} K\right) & \left.=\bar{\sigma}\left(K(I+P K)^{-1}\right)=\bar{\sigma}\left(K S_{o}\right) \quad \text { (para perturbación en salida de planta, } d\right)
\end{aligned}
$$

sean pequeñas, particularmente en el rango de bajas frecuencias donde $d$ y $d_{i}$ suelen ser significativas.

Nótese que

$$
\begin{aligned}
& \underline{\sigma}(P K)-1 \leq \underline{\sigma}(I+P K) \leq \underline{\sigma}(P K)+1 \\
& \underline{\sigma}(K P)-1 \leq \underline{\sigma}(I+K P) \leq \underline{\sigma}(K P)+1
\end{aligned}
$$

entonces

$$
\begin{aligned}
& \frac{1}{\underline{\sigma}(P K)+1} \leq \bar{\sigma}\left(S_{o}\right) \leq \frac{1}{\underline{\sigma}(P K)-1}, \quad \text { si } \quad \underline{\sigma}(P K)>1 \\
& \frac{1}{\underline{\sigma}(K P)+1} \leq \bar{\sigma}\left(S_{i}\right) \leq \frac{1}{\underline{\sigma}(K P)-1}, \quad \text { si } \quad \underline{\sigma}(K P)>1
\end{aligned}
$$

Estas ecuaciones implican que

$$
\begin{aligned}
& \bar{\sigma}\left(S_{o}\right) \ll 1 \Longleftrightarrow \underline{\sigma}(P K) \gg 1 \\
& \bar{\sigma}\left(S_{i}\right) \ll 1 \Longleftrightarrow \underline{\sigma}(K P) \gg 1
\end{aligned}
$$

Ahora suponga que $P$ y $K$ son invertibles; entonces

$$
\begin{aligned}
& \underline{\sigma}(P K) \gg 1 \text { o } \underline{\sigma}(K P) \gg 1 \Longleftrightarrow \bar{\sigma}\left(S_{o} P\right)=\bar{\sigma}\left((I+P K)^{-1} P\right) \approx \bar{\sigma}\left(K^{-1}\right)=\frac{1}{\underline{\sigma}(K)} \\
& \underline{\sigma}(P K) \gg 1 \text { o } \underline{\sigma}(K P) \gg 1 \Longleftrightarrow \bar{\sigma}\left(K S_{o}\right)=\bar{\sigma}\left(K(I+P K)^{-1}\right) \approx \bar{\sigma}\left(P^{-1}\right)=\frac{1}{\underline{\sigma}(P)}
\end{aligned}
$$

Por lo tanto, un buen desempeño en la salida de planta $(y)$ requiere, en general, una gran ganancia de lazo de salida $\underline{\sigma}\left(L_{o}\right)=\underline{\sigma}(P K) \gg 1$ en el rango de frecuencias donde $d$ es significativa para desensibilizar $d$, y una ganancia del controlador suficientemente grande $\underline{\sigma}(K) \gg 1$ en el rango de frecuencias donde $d_{i}$ es significativa para desensibilizar $d_{i}$. De forma similar, un buen desempeño en la entrada de planta $\left(u_{p}\right)$ requiere, en general, gran ganancia de lazo de entrada $\underline{\sigma}\left(L_{i}\right)=\underline{\sigma}(K P) \gg 1$ en el rango de frecuencias donde $d_{i}$ es significativa para desensibilizar $d_{i}$, y ganancia de planta suficientemente grande $\underline{\sigma}(P) \gg 1$ en el rango de frecuencias donde $d$ es significativa, la cual no puede cambiarse mediante diseño del controlador, para desensibilizar $d$. [En general, $S_{o} \neq S_{i}$ salvo que $K$ y $P$ sean cuadradas y diagonales, lo cual es cierto si $P$ es un sistema escalar. Por lo tanto, un $\bar{\sigma}\left(S_{o}\right)$ pequeño no implica necesariamente un $\bar{\sigma}\left(S_{i}\right)$ pequeño; en otras palabras, buen rechazo de perturbaciones en la salida no implica necesariamente buen rechazo de perturbaciones en la entrada de planta.]

Así, el buen diseño de lazo de realimentación multivariable se reduce a lograr altas ganancias de lazo (y posiblemente del controlador) en el rango de frecuencia necesario.

A pesar de la simplicidad de esta afirmación, el diseño por realimentación no es en absoluto trivial. Esto ocurre porque las ganancias de lazo no pueden hacerse arbitrariamente altas sobre rangos de frecuencia arbitrariamente amplios. Más bien, deben satisfacer ciertos compromisos de desempeño y limitaciones de diseño. Un compromiso importante de desempeño, por ejemplo, concierne la reducción de errores por referencia/perturbación frente a la estabilidad bajo incertidumbre de modelo. Suponga que el modelo de planta se perturba a $(I+\Delta) P$ con $\Delta$ estable, y suponga que el sistema nominal es estable (es decir, el sistema en lazo cerrado con $\Delta=0$ es estable). Entonces el sistema perturbado en lazo cerrado es estable si

$$
\operatorname{det}(I+(I+\Delta) P K)=\operatorname{det}(I+P K) \operatorname{det}\left(I+\Delta T_{o}\right)
$$

no tiene ceros en semiplano derecho. Esto, en general, equivale a requerir que $\left\|\Delta T_{o}\right\|$ sea pequeño o que $\bar{\sigma}\left(T_{o}\right)$ sea pequeño en aquellas frecuencias donde $\Delta$ es significativa, típicamente en alta frecuencia, lo cual a su vez implica que la ganancia de lazo, $\bar{\sigma}\left(L_{o}\right)$, debe ser pequeña en esas frecuencias.

Otro compromiso adicional aparece con la reducción del error por ruido de sensor. El conflicto entre rechazo de perturbaciones y reducción de ruido de sensor es evidente en la ecuación (6.1). Valores grandes de $\underline{\sigma}\left(L_{o}(j \omega)\right)$ sobre un amplio rango de frecuencias hacen pequeño el error debido a $d$. Sin embargo, también hacen grande el error debido a $n$ porque ese ruido "pasa" en ese mismo rango de frecuencias, es decir,

$$
y=T_{o}(r-n)+S_{o} P d_{i}+S_{o} d \approx(r-n)
$$

Nótese que $n$ suele ser significativo en el rango de alta frecuencia. Peor aún, ganancias de lazo grandes fuera del ancho de banda de $P$, es decir, $\underline{\sigma}\left(L_{o}(j \omega)\right) \gg 1$ o $\underline{\sigma}\left(L_{i}(j \omega)\right) \gg 1$ mientras $\bar{\sigma}(P(j \omega)) \ll 1$, pueden hacer que la actividad de control $(u)$ sea inaceptable, lo que puede provocar saturación de actuadores. Esto se sigue de

$$
u=K S_{o}(r-n-d)-T_{i} d_{i}=S_{i} K(r-n-d)-T_{i} d_{i} \approx P^{-1}(r-n-d)-d_{i}
$$

Aquí hemos supuesto, por conveniencia, que $P$ es cuadrada e invertible. La ecuación resultante muestra que las perturbaciones y el ruido de sensor en realidad se amplifican en $u$ cuando el rango de frecuencia excede significativamente el ancho de banda de $P$, ya que para $\omega$ tal que $\bar{\sigma}(P(j \omega)) \ll 1$ se tiene

$$
\underline{\sigma}\left[P^{-1}(j \omega)\right]=\frac{1}{\bar{\sigma}[P(j \omega)]} \gg 1
$$

De forma similar, la ganancia del controlador, $\bar{\sigma}(K)$, también debe mantenerse no demasiado grande en el rango de frecuencias donde la ganancia de lazo es pequeña, para no saturar los actuadores. Esto se debe a que, para ganancia de lazo pequeña $\bar{\sigma}\left(L_{o}(j \omega)\right) \ll 1$ o $\bar{\sigma}\left(L_{i}(j \omega)\right) \ll 1$,

$$
u=K S_{o}(r-n-d)-T_{i} d_{i} \approx K(r-n-d)
$$

Por lo tanto, es deseable mantener $\bar{\sigma}(K)$ no demasiado grande cuando la ganancia de lazo es pequeña.

Para resumir la discusión anterior, observamos que un buen desempeño requiere, en cierto rango de frecuencias, típicamente un rango de bajas frecuencias $\left(0, \omega_{l}\right)$,

$$
\underline{\sigma}(P K) \gg 1, \quad \underline{\sigma}(K P) \gg 1, \quad \underline{\sigma}(K) \gg 1
$$

y una buena robustez y buen rechazo de ruido de sensor requieren, en cierto rango de frecuencias, típicamente un rango de altas frecuencias $\left(\omega_{h}, \infty\right)$,

$$
\bar{\sigma}(P K) \ll 1, \quad \bar{\sigma}(K P) \ll 1, \quad \bar{\sigma}(K) \leq M
$$

donde $M$ no es demasiado grande. Estos requerimientos de diseño se muestran gráficamente en la Figura 6.2. Las frecuencias específicas $\omega_{l}$ y $\omega_{h}$ dependen de la aplicación particular y del conocimiento que se tenga sobre las características de perturbación, las incertidumbres de modelado y los niveles de ruido de sensor.

\insertimage[\label{img:ganancia-lazo-deseada}]{2024_06_29_ce7c73b22613114bd4d6g-099_681_1021_817_541}{width=0.35\textwidth}{Ganancia de lazo deseada.}

\section{Desempeño ponderado $\mathcal{H}_2$ y $\mathcal{H}_{\infty}$}
En esta sección, consideramos cómo formular algunos objetivos de desempeño en problemas matemáticamente tratables. Como se mostró en la Sección 6.1, los objetivos de desempeño de un sistema con realimentación suelen especificarse en términos de requisitos sobre funciones de sensibilidad y/o sensibilidad complementaria, o en términos de otras funciones de transferencia en lazo cerrado. Por ejemplo, para un sistema escalar, los criterios de desempeño pueden especificarse exigiendo

$$
\begin{cases}|S(j \omega)| \leq \varepsilon, & \forall \omega \leq \omega_{0} \\ |S(j \omega)| \leq M, & \forall \omega>\omega_{0}\end{cases}
$$

donde $S(j \omega)=1 /(1+P(j \omega) K(j \omega))$. Sin embargo, es mucho más conveniente reflejar los objetivos de desempeño del sistema eligiendo funciones de ponderación apropiadas. Por ejemplo, el objetivo de desempeño anterior puede escribirse como

$$
\left|W_{e}(j \omega) S(j \omega)\right| \leq 1, \quad \forall \omega
$$

con

$$
\left|W_{e}(j \omega)\right|=\left\{\begin{array}{cc}
1 / \varepsilon, & \forall \omega \leq \omega_{0} \\
1 / M, & \forall \omega>\omega_{0}
\end{array}\right.
$$

Para usar $W_{e}$ en el diseño de control, normalmente se emplea una función de transferencia racional $W_{e}(s)$ que aproxime la respuesta en frecuencia anterior.

La ventaja de usar especificaciones de desempeño ponderadas es evidente en el diseño multivariable. Primero, algunos componentes de una señal vectorial suelen ser más importantes que otros. Segundo, cada componente de la señal puede no medirse en las mismas unidades; por ejemplo, algunos componentes de la señal de error de salida pueden medirse en términos de longitud y otros en términos de voltaje. Por tanto, las funciones de ponderación son esenciales para hacer comparables estos componentes. Además, puede interesarnos principalmente rechazar errores en cierto rango de frecuencias (por ejemplo, bajas frecuencias); por ello deben elegirse ponderaciones dependientes de la frecuencia.

\insertimage[\label{img:configuracion-realimentacion-con-ponderaciones}]{2024_06_29_ce7c73b22613114bd4d6g-100_486_1182_1123_466}{width=0.35\textwidth}{Configuración estándar de realimentación con ponderaciones.}

En general, modificaremos el diagrama estándar de la Figura 6.1 al de la Figura 6.3. Las funciones de ponderación en la Figura 6.3 se eligen para reflejar los objetivos de diseño y el conocimiento de perturbaciones y ruido de sensor. Por ejemplo, $W_{d}$ y $W_{i}$ pueden elegirse para reflejar el contenido en frecuencia de las perturbaciones $d$ y $d_{i}$, o pueden usarse para modelar el espectro de potencia de perturbación según la naturaleza de las señales involucradas en los sistemas prácticos. La matriz de ponderación $W_{n}$ se usa para modelar el contenido en frecuencia del ruido de sensor, mientras que $W_{e}$ puede usarse para reflejar requisitos sobre la forma de ciertas funciones de transferencia en lazo cerrado (por ejemplo, la forma de la función de sensibilidad de salida). De forma similar, $W_{u}$ puede usarse para reflejar restricciones sobre señales de control o de actuador, y el precompensador punteado $W_{r}$ es un elemento opcional usado para imponer una forma deliberada al comando o para representar un sistema con realimentación no unitaria en forma equivalente de realimentación unitaria.

De hecho, es esencial usar matrices de ponderación apropiadas para poder utilizar la teoría de control óptimo discutida en este libro (es decir, teoría $\mathcal{H}_{2}$ y $\mathcal{H}_{\infty}$). Por ello, un paso muy importante en el proceso de diseño del controlador es elegir adecuadamente los pesos $W_{e}, W_{d}, W_{u}$ y posiblemente $W_{n}, W_{i}, W_{r}$. La elección adecuada de pesos para un problema práctico particular no es trivial. En muchas ocasiones, como en el caso escalar, los pesos se eligen puramente como parámetro de diseño sin base física directa; por tanto, pueden tratarse como parámetros de ajuste escogidos por el diseñador para lograr el mejor compromiso entre objetivos en conflicto. La selección de las matrices de ponderación debe guiarse por las entradas esperadas del sistema y la importancia relativa de las salidas.

Así, el diseño de control puede verse como el proceso de elegir un controlador $K$ tal que ciertas señales ponderadas se hagan pequeñas en algún sentido. Existen muchas formas de definir la pequeñez de una señal o de una matriz de transferencia, como discutimos en el capítulo anterior. Distintas definiciones conducen a distintos métodos de síntesis de control, y algunos son mucho más difíciles que otros. Un ingeniero de control debe juzgar la complejidad matemática frente a los requerimientos de ingeniería.

A continuación introducimos dos clases de formulaciones de desempeño: criterios $\mathcal{H}_{2}$ y $\mathcal{H}_{\infty}$. Para simplificar la presentación, supondremos que $d_{i}=0$ y $n=0$.

\subsection*{Desempeño $\mathcal{H}_{2}$}
Suponga, por ejemplo, que la perturbación $\tilde{d}$ puede modelarse aproximadamente como un impulso con dirección de entrada aleatoria; es decir,

$$
\tilde{d}(t)=\eta \delta(t)
$$

y

$$
E\left(\eta \eta^{*}\right)=I
$$

donde $E$ denota la esperanza. Podemos elegir minimizar la energía esperada del error $e$ debida a la perturbación $\tilde{d}$:

$$
E\left\{\|e\|_{2}^{2}\right\}=E\left\{\int_{0}^{\infty}\|e\|^{2} d t\right\}=\left\|W_{e} S_{o} W_{d}\right\|_{2}^{2}
$$

En general, un controlador que minimice solo el criterio anterior puede llevar a una señal de control $u$ muy grande, lo que podría causar saturación de actuadores y muchos otros problemas indeseables. Por ello, para un diseño realista, es necesario incluir la señal de control $u$ en la función de costo. Así, nuestro criterio de diseño usualmente sería algo como

$$
E\left\{\|e\|_{2}^{2}+\rho^{2}\|\tilde{u}\|_{2}^{2}\right\}=\left\|\left[\begin{array}{c}
W_{e} S_{o} W_{d} \\
\rho W_{u} K S_{o} W_{d}
\end{array}\right]\right\|_{2}^{2}
$$

con alguna elección apropiada de matriz de ponderación $W_{u}$ y escalar $\rho$. El parámetro $\rho$ define claramente el compromiso discutido antes entre buen rechazo de perturbaciones en la salida y esfuerzo de control (o rechazo de perturbación y ruido de sensor en los actuadores). Nótese que puede fijarse $\rho=1$ mediante una elección apropiada de $W_{u}$. Este problema puede verse como minimizar la energía consumida por el sistema para rechazar la perturbación $d$.

Este tipo de problema fue el paradigma dominante en las décadas de 1960 y 1970 y suele denominarse control lineal cuadrático gaussiano, o simplemente LQG. (También nos referiremos a estos problemas como problemas de sensibilidad mixta $\mathcal{H}_{2}$ para mantener consistencia con los problemas $\mathcal{H}_{\infty}$ discutidos después.) El desarrollo de este paradigma estimuló amplios esfuerzos de investigación y es responsable de importantes innovaciones tecnológicas, particularmente en el área de estimación. Las contribuciones teóricas incluyen una comprensión más profunda de sistemas lineales y mejores métodos computacionales para sistemas complejos mediante técnicas en espacio de estados. La principal limitación de esta teoría es la ausencia de un tratamiento formal de la incertidumbre en la planta misma. Al permitir solo ruido aditivo como incertidumbre, la teoría estocástica ignoró este importante asunto práctico. La incertidumbre de planta es particularmente crítica en sistemas con realimentación. (Véase Paganini $[1995,1996]$ para algunos resultados recientes sobre teoría de control robusto $\mathcal{H}_{2}$.)

\subsection*{Desempeño $\mathcal{H}_{\infty}$}
Aunque la norma $\mathcal{H}_{2}$ (o norma $\mathcal{L}_{2}$) puede ser una medida de desempeño significativa y aunque la teoría LQG puede brindar compromisos de diseño eficientes bajo ciertas hipótesis de perturbación y planta, la norma $\mathcal{H}_{2}$ sufre una deficiencia mayor. Esta deficiencia se debe a que el compromiso entre reducción de error por perturbación y reducción de error por ruido de sensor no es la única restricción del diseño por realimentación. El problema es que estos compromisos de desempeño suelen quedar eclipsados por una segunda limitación sobre ganancias altas de lazo: el requisito de tolerancia a incertidumbres. Aunque un controlador pueda diseñarse usando modelos FDLTI, el diseño debe implementarse y operar con una planta física real. Las propiedades de sistemas físicos (en particular, las formas en que se desvían de modelos lineales de dimensión finita) imponen limitaciones estrictas sobre el rango de frecuencias en el cual las ganancias de lazo pueden ser grandes.

Una solución a este problema sería imponer restricciones explícitas a la ganancia de lazo en la función de costo. Por ejemplo, puede elegirse minimizar

$$
\sup _{\|\tilde{d}\|_{2} \leq 1}\|e\|_{2}=\left\|W_{e} S_{o} W_{d}\right\|_{\infty}
$$

sujeto a algunas restricciones sobre energía de control o ancho de banda de control:

$$
\sup _{\|\tilde{d}\|_{2} \leq 1}\|\tilde{u}\|_{2}=\left\|W_{u} K S_{o} W_{d}\right\|_{\infty}
$$

O, con mayor frecuencia, puede introducirse un parámetro $\rho$ y un criterio mixto

$$
\sup _{\|\tilde{d}\|_{2} \leq 1}\left\{\|e\|_{2}^{2}+\rho^{2}\|\tilde{u}\|_{2}^{2}\right\}=\left\|\left[\begin{array}{c}
W_{e} S_{o} W_{d} \\
\rho W_{u} K S_{o} W_{d}
\end{array}\right]\right\|_{\infty}^{2}
$$

Alternativamente, si el margen de estabilidad robusta del sistema es la principal preocupación, debe limitarse la sensibilidad complementaria ponderada. Así, toda la función de costo puede ser

$$
\left\|\left[\begin{array}{c}
W_{e} S_{o} W_{d} \\
\rho W_{1} T_{o} W_{2}
\end{array}\right]\right\|_{\infty}
$$

donde $W_{1}$ y $W_{2}$ son matrices de escalamiento de incertidumbre dependientes de la frecuencia. Estos problemas de diseño suelen llamarse problemas de sensibilidad mixta $\mathcal{H}_{\infty}$. Para un sistema escalar, un problema de minimización en norma $\mathcal{H}_{\infty}$ también puede verse como minimizar la magnitud máxima de la respuesta en régimen permanente del sistema con respecto a entradas sinusoidales de peor caso.

\section{Selección de funciones de ponderación}
La selección de funciones de ponderación para un problema de diseño específico suele involucrar ajustes ad hoc, muchas iteraciones y afinación fina. Es muy difícil dar una fórmula general de ponderaciones que funcione en todos los casos. No obstante, intentaremos dar algunas guías en esta sección observando un problema SISO típico.

Considere un sistema SISO con realimentación como en la Figura 6.1. Entonces el error de seguimiento es $e=r-y=S(r-d)+T n-S P d_{i}$. Por tanto, como hemos discutido antes, debemos mantener $|S|$ pequeño sobre un rango de frecuencias, típicamente bajas frecuencias donde $r$ y $d$ son significativas. Para motivar la elección de la ponderación de desempeño $W_{e}$, sea $L=P K$ un sistema estándar de segundo orden

$$
L=\frac{\omega_{n}^{2}}{s\left(s+2 \xi \omega_{n}\right)}
$$

Es bien conocido en control clásico que la calidad de la respuesta temporal (al escalón) puede cuantificarse por tiempo de subida $t_{r}$, tiempo de establecimiento $t_{s}$ y sobreimpulso porcentual $100 M_{p} \%$. Además, estos índices de desempeño pueden calcularse aproximadamente como

$$
t_{r} \approx \frac{0.6+2.16 \xi}{\omega_{n}}, 0.3 \leq \xi \leq 0.8 ; t_{s} \approx \frac{4}{\xi \omega_{n}} ; M_{p}=e^{-\frac{\pi \xi}{\sqrt{1-\xi^{2}}}}, 0<\xi<1
$$

Los puntos clave a notar son que (1) la rapidez de respuesta del sistema es proporcional a $\omega_{n}$ y (2) el sobreimpulso de la respuesta del sistema está determinado solo por la razón de amortiguamiento $\xi$. Es bien conocido que la frecuencia $\omega_{n}$ y la razón de amortiguamiento $\xi$ pueden capturarse esencialmente en el dominio de la frecuencia mediante la frecuencia de cruce de lazo abierto y el margen de fase, o bien mediante el ancho de banda y el pico resonante de la función de sensibilidad complementaria en lazo cerrado $T$.

Dado que nuestros objetivos de desempeño están estrechamente relacionados con la función de sensibilidad, consideraremos con cierto detalle cómo estos índices temporales o, equivalentemente, $\omega_{n}$ y $\xi$, se relacionan con la respuesta en frecuencia de la función de sensibilidad

$$
S=\frac{1}{1+L}=\frac{s\left(s+2 \xi \omega_{n}\right)}{s^{2}+2 \xi \omega_{n} s+\omega_{n}^{2}}
$$

\insertimage[\label{img:funcion-sensibilidad-s-xi}]{2024_06_29_ce7c73b22613114bd4d6g-104_761_1147_411_478}{width=0.35\textwidth}{Función de sensibilidad $S$ para $\xi=0.05,0.1,0.2,0.5,0.8$, y 1 con frecuencia normalizada $\left(\omega / \omega_{n}\right)$.}

La respuesta en frecuencia de la función de sensibilidad $S$ se muestra en la Figura 6.4. Nótese que $\left|S\left(j \omega_{n} / \sqrt{2}\right)\right|=1$. Podemos considerar el ancho de banda en lazo cerrado como $\omega_{b} \approx \omega_{n} / \sqrt{2}$, ya que más allá de esa frecuencia el sistema en lazo cerrado no podrá seguir la referencia y la perturbación en realidad será amplificada.

Luego, nótese que

$$
M_{s}:=\|S\|_{\infty}=\left|S\left(j \omega_{\max }\right)\right|=\frac{\alpha \sqrt{\alpha^{2}+4 \xi^{2}}}{\sqrt{\left(1-\alpha^{2}\right)^{2}+4 \xi^{2} \alpha^{2}}}
$$

donde $\alpha=\sqrt{0.5+0.5 \sqrt{1+8 \xi^{2}}}$ y $\omega_{\max }=\alpha \omega_{n}$. Por ejemplo, $M_{s}=5.123$ cuando $\xi=0.1$. La relación entre $\xi$ y $M_{s}$ se muestra en la Figura 6.5. Es claro que el sobreimpulso puede ser excesivo si $M_{s}$ es grande. Por lo tanto, un buen diseño de control no debería tener un $M_{s}$ muy grande.

Ahora suponga que se nos dan especificaciones de desempeño en dominio del tiempo; entonces podemos determinar los requisitos correspondientes en dominio de la frecuencia en términos del ancho de banda $\omega_{b}$ y la sensibilidad pico $M_{s}$. Por lo tanto, un buen diseño de control debe producir una función de sensibilidad $S$ que satisfaga simultáneamente los requisitos de $\omega_{b}$ y $M_{s}$, como se muestra en la Figura 6.6. Estos requisitos pueden representarse aproximadamente como

$$
|S(s)| \leq\left|\frac{s}{s / M_{s}+\omega_{b}}\right|, s=j \omega, \forall \omega
$$

\insertimage[\label{img:sensibilidad-pico-vs-amortiguamiento}]{2024_06_29_ce7c73b22613114bd4d6g-105_634_1133_545_491}{width=0.35\textwidth}{Sensibilidad pico $M_{s}$ versus razón de amortiguamiento $\xi$.}

\insertimage[\label{img:ponderacion-desempeno-we-s-deseada}]{2024_06_29_ce7c73b22613114bd4d6g-105_626_1155_1389_474}{width=0.35\textwidth}{Ponderación de desempeño $W_{e}$ y $S$ deseada.}

O, equivalentemente, $\left|W_{e} S\right| \leq 1$ con


\begin{equation*}
W_{e}=\frac{s / M_{s}+\omega_{b}}{s} \tag{6.5}
\end{equation*}


La discusión anterior aplica en principio a la mayoría de diseños de control y, por tanto, la ponderación anterior puede usarse como candidata en un diseño inicial. Dado que el error en estado estacionario respecto a una entrada escalón está dado por $|S(0)|$, es claro que $|S(0)|=0$ si el sistema en lazo cerrado es estable y $\left\|W_{e} S\right\|_{\infty}<\infty$.

Desafortunadamente, las técnicas de control óptimo descritas en este libro no pueden usarse directamente para problemas con tales ponderaciones, ya que esas técnicas suponen que todos los polos inestables del sistema (incluyendo planta y todas las ponderaciones de desempeño y control) son estabilizables por el control y detectables desde las salidas de medición, lo cual claramente no se satisface si $W_{e}$ tiene un polo sobre el eje imaginario, pues $W_{e}$ no es detectable desde la medición. Discutiremos en el Capítulo 14 cómo reformular estos problemas para que las técnicas descritas en este libro puedan aplicarse. Existe una teoría que trata directamente estos problemas, pero es mucho más complicada, tanto teórica como computacionalmente, y no parece ofrecer mucha ventaja.

\insertimage[\label{img:ponderacion-practica-we-s-deseada}]{2024_06_29_ce7c73b22613114bd4d6g-106_626_1152_1178_476}{width=0.35\textwidth}{Ponderación práctica de desempeño $W_{e}$ y $S$ deseada.}

Ahora, en lugar de seguimiento perfecto para entrada escalón, suponga que solo necesitamos que el error en estado estacionario respecto a escalón no sea mayor que $\epsilon$ (es decir, $|S(0)| \leq \epsilon$); entonces basta elegir una ponderación $W_{e}$ que satisfaga $\left|W_{e}(0)\right| \geq 1 / \varepsilon$ de modo que pueda lograrse $\left\|W_{e} S\right\|_{\infty} \leq 1$. Una elección posible de $W_{e}$ puede obtenerse modificando la ponderación de la ecuación (6.5):


\begin{equation*}
W_{e}=\frac{s / M_{s}+\omega_{b}}{s+\omega_{b} \varepsilon} \tag{6.6}
\end{equation*}


Por tanto, para fines prácticos, usualmente puede elegirse un $\varepsilon$ apropiado, como se muestra en la Figura 6.7, para satisfacer las especificaciones de desempeño. Si se desea una transición más abrupta entre baja y alta frecuencia, la ponderación $W_{e}$ puede modificarse como sigue:


\begin{equation*}
W_{e}=\left(\frac{s / \sqrt[k]{M_{s}}+\omega_{b}}{s+\omega_{b} \sqrt[k]{\varepsilon}}\right)^{k} \tag{6.7}
\end{equation*}


para algún entero $k \geq 1$.

La selección de la ponderación de control $W_{u}$ se obtiene de forma similar a partir de la discusión anterior considerando la ecuación de señal de control

$$
u=K S(r-n-d)-T d_{i}
$$

La magnitud de $|K S|$ en el rango de baja frecuencia está esencialmente limitada por el costo permitido de esfuerzo de control y por el límite de saturación de los actuadores; por tanto, en general, la ganancia máxima $M_{u}$ de $K S$ puede ser bastante grande, mientras que la ganancia en alta frecuencia está esencialmente limitada por el ancho de banda del controlador $\left(\omega_{b c}\right)$ y por las frecuencias de ruido (de sensor). Idealmente, se desea atenuar tan rápido como sea posible más allá del ancho de banda de control deseado para que el ruido de alta frecuencia se atenúe lo más posible. Así, una ponderación candidata $W_{u}$ sería


\begin{equation*}
W_{u}=\frac{s+\omega_{b c} / M_{u}}{\omega_{b c}} \tag{6.8}
\end{equation*}


\insertimage[\label{img:ponderacion-control-wu-ks-deseada}]{2024_06_29_ce7c73b22613114bd4d6g-107_646_1152_1390_476}{width=0.35\textwidth}{Ponderación de control $W_{u}$ y $K S$ deseada.}

Sin embargo, nuevamente, las técnicas de diseño óptimo desarrolladas en este libro no pueden aplicarse directamente a un problema con una ponderación de control impropia. Por ello, introduciremos un polo alejado para hacer $W_{u}$ propia:


\begin{equation*}
W_{u}=\frac{s+\omega_{b c} / M_{u}}{\varepsilon_{1} s+\omega_{b c}} \tag{6.9}
\end{equation*}


para un $\varepsilon_{1}>0$ pequeño, como se muestra en la Figura 6.8. De forma similar, si se desea una atenuación más rápida, puede elegirse


\begin{equation*}
W_{u}=\left(\frac{s+\omega_{b c} / \sqrt[k]{M_{u}}}{\sqrt[k]{\varepsilon_{1}} s+\omega_{b c}}\right)^{k} \tag{6.10}
\end{equation*}


para algún entero $k \geq 1$.

Las ponderaciones para problemas MIMO pueden elegirse inicialmente como matrices diagonales, con cada término diagonal escogido en la forma anterior.

\section{Relación de ganancia y fase de Bode}
Un problema importante que surge frecuentemente es el nivel de desempeño que puede lograrse en diseño por realimentación. En la Sección 6.1 se mostró que los objetivos de diseño por realimentación son inherentemente conflictivos y que debe hacerse un compromiso entre distintos objetivos. También se sabe que requisitos fundamentales, como estabilidad y robustez, imponen limitaciones inherentes a las propiedades de realimentación independientemente del método de diseño, y esas limitaciones se vuelven más severas en presencia de ceros y polos en semiplano derecho en la función de transferencia de lazo abierto.

En teoría clásica de realimentación, la relación integral ganancia-fase de Bode (véase Bode [1945]) se ha usado como herramienta importante para expresar restricciones de diseño en sistemas escalares. Esta relación integral dice que la fase de una función de transferencia estable y de fase mínima está determinada de manera única por su magnitud. Más precisamente, sea $L(s)$ una función de transferencia estable y de fase mínima; entonces


\begin{equation*}
\angle L\left(j \omega_{0}\right)=\frac{1}{\pi} \int_{-\infty}^{\infty} \frac{d \ln |L|}{d \nu} \ln \operatorname{coth} \frac{|\nu|}{2} d \nu \tag{6.11}
\end{equation*}


donde $\nu:=\ln \left(\omega / \omega_{0}\right)$. La función $\ln \operatorname{coth} \frac{|\nu|}{2}=\ln \frac{e^{|\nu| / 2}+e^{-|\nu| / 2}}{e^{|\nu| / 2}-e^{-|\nu| / 2}}$ está graficada en la Figura 6.9.

Nótese que $\ln \operatorname{coth} \frac{|\nu|}{2}$ decrece rápidamente conforme $\omega$ se aleja de $\omega_{0}$ y, por lo tanto, la integral depende principalmente del comportamiento de $\frac{d \ln |L(j \omega)|}{d \nu}$ cerca de la frecuencia $\omega_{0}$. Esto queda claro de la siguiente integración:

$$
\frac{1}{\pi} \int_{-\alpha}^{\alpha} \ln \operatorname{coth} \frac{|\nu|}{2} d \nu=\left\{\begin{array}{ll}
1.1406(\mathrm{rad}), & \alpha=\ln 3 \\
1.3146(\mathrm{rad}), & \alpha=\ln 5 \\
1.443(\mathrm{rad}), & \alpha=\ln 10
\end{array}= \begin{cases}65.3^{\circ}, & \alpha=\ln 3 \\
75.3^{\circ}, & \alpha=\ln 5 \\
82.7^{\circ}, & \alpha=\ln 10\end{cases}\right.
$$

\insertimage[\label{img:funcion-log-coth}]{2024_06_29_ce7c73b22613114bd4d6g-109_762_895_405_607}{width=0.35\textwidth}{Función $\ln \operatorname{coth}\frac{|\nu|}{2}$ usada en la relación ganancia-fase de Bode.}

Nótese que $\frac{d \ln |L(j \omega)|}{d \nu}$ es la pendiente del diagrama de Bode, generalmente negativa para casi todas las frecuencias. Se sigue que $\angle L\left(j \omega_{0}\right)$ será grande si la ganancia $L$ se atenúa lentamente cerca de $\omega_{0}$, y pequeña si se atenúa rápidamente cerca de $\omega_{0}$. Por ejemplo, suponga que la pendiente $\frac{d \ln |L(j \omega)|}{d \nu}=-\ell$, es decir, $\left(-20 \ell \mathrm{dB}\right.$ por década), en la vecindad de $\omega_{0}$; entonces es razonable esperar

$$
\angle L\left(j \omega_{0}\right)< \begin{cases}-\ell \times 65.3^{\circ}, & \text { si la pendiente de } L=-\ell \text { para } \frac{1}{3} \leq \frac{\omega}{\omega_{0}} \leq 3 \\ -\ell \times 75.3^{\circ}, & \text { si la pendiente de } L=-\ell \text { para } \frac{1}{5} \leq \frac{\omega}{\omega_{0}} \leq 5 \\ -\ell \times 82.7^{\circ}, & \text { si la pendiente de } L=-\ell \text { para } \frac{1}{10} \leq \frac{\omega}{\omega_{0}} \leq 10 .\end{cases}
$$

El comportamiento de $\angle L(j \omega)$ es particularmente importante cerca de la frecuencia de cruce $\omega_{c}$, donde $\left|L\left(j \omega_{c}\right)\right|=1$, ya que $\pi+\angle L\left(j \omega_{c}\right)$ es el margen de fase del sistema con realimentación. Además, la diferencia de retorno está dada por

$$
\left|1+L\left(j \omega_{c}\right)\right|=\left|1+L^{-1}\left(j \omega_{c}\right)\right|=2\left|\sin \frac{\pi+\angle L\left(j \omega_{c}\right)}{2}\right|,
$$

la cual no debe ser demasiado pequeña para tener buena robustez de estabilidad. Si $\pi+\angle L\left(j \omega_{c}\right)$ se fuerza a ser muy pequeño por una atenuación rápida de ganancia, el sistema con realimentación amplificará perturbaciones y mostrará poca tolerancia a incertidumbre en y cerca de $\omega_{c}$. Dado que generalmente se requiere que la función de transferencia de lazo $L$ caiga lo más rápido posible en alta frecuencia, es razonable esperar que $\angle L\left(j \omega_{c}\right)$ sea como máximo $-\ell \times 90^{\circ}$ si la pendiente de $L(j \omega)$ es $-\ell$ cerca de $\omega_{c}$. Así, es importante mantener la pendiente de $L$ cerca de $\omega_{c}$ no mucho menor que -1 en un rango razonablemente amplio de frecuencias para garantizar un desempeño razonable. El conflicto entre tasa de atenuación y calidad de lazo cerca del cruce es, por tanto, claramente evidente.

La relación ganancia-fase de Bode puede extenderse fácilmente a funciones de transferencia estables y de fase no mínima. Sean $z_{1}, z_{2}, \ldots, z_{k}$ los ceros en semiplano derecho de $L(s)$, entonces $L$ puede factorizarse como

$$
L(s)=\frac{-s+z_{1}}{s+z_{1}} \frac{-s+z_{2}}{s+z_{2}} \cdots \frac{-s+z_{k}}{s+z_{k}} L_{\mathrm{mp}}(s)
$$

donde $L_{\mathrm{mp}}$ es estable y de fase mínima, y $|L(j \omega)|=\left|L_{\mathrm{mp}}(j \omega)\right|$. Por lo tanto,

$$
\begin{gathered}
\angle L\left(j \omega_{0}\right)=\angle L_{\mathrm{mp}}\left(j \omega_{0}\right)+\angle \prod_{i=1}^{k} \frac{-j \omega_{0}+z_{i}}{j \omega_{0}+z_{i}} \\
=\frac{1}{\pi} \int_{-\infty}^{\infty} \frac{d \ln \left|L_{\mathrm{mp}}\right|}{d \nu} \ln \operatorname{coth} \frac{|\nu|}{2} d \nu+\sum_{i=1}^{k} \angle \frac{-j \omega_{0}+z_{i}}{j \omega_{0}+z_{i}}
\end{gathered}
$$

lo cual da


\begin{equation*}
\angle L\left(j \omega_{0}\right)=\frac{1}{\pi} \int_{-\infty}^{\infty} \frac{d \ln |L|}{d \nu} \ln \operatorname{coth} \frac{|\nu|}{2} d \nu+\sum_{i=1}^{k} \angle \frac{-j \omega_{0}+z_{i}}{j \omega_{0}+z_{i}} \tag{6.12}
\end{equation*}


Como $\angle \frac{-j \omega_{0}+z_{i}}{j \omega_{0}+z_{i}} \leq 0$ para cada $i$, un cero de fase no mínima aporta retraso de fase adicional e impone limitaciones sobre la tasa de caída de la ganancia en lazo abierto. Por ejemplo, suponga que $L$ tiene un cero en $z>0$; entonces

$$
\phi_{1}\left(\omega_{0} / z\right):=\left.\angle \frac{-j \omega_{0}+z}{j \omega_{0}+z}\right|_{\omega_{0}=z, z / 2, z / 4}=-90^{\circ},-53.13^{\circ},-28^{\circ}
$$

como se muestra en la Figura 6.10. Dado que la pendiente de $|L|$ cerca de la frecuencia de cruce es, en general, no mayor que -1, lo que significa que la fase debida a la parte de fase mínima $L_{\mathrm{mp}}$ de $L$ será, en general, no mayor que $-90^{\circ}$, la frecuencia de cruce (o ancho de banda en lazo cerrado) debe satisfacer


\begin{equation*}
\omega_{c}<z / 2 \tag{6.13}
\end{equation*}


para garantizar la estabilidad en lazo cerrado y un desempeño razonable.

Ahora suponga que $L$ tiene un par de ceros complejos en semiplano derecho en $z=x \pm j y$ con $x>0$; entonces

$$
\begin{aligned}
& \phi_2\left(\omega_0 /|z|\right):=\left.\angle \frac{-j \omega_0+z}{j \omega_0+z} \frac{-j \omega_0+\bar{z}}{j \omega_0+\bar{z}}\right|_{\omega_0=|z|,|z| / 2,|z| / 3,|z| / 4} \\
& \approx\left\{\begin{array}{rrrrl}
-180^{\circ}, & -106.26^{\circ}, & -73.7^{\circ}, & -56^{\circ}, & \operatorname{Re}(z) \gg \Im(z) \\
-180^{\circ}, & -86.7^{\circ}, & -55.9^{\circ}, & -41.3^{\circ}, & \operatorname{Re}(z) \approx \Im(z) \\
-360^{\circ}, & 0^{\circ}, & 0^{\circ}, & 0^{\circ}, & \operatorname{Re}(z) \ll \Im(z)
\end{array}\right.
\end{aligned}
$$

\insertimage[\label{img:fase-cero-real-panel-2}]{2024_06_29_ce7c73b22613114bd4d6g-111_721_811_453_646}{width=0.35\textwidth}{Fase $\phi_{1}\left(\omega_{0} / z\right)$ debida a un cero real $z>0$.}

\insertimage[\label{img:fase-cero-complejo}]{2024_06_29_ce7c73b22613114bd4d6g-111_726_811_1385_646}{width=0.35\textwidth}{Fase $\phi_{2}\left(\omega_{0} /|z|\right)$ debida a un par de ceros complejos: $z=x \pm j y$ y $x>0$.}

como se muestra en la Figura 6.11. En este caso concluimos que la frecuencia de cruce debe satisfacer

\[
\omega_{c}< \begin{cases}|z| / 4, & \operatorname{Re}(z) \gg \Im(z)  \tag{6.14}\\ |z| / 3, & \operatorname{Re}(z) \approx \Im(z) \\ |z|, & \operatorname{Re}(z) \ll \Im(z)\end{cases}
\]

para garantizar estabilidad en lazo cerrado y desempeño razonable.

\section{Integral de sensibilidad de Bode}
En esta sección, consideramos las limitaciones de diseño impuestas por restricciones de ancho de banda y por polos y ceros en semiplano derecho usando la integral de sensibilidad de Bode y la integral de Poisson. Sea $L$ la función de transferencia de lazo abierto con al menos dos polos más que ceros, y sean $p_{1}, p_{2}, \ldots, p_{m}$ los polos en semiplano derecho abierto de $L$. Entonces se cumple la siguiente integral de sensibilidad de Bode:


\begin{equation*}
\int_{0}^{\infty} \ln |S(j \omega)| d \omega=\pi \sum_{i=1}^{m} \operatorname{Re}\left(p_{i}\right) \tag{6.15}
\end{equation*}


En el caso en que $L$ sea estable, la integral se simplifica a


\begin{equation*}
\int_{0}^{\infty} \ln |S(j \omega)| d \omega=0 \tag{6.16}
\end{equation*}


Estas integrales muestran que existirá un rango de frecuencias sobre el cual la magnitud de la función de sensibilidad excede uno si se la mantiene por debajo de uno en otras frecuencias, como se ilustra en la Figura 6.12. Este es el llamado efecto colchón de agua (water-bed effect).

\insertimage[\label{img:efecto-colchon-agua-sensibilidad}]{2024_06_29_ce7c73b22613114bd4d6g-112_518_897_1633_603}{width=0.35\textwidth}{Efecto colchón de agua de la función de sensibilidad.}

Suponga que el sistema con realimentación está diseñado de modo que el nivel de reducción de sensibilidad viene dado por

$$
|S(j \omega)| \leq \epsilon<1, \quad \forall \omega \in\left[0, \omega_{l}\right]
$$

donde $\epsilon>0$ es una constante dada.

Las restricciones de ancho de banda en diseño por realimentación típicamente requieren que la función de transferencia de lazo abierto sea pequeña por encima de cierta frecuencia especificada y que decrezca con una tasa equivalente a más de un exceso polo-cero por encima de esa frecuencia. Estas restricciones suelen ser necesarias para garantizar robustez de estabilidad pese a incertidumbre de modelado en la planta, particularmente en alta frecuencia. Una forma de cuantificar tales restricciones de ancho de banda es exigir que la función de transferencia de lazo abierto satisfaga

$$
|L(j \omega)| \leq \frac{M_{h}}{\omega^{1+\beta}} \leq \tilde{\epsilon}<1, \quad \forall \omega \in\left[\omega_{h}, \infty\right)
$$

donde $\omega_{h}>\omega_{l}$, y $M_{h}>0, \beta>0$ son constantes dadas.

Nótese que para $\omega \geq \omega_{h}$,

$$
|S(j \omega)| \leq \frac{1}{1-|L(j \omega)|} \leq \frac{1}{1-\frac{M_{h}}{\omega^{1+\beta}}}
$$

y

$$
\begin{aligned}
-\int_{\omega_{h}}^{\infty} \ln \left(1-\frac{M_{h}}{\omega^{1+\beta}}\right) d \omega & =\sum_{i=1}^{\infty} \int_{\omega_{h}}^{\infty} \frac{1}{i}\left(\frac{M_{h}}{\omega^{1+\beta}}\right)^{i} d \omega \\
& =\sum_{i=1}^{\infty} \frac{1}{i} \frac{\omega_{h}}{i(1+\beta)-1}\left(\frac{M_{h}}{\omega_{h}^{1+\beta}}\right)^{i} \\
& \leq \frac{\omega_{h}}{\beta} \sum_{i=1}^{\infty} \frac{1}{i}\left(\frac{M_{h}}{\omega_{h}^{1+\beta}}\right)^{i}=-\frac{\omega_{h}}{\beta} \ln \left(1-\frac{M_{h}}{\omega_{h}^{1+\beta}}\right) \\
& \leq-\frac{\omega_{h}}{\beta} \ln (1-\tilde{\epsilon})
\end{aligned}
$$

Entonces

$$
\begin{gathered}
\pi \sum_{i=1}^{m} \operatorname{Re}\left(p_{i}\right)=\int_{0}^{\infty} \ln |S(j \omega)| d \omega \\
=\int_{0}^{\omega_{l}} \ln |S(j \omega)| d \omega+\int_{\omega_{l}}^{\omega_{h}} \ln |S(j \omega)| d \omega+\int_{\omega_{h}}^{\infty} \ln |S(j \omega)| d \omega \\
\leq \omega_{l} \ln \epsilon+\left(\omega_{h}-\omega_{l}\right) \max _{\omega \in\left[\omega_{l}, \omega_{h}\right]} \ln |S(j \omega)|-\int_{\omega_{h}}^{\infty} \ln \left(1-\frac{M_{h}}{\omega^{1+\beta}}\right) d \omega \\
\leq \omega_{l} \ln \epsilon+\left(\omega_{h}-\omega_{l}\right) \max _{\omega \in\left[\omega_{l}, \omega_{h}\right]} \ln |S(j \omega)|-\frac{\omega_{h}}{\beta} \ln (1-\tilde{\epsilon})
\end{gathered}
$$

lo cual da

$$
\max _{\omega \in\left[\omega_{l}, \omega_{h}\right]}|S(j \omega)| \geq e^{\alpha}\left(\frac{1}{\epsilon}\right)^{\frac{\omega_{l}}{\omega_{h}-\omega_{l}}}(1-\tilde{\epsilon})^{\frac{\omega_{h}}{\beta\left(\omega_{h}-\omega_{l}\right)}}
$$

donde

$$
\alpha=\frac{\pi \sum_{i=1}^{m} \operatorname{Re}\left(p_{i}\right)}{\omega_{h}-\omega_{l}}
$$

La cota inferior anterior muestra que la sensibilidad puede ser muy significativa en la banda de transición.

A continuación, usando la relación integral de Poisson, investigamos las restricciones de diseño sobre propiedades de sensibilidad impuestas por ceros de fase no mínima del lazo abierto. Suponga que $L$ tiene al menos un polo más que ceros y suponga que $z=x_{0}+j y_{0}$ con $x_{0}>0$ es un cero en semiplano derecho de $L$. Entonces


\begin{equation*}
\int_{-\infty}^{\infty} \ln |S(j \omega)| \frac{x_{0}}{x_{0}^{2}+\left(\omega-y_{0}\right)^{2}} d \omega=\pi \ln \prod_{i=1}^{m}\left|\frac{z+p_{i}}{z-p_{i}}\right| \tag{6.17}
\end{equation*}


Esta integral implica que la capacidad de reducción de sensibilidad del sistema puede estar severamente limitada por polos inestables de lazo abierto y ceros de fase no mínima, especialmente cuando estos polos y ceros están cercanos entre sí.

Defina

$$
\theta(z):=\int_{-\omega_{l}}^{\omega_{l}} \frac{x_{0}}{x_{0}^{2}+\left(\omega-y_{0}\right)^{2}} d \omega
$$

Entonces

$$
\begin{gathered}
\pi \ln \prod_{i=1}^{m}\left|\frac{z+p_{i}}{z-p_{i}}\right|=\int_{-\infty}^{\infty} \ln |S(j \omega)| \frac{x_{0}}{x_{0}^{2}+\left(\omega-y_{0}\right)^{2}} d \omega \\
\leq(\pi-\theta(z)) \ln \|S(j \omega)\|_{\infty}+\theta(z) \ln (\epsilon)
\end{gathered}
$$

lo cual da

$$
\|S(s)\|_{\infty} \geq\left(\frac{1}{\epsilon}\right)^{\frac{\theta(z)}{\pi-\theta(z)}}\left(\prod_{i=1}^{m}\left|\frac{z+p_{i}}{z-p_{i}}\right|\right)^{\frac{\pi}{\pi-\theta(z)}}
$$

Esta cota inferior sobre la sensibilidad máxima muestra que, para un sistema de fase no mínima, su sensibilidad debe incrementarse significativamente por encima de uno en ciertas frecuencias si se desea reducción de sensibilidad en otras frecuencias.

\section{Restricciones de analitcidad}
Sean $p_{1}, p_{2}, \ldots, p_{m}$ y $z_{1}, z_{2}, \ldots, z_{k}$ los polos y ceros en semiplano derecho abierto de $L$, respectivamente. Suponga que el sistema en lazo cerrado es estable. Entonces

$$
S\left(p_{i}\right)=0, \quad T\left(p_{i}\right)=1, i=1,2, \ldots, m
$$

y

$$
S\left(z_{j}\right)=1, \quad T\left(z_{j}\right)=0, j=1,2, \ldots, k
$$

La estabilidad interna del sistema con realimentación se garantiza al satisfacer estas condiciones de analiticidad (o interpolación). Por otro lado, estas condiciones también imponen limitaciones severas sobre el desempeño alcanzable del sistema con realimentación.

Suponga que $S=(I+L)^{-1}$ y $T=L(I+L)^{-1}$ son estables. Entonces $p_{1}, p_{2}, \ldots, p_{m}$ son los ceros en semiplano derecho de $S$ y $z_{1}, z_{2}, \ldots, z_{k}$ son los ceros en semiplano derecho de $T$. Sea

$$
B_{p}(s)=\prod_{i=1}^{m} \frac{s-p_{i}}{s+p_{i}}, \quad B_{z}(s)=\prod_{j=1}^{k} \frac{s-z_{j}}{s+z_{j}}
$$

Entonces $\left|B_{p}(j \omega)\right|=1$ y $\left|B_{z}(j \omega)\right|=1$ para todas las frecuencias y, además,

$$
B_{p}^{-1}(s) S(s) \in \mathcal{H}_{\infty}, \quad B_{z}^{-1}(s) T(s) \in \mathcal{H}_{\infty}
$$

Por lo tanto, por el teorema del módulo máximo,

$$
\|S(s)\|_{\infty}=\left\|B_{p}^{-1}(s) S(s)\right\|_{\infty} \geq\left|B_{p}^{-1}(z) S(z)\right|
$$

para cualquier $z$ con $\operatorname{Re}(z)>0$. Sea $z$ un cero en semiplano derecho de $L$; entonces

$$
\|S(s)\|_{\infty} \geq\left|B_{p}^{-1}(z)\right|=\prod_{i=1}^{m}\left|\frac{z+p_{i}}{z-p_{i}}\right|
$$

De forma similar, puede obtenerse

$$
\|T(s)\|_{\infty} \geq\left|B_{z}^{-1}(p)\right|=\prod_{j=1}^{k}\left|\frac{p+z_{j}}{p-z_{j}}\right|
$$

donde $p$ es un polo en semiplano derecho de $L$.

El problema ponderado puede tratarse de manera análoga. Sea $W_{e}$ un peso tal que $W_{e} S$ sea estable. Entonces

$$
\left\|W_{e}(s) S(s)\right\|_{\infty} \geq\left|W_{e}(z)\right| \prod_{i=1}^{m}\left|\frac{z+p_{i}}{z-p_{i}}\right|
$$

Ahora suponga $W_{e}(s)=\frac{s / M_{s}+\omega_{b}}{s+\omega_{b} \epsilon},\left\|W_{e} S\right\|_{\infty} \leq 1$, y que $z$ es un cero real en semiplano derecho. Entonces

$$
\frac{z / M_{s}+\omega_{b}}{z+\omega_{b} \epsilon} \leq \prod_{i=1}^{m}\left|\frac{z-p_{i}}{z+p_{i}}\right|=: \alpha
$$

lo cual da

$$
\omega_{b} \leq \frac{z}{1-\alpha \epsilon}\left(\alpha-\frac{1}{M_{s}}\right) \approx z\left(\alpha-\frac{1}{M_{s}}\right)
$$

donde $\alpha=1$ si $L$ no tiene polos en semiplano derecho. Esto muestra que el ancho de banda en lazo cerrado debe ser mucho menor que el cero en semiplano derecho. Conclusiones similares se obtienen para ceros complejos en semiplano derecho.

\section{Notas y referencias}
El diseño por conformación de lazo es bien conocido para sistemas SISO en teoría clásica de control. La idea fue extendida a sistemas MIMO por Doyle y Stein [1981] usando técnica de diseño LQG. Las limitaciones del diseño por conformación de lazo se discuten en detalle en Stein y Doyle [1991]. El Capítulo 16 presenta otro método de conformación de lazo usando teoría de control $\mathcal{H}_{\infty}$, con potencial para superar limitaciones del método LQG/LTR. Algunas discusiones adicionales sobre elección de funciones de ponderación pueden encontrarse en Skogestad y Postlethwaite [1996]. Los compromisos y limitaciones de diseño para sistemas SISO se discuten en detalle en Bode [1945], Horowitz [1963], y Doyle, Francis y Tannenbaum [1992]. La monografía de Freudenberg y Looze [1988] contiene muchas generalizaciones multivariables. La generalización multivariable de la relación integral de Bode puede encontrarse en Chen [1995], en la cual se basa la Sección 6.5. Algunos resultados relacionados pueden encontrarse en Boyd y Desoer [1985]. Resultados adicionales relacionados se encuentran en un libro más reciente de Seron, Braslavsky y Goodwin [1997].

\section{Problemas}
Problema 6.1 Sea $P$ una planta en lazo abierto. Se desea diseñar un controlador tal que el sobreimpulso sea $\leq 10 \%$ y el tiempo de establecimiento $\leq 10$ s. Estime la sensibilidad pico permitida $M_{s}$ y el ancho de banda en lazo cerrado.

Problema 6.2 Sea $L_{1}=\frac{1}{s(s+1)^{2}}$ una función de transferencia de lazo abierto de un sistema de realimentación unitaria. Encuentre el margen de fase, sobreimpulso, tiempo de establecimiento y el correspondiente $M_{s}$.

Problema 6.3 Repita el Problema 6.2 con

$$
L_{2}=\frac{100(s+10)}{(s+1)(s+2)(s+20)}
$$

Problema 6.4 Sea $P=\frac{10(1-s)}{s(s+10)}$. Use el método clásico de conformación de lazo para diseñar un controlador de adelanto o atraso de primer orden tal que el sistema tenga al menos $30^{\circ}$ de margen de fase y la mayor frecuencia de cruce posible.

Problema 6.5 Use el método del lugar de las raíces para mostrar que un sistema de fase no mínima no puede estabilizarse con un controlador de ganancia muy alta.

Problema 6.6 Sea $P=\frac{5}{(1-s)(s+2)}$. Diseñe un controlador de adelanto o atraso tal que el sistema tenga al menos $30^{\circ}$ de margen de fase con ganancia de lazo $\geq 2$ para cualquier frecuencia $\omega \leq 0.1$ y el menor ancho de banda posible (o frecuencia de cruce).

Problema 6.7 Use el método del lugar de las raíces para mostrar que un sistema inestable no puede estabilizarse con un controlador de ganancia muy baja.

Problema 6.8 Considere el lazo de realimentación unitaria con controlador propio $K(s)$ y planta estrictamente propia $P(s)$, ambos supuestos cuadrados. Suponga estabilidad interna.

\begin{enumerate}
  \item Sea $w(s)$ una función de ponderación escalar, supuesta en $\mathcal{R} \mathcal{H}_{\infty}$. Defina
\end{enumerate}

$$
\epsilon=\left\|w(I+P K)^{-1}\right\|_{\infty}, \quad \delta=\left\|K(I+P K)^{-1}\right\|_{\infty}
$$

así, $\epsilon$ mide, por ejemplo, atenuación de perturbaciones y $\delta$ mide, por ejemplo, esfuerzo de control. Derive la siguiente desigualdad, que muestra que $\epsilon$ y $\delta$ no pueden ser ambos pequeños simultáneamente en general. Para todo $\operatorname{Re} s_{0} \geq 0$

$$
\left|w\left(s_{0}\right)\right| \leq \epsilon+\left|w\left(s_{0}\right)\right| \sigma_{\min }\left[P\left(s_{0}\right)\right] \delta
$$

\begin{enumerate}
  \setcounter{enumi}{1}
  \item Si queremos muy buena atenuación de perturbaciones en una frecuencia particular, podría suponerse que necesitamos alta ganancia del controlador en esa frecuencia. Fije $\omega$ con $j \omega$ no polo de $P(s)$, y suponga
\end{enumerate}

$$
\epsilon:=\sigma_{\max }\left[(I+P K)^{-1}(j \omega)\right]<1
$$

Derive una cota inferior para $\sigma_{\min }[K(j \omega)]$. Esta cota inferior debe divergir cuando $\epsilon \rightarrow 0$.

Problema 6.9 Suponga que $P$ es propia y tiene un cero en semiplano derecho en $s=z>0$. Suponga que $y=\frac{w}{1+P K}$, donde $w$ es un escalón unitario en el tiempo $t=0$, y que nuestra especificación de desempeño es

$$
|y(t)| \leq \begin{cases}\alpha, & \text { si } \quad 0 \leq t \leq T \\ \beta, & \text { si } \quad T<t\end{cases}
$$

para algunos $\alpha>1>\beta>0$. Muestre que, para que exista un controlador LTI propio e internamente estabilizante $K$ que cumpla la especificación, debe cumplirse

$$
\ln \left(\frac{\alpha-\beta}{\alpha-1}\right) \leq z T
$$

¿Qué compromisos implica esto?

Problema 6.10 Sea $K$ un controlador estabilizante para la planta

$$
\begin{gathered}
P=\frac{s-\alpha}{(s-\beta)(s+\gamma)} \\
\alpha>0, \beta>0, \gamma \geq 0 . \text { Suponga }|S(j \omega)| \leq \delta<1, \forall \omega \in\left[-\omega_{0}, \omega_{0}\right] \text { donde } \\
S(s)=\frac{1}{1+P K}
\end{gathered}
$$

Encuentre una cota inferior para $\|S\|_{\infty}$ y calcule dicha cota para $\alpha=1, \beta=2, \gamma=10$, $\delta=0.2$, y $\omega_{0}=1$.

7. Reducción de modelos equilibrados 

8. Incertidumbre y robustez 


% FIN DEL DOCUMENTO
\end{document}